% La correspondencia y los intercambios escritos — Espagnol 1ère A — Cours signé OnBuch+
% Programme officiel MINESEC. Compiler avec : tectonic EA19-correspondencia-intercambios-escritos.tex
\newcommand{\DOCMATIERE}{Espagnol}
\newcommand{\DOCNIVEAU}{1ère A}
\newcommand{\DOCMODULE}{Módulo 5 : Les mass médias et la communication}
\newcommand{\DOCLECON}{EA19}
\newcommand{\DOCTITRE}{La correspondencia y los intercambios escritos}
% OnBuch+ — préambule « cours signés » ENRICHI (compilé avec tectonic / XeLaTeX)
% Système visuel « Pop » d'OnBuch+ : fond crème (jamais blanc), bordures encre
% épaisses, ombres « dures » décalées, titres Archivo Black en capitales.
% Version MATHÉMATIQUES : graphes (pgfplots/tikz), boîtes pédagogiques
% (définition, propriété, méthode, attention, le savais-tu, exercice résolu,
% exercice, TP, bilan), page de garde et signature OnBuch+ ancrée en bas.
%
% Macros attendues dans main.tex AVANT \input{preamble} :
%   \DOCMATIERE  \DOCNIVEAU  \DOCMODULE  \DOCLECON (numéro)  \DOCTITRE
\documentclass[11pt]{article}

\usepackage{fontspec}
\usepackage[spanish,french]{babel} % français = langue principale ; l'espagnol via \esp{..} / espagnol
\usepackage[a4paper,margin=2cm,top=3.4cm,bottom=2.8cm,headheight=1.4cm,footskip=1.3cm]{geometry}
\usepackage{amsmath,amssymb,amsfonts}
\usepackage{mathtools} % \xmapsto, \coloneqq... souvent utilisés par les modèles
\usepackage{cancel} % simplifications barrées (bilans de forces)
% \abs{z} (module, valeur absolue) et \norm{u} (norme), utilisés par les modèles
\providecommand{\abs}{}\let\abs\relax\DeclarePairedDelimiter{\abs}{\lvert}{\rvert}
\providecommand{\norm}{}\let\norm\relax\DeclarePairedDelimiter{\norm}{\lVert}{\rVert}
\usepackage[table]{xcolor} % [table] : fournit aussi \rowcolors (alternance de lignes)
\usepackage{pagecolor}
\usepackage{tikz}
\usetikzlibrary{arrows.meta,positioning,calc,shapes.geometric,shapes.misc,shapes.arrows,decorations.pathmorphing,decorations.pathreplacing,decorations.markings,patterns,shadows,mindmap,trees,fit,backgrounds,matrix,intersections,angles,quotes}
\usepackage{pgfplots}
\usepackage[version=4]{mhchem} % équations nucléaires \ce{^{235}_{92}U}
\usepackage[european,siunitx]{circuitikz} % schémas électriques
\pgfplotsset{compat=1.18}
\usepgfplotslibrary{fillbetween,groupplots} % aires entre courbes (calcul intégral)
\usepackage{enumitem}
\usepackage{fancyhdr}
% [most] charge toutes les bibliothèques tcolorbox (listings, minted, raster,
% documentation...) et gonfle inutilement la mémoire TeX sur un document long
% (~300 boîtes) : on ne charge que ce qui est réellement utilisé.
\usepackage[skins,breakable]{tcolorbox}
\usepackage{graphicx}
\usepackage{colortbl}
\usepackage{booktabs}
\usepackage{tabularx}
\usepackage{diagbox} % case d'en-tête diagonale des tableaux à double entrée
% Colonnes extensibles centrée / gauche / droite (C, L, R), souvent utilisées par les modèles
\newcolumntype{C}{>{\centering\arraybackslash}X}
\newcolumntype{L}{>{\raggedright\arraybackslash}X}
\newcolumntype{R}{>{\raggedleft\arraybackslash}X}
\usepackage{array}
\usepackage{multicol}
\usepackage{multirow}
\usepackage{siunitx}
% parse-numbers=false : accepte aussi \SI{2,5 \times 10^{-3}}{...}
\sisetup{locale=FR,output-decimal-marker={,},group-minimum-digits=4,parse-numbers=false}
\DeclareSIUnit\ohm{\ensuremath{\symup{\Omega}}} % U+2126 absent de la police
\DeclareSIUnit\division{div} % réglages d'oscilloscope (ms/div, V/div)
\DeclareSIUnit\tr{tr} % tours (vitesse de rotation, ex. \SI{1500}{\tr\per\minute})
\DeclareSIUnit\molar{M} % concentration molaire (ex. \SI{3}{\molar}, \SI{1}{\micro\molar})
\DeclareSIUnit\an{an} % année (le modèle confond parfois avec \year, un compteur TeX) — ex. \SI{5}{\kilo\meter\per\an}
\DeclareSIUnit\FCFA{FCFA} % franc CFA (ex. \SI{500}{\FCFA\per\kilo\gram})
\DeclareSIUnit\degreeBrix{\textdegree Brix} % teneur en sucre d'un sirop/jus (réfractométrie)
\DeclareSIUnit\month{mois} % le modèle utilise parfois \month, un compteur TeX, comme unité de durée
\usepackage{qrcode}
% \degree : commande fréquemment utilisée par les modèles pour un angle ou une
% température en dehors de \SI{}{} (non fournie par les paquets chargés).
\providecommand{\degree}{^{\circ}}
% \qed (fin de démonstration), utilisé par les modèles sans amsthm
\providecommand{\qed}{\hfill$\square$}
% \barre : double barre des valeurs interdites dans un tableau de variations (tkz-tab)
\providecommand{\barre}{\ensuremath{\|}}
% environnement proof (amsthm non chargé) utilisé par les modèles
\ifdefined\proof\else\newenvironment{proof}[1][Démonstration]{\par\noindent\textit{#1.}\ }{\qed\par}\fi
\usepackage{lastpage}
\usepackage{hyperref}
\hypersetup{hidelinks}

% ============================================================
% Palette « Pop » OnBuch+ — mêmes hex que la classe P (plus_ui.dart)
% ============================================================
\definecolor{poporange}{HTML}{F59321}
\definecolor{popdark}{HTML}{E07A0C}
\definecolor{popcream}{HTML}{FFF6E8}
\definecolor{popink}{HTML}{1C1714}
\definecolor{poppurple}{HTML}{7A5AE0}
\definecolor{popgreen}{HTML}{2E9E6A}
\definecolor{poppink}{HTML}{E0457B}
\definecolor{popblue}{HTML}{2F7DE1}
\definecolor{popgold}{HTML}{C98A12}
\definecolor{popmuted}{HTML}{8A7F76}
\definecolor{popline}{HTML}{EADFD2}
\definecolor{popgray}{HTML}{8A7F76}
\colorlet{popgrey}{popgray}
% Alias tolérants (le modèle nomme parfois les couleurs différemment)
\colorlet{popwhite}{white}
\colorlet{popblack}{popink}
\colorlet{popred}{poppink}
\colorlet{popyellow}{popgold}
% Teintes claires (fonds de tableaux/encadrés) — le modèle nomme parfois
% les couleurs avec un suffixe L (ex. popblueL) sans les avoir définies.
\colorlet{poporangeL}{poporange!20}
\colorlet{popdarkL}{popdark!20}
\colorlet{poppurpleL}{poppurple!20}
\colorlet{popgreenL}{popgreen!20}
\colorlet{poppinkL}{poppink!20}
\colorlet{popblueL}{popblue!20}
\colorlet{popgoldL}{popgold!20}
\colorlet{popredL}{popred!20}
\colorlet{popgrayL}{popgray!20}
% Teintes claires pour fonds de tableaux / zones
\colorlet{popblueL}{popblue!10!white}
\colorlet{poporangeL}{poporange!14!white}
\colorlet{popgreenL}{popgreen!12!white}
\colorlet{poppurpleL}{poppurple!12!white}
\colorlet{poppinkL}{poppink!12!white}
\colorlet{popgoldL}{popgold!14!white}

\pagecolor{popcream}

% ============================================================
% Typographie
% ============================================================
\defaultfontfeatures{Ligatures=TeX}
\setmainfont{PlusJakartaSans-Medium.ttf}[Path=../../../fonts/,BoldFont=PlusJakartaSans-Bold.ttf]
% Symboles, API et emojis absents de la police principale : repli sur DejaVu Sans (caractères actifs)
\newfontface\symfont{DejaVuSans.ttf}[Path=../../../fonts/]
\makeatletter
\newcommand\symactive[1]{\catcode#1=13 \begingroup\lccode`\~=#1 \lowercase{\endgroup\def~}{{\symfont\char#1}}}
\newcommand\symblank[1]{\catcode#1=13 \begingroup\lccode`\~=#1 \lowercase{\endgroup\def~}{}}
\newcommand\symhyph[1]{\catcode#1=13 \begingroup\lccode`\~=#1 \lowercase{\endgroup\def~}{\mbox{-}}}
\newcount\symc
\newcommand\symrange[2]{\symc=#1 \@whilenum\symc<#2 \do{\expandafter\symactive\expandafter{\the\symc}\advance\symc by 1 }}
\newcommand\symblankrange[2]{\symc=#1 \@whilenum\symc<#2 \do{\expandafter\symblank\expandafter{\the\symc}\advance\symc by 1 }}
\makeatother
\symrange{"0250}{"02FF}\symactive{"2713}\symactive{"2714}\symactive{"2717}\symactive{"2718}\symactive{"2610}\symactive{"2611}\symactive{"2612}
\symhyph{"2011}\symhyph{"2E1A}\symblankrange{"1F300}{"1FAFF}\symblank{"FE0F}
% Maths en sans-serif (Fira Math) avec chiffres et lettres droites de Plus
% Jakarta, pour que les formules se fondent dans le texte.
\usepackage[math-style=ISO,bold-style=ISO]{unicode-math}
\setmathfont{FiraMath-Regular.otf}[Path=../../../fonts/]
\setmathfont{PlusJakartaSans-Medium.ttf}[Path=../../../fonts/,range={up/num,up/{Latin,latin},\mathup}]
\usepackage{icomma} % virgule décimale sans espace en mode math
% Caractères mathématiques Unicode absents des polices de texte (Plus Jakarta,
% Archivo Black) : rendus par leur équivalent mathématique, en texte comme en
% formule — sinon ils s'affichent en carré vide (ex. « Arithmétique dans ℤ »).
\usepackage{newunicodechar}
\newunicodechar{ℝ}{\ensuremath{\mathbb{R}}}
\newunicodechar{ℤ}{\ensuremath{\mathbb{Z}}}
\newunicodechar{ℂ}{\ensuremath{\mathbb{C}}}
\newunicodechar{ℕ}{\ensuremath{\mathbb{N}}}
\newunicodechar{ℚ}{\ensuremath{\mathbb{Q}}}
\newunicodechar{𝔻}{\ensuremath{\mathbb{D}}}
\newunicodechar{α}{\ensuremath{\alpha}}
\newunicodechar{β}{\ensuremath{\beta}}
\newunicodechar{γ}{\ensuremath{\gamma}}
\newunicodechar{δ}{\ensuremath{\delta}}
\newunicodechar{ε}{\ensuremath{\varepsilon}}
\newunicodechar{θ}{\ensuremath{\theta}}
\newunicodechar{λ}{\ensuremath{\lambda}}
\newunicodechar{μ}{\ensuremath{\mu}}
\newunicodechar{σ}{\ensuremath{\sigma}}
\newunicodechar{τ}{\ensuremath{\tau}}
\newunicodechar{φ}{\ensuremath{\varphi}}
\newunicodechar{ψ}{\ensuremath{\psi}}
\newunicodechar{ω}{\ensuremath{\omega}}
\newunicodechar{Σ}{\ensuremath{\Sigma}}
% majuscules produites par \MakeUppercase dans les titres (α-aminés -> Α)
\newunicodechar{Α}{\ensuremath{\alpha}}
\newunicodechar{Β}{\ensuremath{\beta}}
\newunicodechar{Γ}{\ensuremath{\Gamma}}
\newunicodechar{Δ}{\ensuremath{\Delta}}
\newunicodechar{Ε}{\ensuremath{\varepsilon}}
\newunicodechar{Θ}{\ensuremath{\Theta}}
\newunicodechar{Λ}{\ensuremath{\Lambda}}
\newunicodechar{Μ}{\ensuremath{\mu}}
\newunicodechar{Τ}{\ensuremath{\tau}}
\newunicodechar{Φ}{\ensuremath{\Phi}}
\newunicodechar{Ψ}{\ensuremath{\Psi}}
\newunicodechar{Ω}{\ensuremath{\Omega}}
\newunicodechar{↦}{\ensuremath{\mapsto}}
\newunicodechar{∈}{\ensuremath{\in}}
\newunicodechar{∉}{\ensuremath{\notin}}
\newunicodechar{∧}{\ensuremath{\wedge}}
\newunicodechar{⇔}{\ensuremath{\Leftrightarrow}}
\newunicodechar{⟺}{\ensuremath{\Longleftrightarrow}}
\newunicodechar{⇒}{\ensuremath{\Rightarrow}}
\newunicodechar{⟹}{\ensuremath{\Longrightarrow}}
\newunicodechar{∘}{\ensuremath{\circ}}
\newunicodechar{∪}{\ensuremath{\cup}}
\newunicodechar{∩}{\ensuremath{\cap}}
\newunicodechar{≡}{\ensuremath{\equiv}}
\newunicodechar{⊂}{\ensuremath{\subset}}
\newunicodechar{⊥}{\ensuremath{\perp}}
\newunicodechar{∃}{\ensuremath{\exists}}
\newunicodechar{∀}{\ensuremath{\forall}}
\newunicodechar{∛}{\ensuremath{\sqrt[3]{\,}}}
\newunicodechar{∜}{\ensuremath{\sqrt[4]{\,}}}
\newunicodechar{⃗}{\ensuremath{{}^{\to}}}
\newunicodechar{ⁿ}{\ifmmode^{n}\else\textsuperscript{\ensuremath{n}}\fi}
\newunicodechar{ˣ}{\ifmmode^{x}\else\textsuperscript{\ensuremath{x}}\fi}
\newunicodechar{ᵇ}{\ifmmode^{b}\else\textsuperscript{\ensuremath{b}}\fi}
\newunicodechar{ʳ}{\ifmmode^{r}\else\textsuperscript{\ensuremath{r}}\fi}
\newunicodechar{ᵘ}{\ifmmode^{u}\else\textsuperscript{\ensuremath{u}}\fi}
\newunicodechar{ʸ}{\ifmmode^{y}\else\textsuperscript{\ensuremath{y}}\fi}
\newunicodechar{ᵃ}{\ifmmode^{a}\else\textsuperscript{\ensuremath{a}}\fi}
\newunicodechar{ᵉ}{\ifmmode^{e}\else\textsuperscript{\ensuremath{e}}\fi}
\newunicodechar{ᵏ}{\ifmmode^{k}\else\textsuperscript{\ensuremath{k}}\fi}
\newunicodechar{ᵖ}{\ifmmode^{p}\else\textsuperscript{\ensuremath{p}}\fi}
\newunicodechar{ᴺ}{\ifmmode^{N}\else\textsuperscript{\ensuremath{N}}\fi}
\newunicodechar{ᶜ}{\ifmmode^{c}\else\textsuperscript{\ensuremath{c}}\fi}
\newunicodechar{ʰ}{\ifmmode^{h}\else\textsuperscript{\ensuremath{h}}\fi}
\newunicodechar{ᵠ}{\ifmmode^{q}\else\textsuperscript{\ensuremath{q}}\fi}
\newunicodechar{ₙ}{\ifmmode_{n}\else\textsubscript{\ensuremath{n}}\fi}
\newunicodechar{ₐ}{\ifmmode_{a}\else\textsubscript{\ensuremath{a}}\fi}
\newunicodechar{ᵢ}{\ifmmode_{i}\else\textsubscript{\ensuremath{i}}\fi}
\newunicodechar{ᵣ}{\ifmmode_{r}\else\textsubscript{\ensuremath{r}}\fi}
\newunicodechar{ₖ}{\ifmmode_{k}\else\textsubscript{\ensuremath{k}}\fi}
\newunicodechar{ᵦ}{\ifmmode_{\beta}\else\textsubscript{\ensuremath{\beta}}\fi}
\newunicodechar{ₓ}{\ifmmode_{x}\else\textsubscript{\ensuremath{x}}\fi}
\newfontfamily\popdisplay{ArchivoBlack-Regular.ttf}[Path=../../../fonts/]
\newfontfamily\poplogo{SpaceGrotesk-Bold.ttf}[Path=../../../fonts/]
\newfontfamily\popsemi{PlusJakartaSans-SemiBold.ttf}[Path=../../../fonts/]
\newfontfamily\popxbold{PlusJakartaSans-ExtraBold.ttf}[Path=../../../fonts/]
\newfontfamily\popxboldls{PlusJakartaSans-ExtraBold.ttf}[Path=../../../fonts/,LetterSpace=3.0]

\setlist[enumerate]{leftmargin=1.7em,itemsep=5pt,topsep=4pt}
\setlist[itemize]{leftmargin=1.7em,itemsep=4pt,topsep=4pt,label={\color{poporange}\textbullet}}
\setlength{\parindent}{0pt}
\setlength{\parskip}{5pt}
\renewcommand{\arraystretch}{1.25}

% pgfplots : style Pop par défaut
\pgfplotsset{
  every axis/.append style={axis line style={popink,line width=1pt},tick style={popink},
    grid style={popline,line width=0.5pt},label style={font=\small},tick label style={font=\footnotesize}},
  popcurve/.style={poporange,line width=1.8pt,no markers,smooth},
}
% Même style disponible dans un \draw TikZ simple (hors pgfplots)
\tikzset{popcurve/.style={poporange,line width=1.8pt}}
\tikzset{popgrid/.style={popline,very thin}}
\colorlet{popcurve}{poporange} % le nom du style est parfois utilisé comme couleur

% ============================================================
% Pill / squiggle
% ============================================================
\newcommand{\poppill}[3]{%
  \tikz[baseline=-0.55ex]\node[draw=popink,line width=1.2pt,rounded corners=2.6pt,%
    fill=#2,inner xsep=6.5pt,inner ysep=3pt]{%
    \popxboldls\fontsize{7}{7}\selectfont\color{#3}\MakeUppercase{#1}};%
}
\newcommand{\popsquiggle}[1]{%
  \tikz[baseline=-0.4ex]{
    \draw[#1,line width=2.6pt,line cap=round]
      plot [smooth] coordinates {(0,0.09) (0.34,0.01) (0.68,0.21) (1.02,0.01) (1.36,0.10)};
  }%
}
% Mot-clé surligné (marqueur orange sous le texte)
\newcommand{\cle}[1]{{\popxbold\color{popdark}#1}}

% ============================================================
% En-tête / pied
% ============================================================
\pagestyle{fancy}
\fancyhf{}
\renewcommand{\headrulewidth}{0pt}
\renewcommand{\footrulewidth}{0pt}
\lhead{\raisebox{-0.15ex}{\poplogo\fontsize{13}{13}\selectfont\color{popink}OnBuch\color{poporange}+}}
\rhead{\poppill{\DOCMATIERE\ \textbullet\ \DOCNIVEAU\ \textbullet\ Leçon \DOCLECON}{white}{popink}}
\lfoot{\popxbold\fontsize{7.6}{7.6}\selectfont\color{popmuted}\MakeUppercase{Cours signé OnBuch+}\ \textbullet\ {\popsemi\DOCTITRE}}
\rfoot{\tikz[baseline=-0.6ex]\node[rounded corners=8.5pt,fill=popink,minimum height=17pt,minimum width=17pt,inner xsep=5pt,inner sep=0]{\popxbold\fontsize{8}{8}\selectfont\color{popcream}\thepage\,/\,\pageref*{LastPage}};}

% ============================================================
% Titres
% ============================================================
\newcommand{\chaptitre}[2]{%
  \begin{tcolorbox}[enhanced,colback=poporange,colframe=popink,boxrule=2.6pt,arc=11pt,
      drop shadow=popink,left=15pt,right=15pt,top=12pt,bottom=11pt,before skip=3pt,after skip=18pt]
    \poppill{#2}{popink}{popcream}\par\vspace{9pt}
    {\raggedright\hyphenpenalty=10000\exhyphenpenalty=10000\popdisplay\fontsize{22}{24}\selectfont\color{popcream}\MakeUppercase{#1}\par}
  \end{tcolorbox}
}
% Section numérotée : pastille orange avec le numéro + titre Archivo
\newcounter{popsec}
\newcommand{\coursec}[1]{%
  \stepcounter{popsec}\setcounter{popsub}{0}%
  \par\vspace{16pt}\noindent
  \begin{minipage}{\linewidth}
  \tikz[baseline=-0.7ex]\node[draw=popink,line width=1.6pt,fill=poporange,rounded corners=4pt,minimum size=22pt,inner sep=0]{\popdisplay\fontsize{12}{12}\selectfont\color{popcream}\thepopsec};%
  \hspace{8pt}{\popdisplay\fontsize{15}{17}\selectfont\color{popink}\MakeUppercase{#1}}\par
  \vspace{4pt}\noindent{\color{poporange}\rule{36pt}{3.6pt}}
  \end{minipage}\par\nopagebreak\vspace{8pt}
}
\newcounter{popsub}
\newcommand{\courssub}[1]{%
  \stepcounter{popsub}%
  \par\vspace{9pt}\noindent{\popxbold\fontsize{11.5}{13}\selectfont\color{popink}{\color{poporange}\thepopsec.\thepopsub}\hspace{6pt}#1}\par\nopagebreak\vspace{4pt}
}

% ============================================================
% Boîtes pédagogiques — même recette : fond blanc, bordure épaisse colorée,
% ombre dure encre, badge plein. Titre optionnel [#1].
% ============================================================
\tcbset{popbase/.style={breakable,enhanced,colback=white,boxrule=2.3pt,arc=9pt,
  shadow={3pt}{-3pt}{0pt}{popink},left=14pt,right=14pt,top=11pt,bottom=11pt,before skip=11pt,after skip=15pt,
  fontupper=\popsemi\fontsize{10.6}{14.5}\selectfont\color{popink},
  fontlower=\popsemi\fontsize{10.6}{14.5}\selectfont\color{popink}}}
% Coche dessinée (le glyphe ✓ n'existe pas dans les polices du document)
\newcommand{\popcheck}[1]{\tikz[baseline=-0.5ex]\draw[#1,line width=1.6pt,line cap=round,line join=round](-0.13,0.0)--(-0.04,-0.09)--(0.13,0.1);}
\providecommand{\coche}{\popcheck{popgreen}}
\providecommand{\resultat}[1]{\par\smallskip\noindent\fcolorbox{popgreen}{popgreenL}{\parbox{\dimexpr\linewidth-2\fboxsep-2\fboxrule\relax}{\textbf{Résultat.} #1}}\par\smallskip}
\newcommand{\popboxhead}[3]{\poppill{#1}{#2}{white}\ifx\relax#3\relax\else\hspace{6pt}{\popxbold\fontsize{10.6}{12}\selectfont\color{popink}#3}\fi\par\vspace{7pt}}

\newtcolorbox{aretenir}[1][]{popbase,colframe=popgreen,before upper={\popboxhead{À retenir}{popgreen}{#1}}}
% Texte en espagnol (document de lecture, dialogue, poème, extrait) : cadre
% d'étude. Un titre optionnel (source, auteur) s'affiche dans l'en-tête.
\newtcolorbox{textebox}[1][]{popbase,colframe=popblue,before upper={\popboxhead{Texto}{popblue}{#1}\begin{otherlanguage}{spanish}},after upper={\end{otherlanguage}}}
% Marques ✓ / ✗ (forme correcte / fautive) souvent utilisées par les modèles
\providecommand{\cross}{\textcolor{poppink}{\ensuremath{\times}}}
\providecommand{\xmark}{\cross}\providecommand{\crossmark}{\cross}\providecommand{\cmark}{\ensuremath{\checkmark}}
% Ligne de réponse / justification à compléter (inventée par les modèles)
\providecommand{\justification}[1]{\par\noindent\textit{Justification~:} \dotfill\par}
% \textipa (package tipa non chargé) : la police ne couvre pas l'API -> simple passage
\providecommand{\textipa}[1]{#1}
% Soulignement (ulem non chargé) : \uline, \ul
\providecommand{\uline}[1]{\underline{#1}}\providecommand{\ul}[1]{\underline{#1}}
% \glossaire{\item ...} inventée par les modèles : liste de glossaire
\providecommand{\glossaire}[1]{\begin{itemize}#1\end{itemize}}
% \alert (beamer) inventée par les modèles : texte mis en évidence
\providecommand{\alert}[1]{\textbf{\textcolor{poppink}{#1}}}
% variantes ASCII d'ordinaux écrites en macro
\providecommand{\iere}{\textsuperscript{re}}\providecommand{\ieme}{\textsuperscript{e}}\providecommand{\ere}{\textsuperscript{re}}\providecommand{\eme}{\textsuperscript{e}}\providecommand{\er}{\textsuperscript{er}}\providecommand{\ie}{\textsuperscript{e}}
% Ordinaux français écrits en macro par les modèles : 1\ière, 3\ième
\newcommand{\ière}{\textsuperscript{re}}\newcommand{\ième}{\textsuperscript{e}}
% \exemple (inventée par les modèles) : étiquette « Exemple : »
\providecommand{\exemple}{\textbf{Exemple~:}\ }
% Mot ou phrase en espagnol à l'intérieur d'un paragraphe français
\newcommand{\esp}[1]{\ifmmode\text{\textit{#1}}\else\begingroup\selectlanguage{spanish}\itshape #1\endgroup\fi}
\newtcolorbox{exemplebox}[1][]{popbase,colframe=poporange,before upper={\popboxhead{Exemple}{poporange}{#1}}}
\newtcolorbox{definition}[1][]{popbase,colframe=popblue,before upper={\popboxhead{Définition}{popblue}{#1}}}
\newtcolorbox{propriete}[1][]{popbase,colframe=poppurple,before upper={\popboxhead{Propriété}{poppurple}{#1}}}
\newtcolorbox{methode}[1][]{popbase,colframe=popgold,before upper={\popboxhead{Méthode}{popgold}{#1}}}
\newtcolorbox{attention}[1][]{popbase,colframe=poppink,before upper={\popboxhead{Attention}{poppink}{#1}}}
\newtcolorbox{experience}[1][]{popbase,colframe=popblue,colback=popblueL,before upper={\popboxhead{Expérience}{popblue}{#1}}}
% « Le savais-tu ? » : carte violette pleine (contexte, culture, Cameroun)
\newtcolorbox{savaistu}[1][]{popbase,colback=poppurple,colframe=popink,
  fontupper=\popsemi\fontsize{10.4}{14.2}\selectfont\color{white},
  before upper={\poppill{Le savais-tu\,?}{white}{poppurple}\ifx\relax#1\relax\else\hspace{6pt}{\popxbold\color{white}#1}\fi\par\vspace{7pt}}}
% Exercice résolu : énoncé en haut, \tcblower puis la solution sur fond crème
\newcounter{popexo}\newcounter{popexr}
\newtcolorbox{exoresolu}[1][]{popbase,colframe=poporange,colbacklower=popcream,
  segmentation style={popink,line width=1.2pt,dashed},
  before upper={\stepcounter{popexr}\popboxhead{Exercice résolu \thepopexr}{poporange}{#1}},
  before lower={\poppill{Solution}{popink}{popcream}\par\vspace{6pt}}}
% Exercice (non résolu en place) — la correction est dans la partie Corrigés
\newtcolorbox{exercice}[1][]{popbase,colframe=popink,
  before upper={\stepcounter{popexo}\popboxhead{Exercice \thepopexo}{popink}{#1}}}
% Corrigé
\newtcolorbox{corrige}[1][]{popbase,colframe=popgreen,colback=popgreenL,
  before upper={\popboxhead{Corrigé}{popgreen}{#1}}}
% Objectifs / prérequis / bilan
\newtcolorbox{objectifs}{popbase,colframe=popink,colback=poporangeL,
  before upper={\popboxhead{Objectifs de la leçon}{popink}{}}}
\newtcolorbox{prerequis}{popbase,colframe=popink,colback=popblueL,
  before upper={\popboxhead{Prérequis}{popblue}{}}}
\newtcolorbox{bilan}{popbase,colframe=popink,colback=white,boxrule=2.8pt,
  before upper={\popboxhead{Fiche bilan}{poporange}{L'essentiel à connaître pour l'examen}}}
% Figure encadrée avec légende
\newtcolorbox{popfigure}[1][]{enhanced,breakable=false,colback=white,colframe=popline,boxrule=1.4pt,arc=8pt,
  left=10pt,right=10pt,top=10pt,bottom=8pt,before skip=10pt,after skip=12pt,halign=center,
  lower separated=false,halign lower=center,
  fontlower=\popsemi\fontsize{9}{11.5}\selectfont\color{popmuted},
  #1}
\newcommand{\legende}[1]{\par\vspace{6pt}{\popsemi\fontsize{9}{11.5}\selectfont\color{popmuted}#1}}

% ============================================================
% Signature OnBuch+
% ============================================================
\newcommand{\poplogoinline}[1]{{\poplogo\fontsize{#1}{#1}\selectfont\color{popink}OnBuch\color{poporange}+}}
\newcommand{\signatureonbuch}{%
  \begin{tcolorbox}[enhanced,colback=popink,colframe=popink,boxrule=2.2pt,arc=10pt,
      drop shadow=poporange,left=14pt,right=14pt,top=10pt,bottom=10pt,before skip=0pt,after skip=0pt]
    \tikz[baseline=-0.7ex]\node[circle,fill=popgreen,draw=popcream,line width=1.3pt,minimum size=22pt,inner sep=0]{\popcheck{white}};%
    \hspace{9pt}\begin{minipage}[c]{0.62\linewidth}
      {\popxbold\fontsize{12}{13}\selectfont\color{popcream}Signé\ \textbullet\ {\poplogo OnBuch\color{poporange}+}}\par\vspace{2pt}
      {\raggedright\popsemi\fontsize{8}{10.5}\selectfont\color{popline}\MakeUppercase{Équipe pédagogique OnBuch+}\\\MakeUppercase{Conforme au programme officiel MINESEC}\par}
    \end{minipage}\hfill
    {\poplogo\fontsize{22}{22}\selectfont\color{popcream}OnBuch\color{poporange}+}
  \end{tcolorbox}%
}

% ============================================================
% Page de garde
% #1 titre · #2 matière · #3 niveau · #4 module · #5 n° leçon · #6 durée ·
% #7 description / objectifs (texte ou itemize)
% ============================================================
\newcommand{\pagegarde}[7]{%
  \thispagestyle{empty}
  \newgeometry{a4paper,margin=1.7cm,top=1.6cm,bottom=1.4cm}%
  \begin{tikzpicture}[remember picture,overlay]
    \fill[poporange!18] ([xshift=-3.2cm,yshift=-3.6cm]current page.north east) circle (4.6cm);
    \fill[poppurple!14] ([xshift=2.4cm,yshift=7.6cm]current page.south west) circle (3.8cm);
  \end{tikzpicture}%
  {\poplogo\fontsize{30}{30}\selectfont\color{popink}OnBuch\color{poporange}+}\hfill
  \raisebox{0.6ex}{\poppill{Cours signé \textbullet\ Premium}{popink}{popcream}}\par
  \vspace{1pt}\popsquiggle{poppurple}\par
  \vspace{8pt}
  \begin{tcolorbox}[enhanced,colback=poporange,colframe=popink,boxrule=2.8pt,arc=13pt,
      drop shadow=popink,left=18pt,right=18pt,top=12pt,bottom=13pt,after skip=10pt]
    \poppill{#2 \textbullet\ #3}{popink}{popcream}\hspace{5pt}\poppill{#4}{white}{popink}\par\vspace{8pt}
    {\popdisplay\fontsize{14}{14}\selectfont\color{popink}LEÇON #5}\par\vspace{4pt}
    {\raggedright\hyphenpenalty=10000\exhyphenpenalty=10000\popdisplay\fontsize{25}{27}\selectfont\color{popcream}\MakeUppercase{#1}\par}
    \vspace{8pt}
    {\popxbold\fontsize{9.5}{9.5}\selectfont\color{popink}\MakeUppercase{Durée conseillée : #6}}
  \end{tcolorbox}
  \begin{tcolorbox}[enhanced,colback=white,colframe=popink,boxrule=2.2pt,arc=10pt,
      drop shadow=popink,left=16pt,right=16pt,top=10pt,bottom=10pt,after skip=8pt]
    \poppill{Dans cette leçon}{poppurple}{white}\par\vspace{5pt}
    \popsemi\fontsize{9.3}{12.6}\selectfont\color{popink}#7
  \end{tcolorbox}
  \noindent\begin{minipage}[t]{0.487\linewidth}
    \begin{tcolorbox}[enhanced,colback=popgreen,colframe=popink,boxrule=2.2pt,arc=10pt,
        drop shadow=popink,left=11pt,right=11pt,top=10pt,bottom=10pt,height=3.1cm,valign=center]
      \tcbox[colback=white,colframe=popink,boxrule=1.3pt,arc=5pt,boxsep=0pt,nobeforeafter,
        left=3pt,right=3pt,top=3pt,bottom=3pt,valign=center]{\qrcode[height=1.9cm]{https://play.google.com/store/apps/details?id=com.luvvix.onbuch&hl=fr}}%
      \hspace{8pt}\begin{minipage}[c]{0.5\linewidth}
        {\popdisplay\fontsize{11}{12.5}\selectfont\color{white}\MakeUppercase{Télécharge OnBuch}}\par\vspace{4pt}
        {\raggedright\popsemi\fontsize{8.4}{11}\selectfont\color{white}Gratuit \textbullet\ Google Play\\Tuteur IA, annales, concours, résultats\par}
      \end{minipage}
    \end{tcolorbox}
  \end{minipage}\hfill
  \begin{minipage}[t]{0.487\linewidth}
    \begin{tcolorbox}[enhanced,colback=poppurple,colframe=popink,boxrule=2.2pt,arc=10pt,
        drop shadow=popink,left=11pt,right=11pt,top=10pt,bottom=10pt,height=3.1cm,valign=center]
      \tikz[baseline=-0.7ex]\node[circle,fill=white,draw=popink,line width=1.3pt,minimum size=1.6cm,inner sep=0]{\popdisplay\fontsize{24}{24}\selectfont\color{poppurple}+};%
      \hspace{8pt}\begin{minipage}[c]{0.6\linewidth}
        {\popdisplay\fontsize{11}{12.5}\selectfont\color{white}\MakeUppercase{Passe à OnBuch+}}\par\vspace{4pt}
        {\raggedright\popsemi\fontsize{8.4}{11}\selectfont\color{white}Tous les cours signés, Tuteur IA illimité, fiches PDF — onglet OnBuch+ de l'app\par}
      \end{minipage}
    \end{tcolorbox}
  \end{minipage}
  \vspace{8pt}
  \popcontenu
  \vfill
  \signatureonbuch
  \restoregeometry
  \newpage
}

% Rangée « Ce cours contient » (page de garde)
\newcommand{\poptile}[3]{% #1 couleur #2 gros texte #3 légende
  \begin{tcolorbox}[enhanced,colback=#1,colframe=popink,boxrule=2pt,arc=9pt,shadow={2.5pt}{-2.5pt}{0pt}{popink},
      left=6pt,right=6pt,top=9pt,bottom=9pt,halign=center,valign=center,height=1.75cm,width=\linewidth,nobeforeafter]
    {\popdisplay\fontsize{12.5}{13}\selectfont\color{white}\MakeUppercase{#2}}\par\vspace{3pt}
    {\centering\popsemi\fontsize{7.8}{9.5}\selectfont\color{white}#3\par}
  \end{tcolorbox}}
\newcommand{\popcontenu}{%
  \par\noindent{\popxboldls\fontsize{7.5}{7.5}\selectfont\color{popmuted}CE COURS SIGNÉ CONTIENT}\par\vspace{7pt}
  \noindent\begin{minipage}[t]{0.235\linewidth}\poptile{popblue}{Cours}{complet, illustré, pas à pas}\end{minipage}\hfill
  \begin{minipage}[t]{0.235\linewidth}\poptile{poporange}{Méthodes}{et exercices résolus}\end{minipage}\hfill
  \begin{minipage}[t]{0.235\linewidth}\poptile{poppink}{Exercices}{type examen + corrigés}\end{minipage}\hfill
  \begin{minipage}[t]{0.235\linewidth}\poptile{popgreen}{Fiche bilan}{l'essentiel à retenir}\end{minipage}\par}

% Sommaire
\newtcolorbox{sommaire}{enhanced,colback=white,colframe=popink,boxrule=2.2pt,arc=10pt,
  drop shadow=popink,left=18pt,right=18pt,top=13pt,bottom=13pt,before skip=6pt,after skip=20pt,
  before upper={\poppill{Sommaire}{poppurple}{white}\par\vspace{9pt}\popsemi\fontsize{10.6}{14.5}\selectfont\color{popink}}}

% Signature de fin de cours — poussée tout en bas de la dernière page
\newcommand{\findecours}{%
  \par\vfill\nopagebreak
  \begin{center}\popsquiggle{poporange}\end{center}\vspace{4pt}
  \signatureonbuch
}
\AtBeginDocument{\def\llbracket{\mathopen{[\mkern-3.5mu[}}\def\rrbracket{\mathclose{]\mkern-3.5mu]}}} % intervalles d'entiers (glyphe ⟦ absent de la police)
% Glyphes absents de Fira Math : redéfinis à partir de glyphes disponibles
\AtBeginDocument{%
  \def\setminus{\mathbin{\backslash}}\let\smallsetminus\setminus
  \def\vdots{\mathord{\vbox{\baselineskip3.2pt\lineskiplimit0pt\kern4pt\hbox{.}\hbox{.}\hbox{.}}}}%
  \def\ddots{\mathinner{\mkern1mu\raise6.5pt\hbox{.}\mkern2mu\raise3.6pt\hbox{.}\mkern2mu\raise0.7pt\hbox{.}\mkern1mu}}%
  \def\triangle{\mathord{\scalebox{0.9}{$\symup{\Delta}$}}}%
  \def\nabla{\mathord{\raisebox{\depth}{\scalebox{1}[-1]{$\symup{\Delta}$}}}}%
  \def\checkmark{\popcheck{popdark}}%
  \def\star{\mathbin{\ast}}\def\bigstar{\mathord{\ast}}%
  \def\diamond{\mathbin{\scalebox{0.6}{$\Diamond$}}}%
  \def\bigcup{\mathop{\mathchoice{\raisebox{-0.3em}{\scalebox{1.7}{$\cup$}}}{\scalebox{1.25}{$\cup$}}{\cup}{\cup}}\displaylimits}%
  \def\bigcap{\mathop{\mathchoice{\raisebox{-0.3em}{\scalebox{1.7}{$\cap$}}}{\scalebox{1.25}{$\cap$}}{\cap}{\cap}}\displaylimits}%
  \def\lesseqqgtr{\mathrel{\lessgtr}}\def\bigoplus{\mathop{\oplus}\displaylimits}%
}
\providecommand{\ssi}{si et seulement si} % abréviation fréquente
\AtBeginDocument{\providecommand{\wideparen}{\overparen}} % arc de cercle

%% FIGFIX-BEGIN v4 — correctifs globaux des figures (docs/onbuch-generation/tools/figfix)
\usepackage{adjustbox}\usepackage{etoolbox}
% toute figure TikZ/pgfplots plus large que la ligne est réduite à la largeur disponible
\BeforeBeginEnvironment{tikzpicture}{\begin{adjustbox}{max width=\linewidth}}
\AfterEndEnvironment{tikzpicture}{\end{adjustbox}}
% légendes pgfplots lisibles par-dessus les courbes
\pgfplotsset{every axis/.append style={legend style={fill=white,fill opacity=0.92,text opacity=1,draw=black!15,inner sep=3pt}}}
% étiquettes d'arêtes (syntaxe edge["..."]) avec fond blanc
\tikzset{every edge quotes/.append style={fill=white,inner sep=1.2pt}}
%% FIGFIX-END
\begin{document}

\pagegarde{La correspondencia y los intercambios escritos}{Espagnol}{1ère A}{Módulo 5 : Les mass médias et la communication}{EA19}{2 h}{Cette leçon mixte du module « Les mass médias et la communication » apprend à rédiger lettres amicales, lettres formelles et courriels en espagnol. À travers un texte support, du lexique thématique et des points de grammaire appliqués à la correspondance, l'élève acquiert la compétence clé : adapter le registre au destinataire.
\vspace{4pt}
\begin{multicols}{2}\small
\begin{itemize}
\item À la fin de cette leçon, je sais identifier la structure d'une lettre amicale et d'une lettre formelle en espagnol
\item À la fin de cette leçon, je sais utiliser les formules d'ouverture et de clôture adaptées (Querido/a…, Estimado señor…, Un abrazo, Le saluda atentamente)
\item À la fin de cette leçon, je sais distinguer le registre familier (tuteo) du registre formel (usted) et les employer correctement
\item À la fin de cette leçon, je sais rédiger un courriel ou un message court avec une structure claire (asunto, saludo, cuerpo, despedida)
\item À la fin de cette leçon, je sais rédiger une lettre amicale d'environ 80 mots à un correspondant hispanophone
\item À la fin de cette leçon, je sais employer le vocabulaire de la correspondance et des échanges écrits
\end{itemize}\end{multicols}}

\begin{sommaire}
\begin{multicols}{2}
\begin{enumerate}
\item Texte support : Dos mensajes, dos registros
\item Lexique : el mundo de la correspondencia
\item Grammaire appliquée : tú, usted et l'impératif de politesse
\item Structure de la lettre et du courriel
\item Production guidée : escribo a mi correspondiente
\item Activité d'intégration
\item Méthodes et exercices résolus
\item Exercices
\item Corrigés des exercices
\item Fiche bilan
\end{enumerate}
\end{multicols}
\end{sommaire}

\begin{objectifs}
\begin{itemize}
\item À la fin de cette leçon, je sais identifier la structure d'une lettre amicale et d'une lettre formelle en espagnol
\item À la fin de cette leçon, je sais utiliser les formules d'ouverture et de clôture adaptées (Querido/a…, Estimado señor…, Un abrazo, Le saluda atentamente)
\item À la fin de cette leçon, je sais distinguer le registre familier (tuteo) du registre formel (usted) et les employer correctement
\item À la fin de cette leçon, je sais rédiger un courriel ou un message court avec une structure claire (asunto, saludo, cuerpo, despedida)
\item À la fin de cette leçon, je sais rédiger une lettre amicale d'environ 80 mots à un correspondant hispanophone
\item À la fin de cette leçon, je sais employer le vocabulaire de la correspondance et des échanges écrits
\item À la fin de cette leçon, je sais traduire des phrases de correspondance du français vers l'espagnol et inversement
\item À la fin de cette leçon, je sais éviter les erreurs fréquentes des francophones dans les formules de politesse
\end{itemize}
\end{objectifs}

\begin{prerequis}
\begin{itemize}
\item Les salutations et formules de politesse de base vues en Seconde (Hola, Buenos días, Gracias, Adiós)
\item Le présent de l'indicatif des verbes réguliers et des verbes fréquents (ser, estar, querer, poder)
\item Les pronoms personnels sujets et la différence tú / usted (initiée en Seconde)
\item Le vocabulaire des médias et de la communication (leçons EA17-EA18 du module 5)
\item La structure simple de la phrase affirmative et interrogative en espagnol
\end{itemize}
\end{prerequis}

\newpage

\coursec{Texte support : Dos mensajes, dos registros}

\courssub{Situation de départ : le dilemme de Mariam}

Il est seize heures au lycée de Yaoundé. Mariam, élève de Première A, ouvre sa boîte de réception sur son téléphone et découvre deux nouveaux messages. Le premier vient de son correspondant colombien, Andrés, un garçon de dix-sept ans qui vit à Bogotá : il l'invite à passer les vacances chez lui, en Colombie. Quelle joie ! Mariam veut lui répondre tout de suite, avec chaleur, avec des mots d'amitié. Mais le deuxième message vient du consulat d'Espagne : on lui confirme que sa demande d'informations sur un échange scolaire a bien été reçue, et on lui demande de renvoyer un courriel avec des documents précis. Mariam doit donc écrire deux réponses très différentes, le même jour, dans la même langue.

Peut-elle utiliser les mêmes mots dans les deux messages ? Bien sûr que non : on ne dit pas \esp{Un abrazo} à un consul ! On ne commence pas une lettre à un ami par \esp{Estimado señor}. En espagnol comme en français, chaque destinataire exige un \cle{registre} de langue différent : un registre \cle{familier} pour les amis et la famille, un registre \cle{formel} (ou soutenu) pour les administrations, les professeurs, les personnes inconnues.

\textbf{Problématique :} comment reconnaître le registre d'un message écrit, et comment adapter notre propre écrit au destinataire ?

\courssub{Lecture du texte : Dos mensajes, dos registros}

Lisons d'abord les deux messages que Mariam a reçus. Le premier est une lettre amicale ; le second est un courriel formel.

\begin{textebox}[Texto 1 : una carta amistosa de Andrés]
Querida Mariam:

¡Hola! ¿Qué tal? ¿Cómo estás? Te escribo para invitarte a pasar las vacaciones en Bogotá conmigo y con mi familia. Mi madre está muy contenta y quiere conocerte por fin. Ella cocina muy bien: vas a probar el ajiaco y las arepas, ¡te van a encantar!

Aquí hace un poco de frío porque Bogotá está muy alta, así que trae un abrigo. Podemos visitar el centro histórico, subir al cerro de Monserrate y bailar salsa con mis primos. ¿Te gusta bailar? Yo sé que sí, jaja.

Dime pronto si puedes venir. Pregunta a tus padres y mándame la fecha de tu llegada. Te espero con muchas ganas.

Un abrazo fuerte,

Andrés
\end{textebox}

\begin{textebox}[Texto 2 : un correo formal del consulado]
Estimada señora:

Le escribo para confirmar la recepción de su solicitud de información sobre el intercambio escolar entre su instituto y el colegio Cervantes de Madrid. Según nuestra secretaría, usted debe enviar los siguientes documentos: una fotocopia del pasaporte, un certificado de estudios y una carta de motivación.

Le ruego que adjunte estos documentos a su próximo correo electrónico antes del día quince de este mes. Si necesita más información, puede consultar nuestra página web o llamar a nuestra oficina de lunes a viernes, de nueve a catorce horas.

Sin otro particular, le saluda atentamente,

Consulado de España en Yaundé
\end{textebox}

\textbf{Traduction du texte 1 :} « Chère Mariam : Salut ! Comment ça va ? Comment vas-tu ? Je t'écris pour t'inviter à passer les vacances à Bogotá avec moi et ma famille. Ma mère est très contente et veut enfin te connaître. Elle cuisine très bien : tu vas goûter l'ajiaco et les arepas, tu vas adorer ! Ici il fait un peu froid parce que Bogotá est située très haut, alors apporte un manteau. Nous pouvons visiter le centre historique, monter au mont Monserrate et danser la salsa avec mes cousins. Tu aimes danser ? Moi je sais que oui, haha. Dis-moi vite si tu peux venir. Demande à tes parents et envoie-moi la date de ton arrivée. Je t'attends avec impatience. Je t'embrasse fort, Andrés. »

\textbf{Traduction du texte 2 :} « Madame : Je vous écris pour confirmer la réception de votre demande d'informations sur l'échange scolaire entre votre lycée et le collège Cervantes de Madrid. Selon notre secrétariat, vous devez envoyer les documents suivants : une photocopie du passeport, un certificat de scolarité et une lettre de motivation. Je vous prie de joindre ces documents à votre prochain courriel avant le quinze de ce mois. Si vous avez besoin de plus d'informations, vous pouvez consulter notre page web ou appeler notre bureau du lundi au vendredi, de neuf à quatorze heures. Veuillez agréer, Madame, mes salutations distinguées. Consulat d'Espagne à Yaoundé. »

\begin{definition}[Glossaire : les mots difficiles]
\begin{itemize}
\item \esp{querido/a} : cher, chère (ouverture de lettre amicale).
\item \esp{un abrazo (fuerte)} : une (forte) embrassade, formule de clôture amicale.
\item \esp{estimado/a} : cher, chère (ouverture formelle ; se traduit en français par « Monsieur », « Madame »).
\item \esp{atentamente} : cordialement (clôture formelle).
\item \esp{adjunto / adjuntar} : ci-joint / joindre (un document à un courriel).
\item \esp{solicitar} : demander officiellement ; \esp{la solicitud} : la demande officielle.
\item \esp{le ruego que...} : je vous prie de... (formule de demande très polie, suivie du subjonctif).
\item \esp{el intercambio escolar} : l'échange scolaire.
\item \esp{la carta de motivación} : la lettre de motivation.
\item \esp{te espero con muchas ganas} : je t'attends avec impatience.
\end{itemize}
\end{definition}

\courssub{Compréhension globale : qui écrit ? à qui ? pourquoi ?}

Avant d'analyser les détails, posons les trois questions fondamentales de tout message écrit : l'\cle{émetteur} (qui écrit ?), le \cle{destinataire} (à qui ?) et l'\cle{objectif} (pourquoi ?).

\begin{exercice}[Questions de compréhension globale]
\begin{enumerate}
\item ¿Quién escribe el primer mensaje? ¿A quién? ¿Para qué?
\item ¿Quién escribe el segundo mensaje? ¿A quién? ¿Para qué?
\item ¿Qué relación existe entre Andrés y Mariam? ¿Y entre el consulado y Mariam?
\item ¿Cuál de los dos mensajes es más largo? ¿Por qué?
\end{enumerate}
\end{exercice}

\begin{corrige}[Compréhension globale]
\begin{enumerate}
\item Andrés escribe el primer mensaje a Mariam, su amiga, para invitarla a pasar las vacaciones en Bogotá. (Andrés écrit à son amie Mariam pour l'inviter à passer les vacances à Bogotá.)
\item El consulado de España escribe el segundo mensaje a Mariam para confirmar la recepción de su solicitud y pedirle unos documentos. (Le consulat écrit pour confirmer la réception de sa demande et lui demander des documents.)
\item Andrés y Mariam son amigos, corresponsales: hay una relación de confianza. Entre el consulado y Mariam hay una relación oficial, administrativa. (Relation d'amitié contre relation officielle.)
\item El primer mensaje es más largo porque Andrés cuenta muchas cosas personales; el segundo es breve y directo porque solo da información práctica. (Le message amical est plus long car il raconte des choses personnelles ; le message formel est bref et direct.)
\end{enumerate}
\end{corrige}

\courssub{Compréhension fine : relever les marqueurs de registre}

Observons maintenant les deux textes de près. Un \cle{marqueur de registre} est un mot, une formule ou une structure grammaticale qui indique si le message est familier ou formel.

\begin{exemplebox}[Les formules d'ouverture et de clôture]
\begin{itemize}
\item \esp{Querida Mariam:} « Chère Mariam » : ouverture amicale, avec le prénom seul.
\item \esp{Un abrazo fuerte} : « Je t'embrasse fort » : clôture amicale.
\item \esp{Estimada señora:} « Madame » : ouverture formelle, avec le titre de politesse.
\item \esp{Le saluda atentamente} : « Veuillez agréer mes salutations distinguées » : clôture formelle.
\end{itemize}
\end{exemplebox}

\begin{exemplebox}[Les verbes de demande et les traitements]
\begin{itemize}
\item \esp{Te escribo para invitarte} : « Je t'écris pour t'inviter » : pronom \esp{te} (tutoiement), verbe simple et direct.
\item \esp{Le escribo para confirmar} : « Je vous écris pour confirmer » : pronom \esp{le} (vouvoiement), tournure soutenue.
\item \esp{Dime pronto si puedes venir} : « Dis-moi vite si tu peux venir » : impératif familier \esp{dime}.
\item \esp{Le ruego que adjunte estos documentos} : « Je vous prie de joindre ces documents » : formule de politesse avec subjonctif.
\end{itemize}
\end{exemplebox}

\begin{attention}[Piège : ne pas traduire mot à mot les formules]
Ne traduis jamais \esp{Estimado señor} par « Estimé monsieur » : en français, dans une lettre formelle, on dit simplement « Monsieur » ou « Madame ». De même, \esp{Le saluda atentamente} ne se traduit pas par « Il vous salue attentivement » mais par « Veuillez agréer mes salutations distinguées » ou « Cordialement ». Les formules de politesse sont des blocs fixes : on les traduit par l'équivalent français, pas par le sens littéral. Autre piège : \esp{solicitar} ne veut pas dire « solliciter » au sens français familier de « importuner », mais « demander officiellement ».
\end{attention}

\courssub{Comparaison des deux textes : deux mondes, deux langues}

Comparons maintenant les deux messages sur plusieurs critères : la longueur des phrases, le ton, le vocabulaire.

\begin{center}
\begin{tabularx}{\linewidth}{|X|X|X|}
\hline
\rowcolor{popblueL} Critère & Carta amistosa & Correo formal \\
\hline
Saludo & \esp{Querida Mariam} & \esp{Estimada señora} \\
\hline
Tratamiento & \esp{tú} (\esp{estás, puedes, dime}) & \esp{usted} (\esp{debe, adjunte, puede}) \\
\hline
Despedida & \esp{Un abrazo fuerte} & \esp{Le saluda atentamente} \\
\hline
Frases & courtes, exclamatives (\esp{¡Hola! ¿Qué tal?}) & longues, complexes (\esp{Le ruego que...}) \\
\hline
Tono & affectif, joyeux, personnel & neutre, précis, administratif \\
\hline
Vocabulario & quotidien (\esp{abrigo, bailar, primos}) & technique (\esp{solicitud, adjuntar, certificado}) \\
\hline
\end{tabularx}
\end{center}

\begin{popfigure}
\begin{tikzpicture}[font=\small]
\node[draw=popblue, fill=popblueL, rounded corners, align=center, text width=5.5cm] (carta) at (0,0)
{\textbf{Carta amistosa}\\[2pt]
Destinatario: un amigo\\
Saludo: \esp{Querido/a...}\\
Tratamiento: \esp{tú}\\
Despedida: \esp{Un abrazo}};
\node[draw=poporange, fill=poporangeL, rounded corners, align=center, text width=5.5cm] (correo) at (8,0)
{\textbf{Correo formal}\\[2pt]
Destinatario: una institución\\
Saludo: \esp{Estimado señor/a}\\
Tratamiento: \esp{usted}\\
Despedida: \esp{Atentamente}};
\draw[-{Stealth}, thick, popmuted] (carta) -- (correo) node[midway, above, align=center]{registros\\opuestos};
\end{tikzpicture}
\legende{Deux colonnes comparatives : la lettre amicale et le courriel formel se distinguent par le destinataire, le salut, le traitement et la formule de clôture.}
\end{popfigure}

\begin{savaistu}[La correspondencia en el mundo hispánico]
En Espagne et en Amérique latine, la lettre formelle reste très codifiée : on utilise souvent \esp{Muy señor mío} (« Monsieur ») ou \esp{A quien corresponda} (« À qui de droit »). En Guinée équatoriale, seul pays africain hispanophone, les élèves apprennent aussi ces formules à l'école : écrire une lettre administrative correcte est une compétence très valorisée. Et au Cameroun, pays bilingue français-anglais, savoir rédiger une correspondance en espagnol ouvre les portes des échanges avec toute l'Amérique latine !
\end{savaistu}

\begin{popfigure}
\begin{tikzpicture}[mindmap, grow cyclic, every node/.style={concept}, concept color=popblueL, text=popdark,
level 1/.append style={level distance=3.2cm, sibling angle=90},
level 2/.append style={level distance=2.2cm, sibling angle=60}]
\node[concept color=popblue!50]{\textbf{La correspondencia}}
child[concept color=popgreenL]{node[align=center]{amicale\\ \esp{Querido/a}\\ \esp{Un abrazo}}}
child[concept color=poporangeL]{node[align=center]{formelle\\ \esp{Estimado señor}\\ \esp{Atentamente}}}
child[concept color=poppurpleL]{node[align=center]{courriel\\ \esp{Asunto}\\ \esp{Adjunto}}}
child[concept color=poppinkL]{node[align=center]{message\\ court, rapide\\ \esp{¡Hola!}}};
\end{tikzpicture}
\legende{Carte mentale : les quatre grands types de correspondance et leurs marqueurs.}
\end{popfigure}

\courssub{Méthode : commenter un texte de correspondance}

Au baccalauréat et en classe, on te demandera souvent de commenter un document écrit. Voici la démarche en trois temps.

\begin{methode}[Commenter un texte de correspondance]
\begin{enumerate}
\item \textbf{Introduction : présenter le document.} Indique le type de texte (lettre amicale, courriel formel), l'émetteur, le destinataire et l'objectif. Exemple : « Ce document est une lettre amicale écrite par Andrés, un jeune Colombien, à son amie camerounaise Mariam, pour l'inviter à passer les vacances à Bogotá. »
\item \textbf{Développement : analyser en deux axes.} Axe 1 : le contenu (que dit le texte ? quelles informations ?). Axe 2 : le registre (comment le texte le dit-il ? relève les formules d'ouverture et de clôture, le traitement \esp{tú}/\esp{usted}, le ton, le vocabulaire).
\item \textbf{Conclusion : bilan et ouverture.} Résume ce que le texte révèle (une amitié, une démarche administrative) et ouvre éventuellement sur une comparaison avec un autre texte ou avec le français.
\end{enumerate}
\end{methode}

\begin{exoresolu}[Commentaire guidé : le contraste des deux registres]
\textbf{Énoncé :} Rédige un court commentaire (en français) comparant les deux messages reçus par Mariam. Suis la méthode : introduction, deux axes, conclusion.
\tcblower
\textbf{Introduction.} Nous étudions deux messages reçus le même jour par Mariam, élève de Première à Yaoundé : une lettre amicale d'Andrés, son correspondant colombien, et un courriel formel du consulat d'Espagne. Les deux textes ont des objectifs différents : inviter et demander des documents.

\textbf{Axe 1 : le contenu.} La lettre d'Andrés est riche en détails personnels : la famille, la cuisine colombienne (\esp{ajiaco, arepas}), les visites, la salsa. Le courriel du consulat, lui, ne contient que des informations pratiques : la liste des documents, la date limite, les horaires du bureau.

\textbf{Axe 2 : le registre.} Andrés tutoie Mariam (\esp{¿Cómo estás? Dime pronto}), ouvre par \esp{Querida Mariam} et clôt par \esp{Un abrazo fuerte} : le ton est chaleureux, avec des exclamations et même un rire (\esp{jaja}). Le consulat vouvoie Mariam (\esp{usted debe, le ruego que adjunte}), ouvre par \esp{Estimada señora} et clôt par \esp{Le saluda atentamente} : le ton est neutre et les phrases sont longues et précises.

\textbf{Conclusion.} Ces deux textes montrent qu'une même langue possède plusieurs registres : l'espagnol d'Andrés est celui de l'amitié, celui du consulat est celui de l'administration. Savoir passer de l'un à l'autre, comme le fait le français entre « Salut ! » et « Veuillez agréer... », est la compétence clé de la correspondance.
\end{exoresolu}

\begin{exercice}[À ton tour]
Releve dans les deux textes : 1) deux formules d'ouverture ; 2) deux formules de clôture ; 3) trois verbes conjugués avec \esp{tú} ; 4) trois verbes conjugués avec \esp{usted} ; 5) deux informations pratiques données par le consulat. Puis réponds : pourquoi Andrés utilise-t-il des points d'exclamation alors que le consulat n'en utilise aucun ?
\end{exercice}

\begin{corrige}[À ton tour]
1) Ouvertures : \esp{Querida Mariam} / \esp{Estimada señora}. 2) Clôtures : \esp{Un abrazo fuerte} / \esp{Le saluda atentamente}. 3) Verbes avec \esp{tú} : \esp{estás, puedes (venir), trae, pregunta}. 4) Verbes avec \esp{usted} : \esp{debe (enviar), adjunte, puede (consultar)}. 5) Informations pratiques : la date limite (\esp{antes del día quince}) et les horaires (\esp{de nueve a catorce horas}). Andrés utilise des exclamations parce qu'il exprime sa joie et son affection ; le consulat reste neutre parce qu'un message administratif ne doit pas montrer d'émotion.
\end{corrige}

\begin{aretenir}[L'essentiel de la section]
\begin{itemize}
\item Tout message écrit a un émetteur, un destinataire et un objectif : c'est le point de départ de tout commentaire.
\item Le registre familier s'ouvre avec \esp{Querido/a}, tutoie (\esp{tú}) et se clôt avec \esp{Un abrazo}.
\item Le registre formel s'ouvre avec \esp{Estimado señor/señora}, vouvoie (\esp{usted}) et se clôt avec \esp{Le saluda atentamente}.
\item Les formules de politesse ne se traduisent pas mot à mot : on cherche l'équivalent français (« Monsieur », « Cordialement »).
\item Pour commenter : introduction (présentation), développement en deux axes (contenu, registre), conclusion.
\end{itemize}
\end{aretenir}

\coursec{Lexique : el mundo de la correspondencia}

Pour rédiger nos propres messages, il faut d'abord maîtriser le vocabulaire de la correspondance. Classons-le en quatre familles.

\begin{definition}[Les mots de la lettre et du courrier]
\begin{itemize}
\item \esp{la carta} : la lettre ; \esp{el sobre} : l'enveloppe ; \esp{el sello} : le timbre.
\item \esp{el remitente} : l'expéditeur ; \esp{el destinatario} : le destinataire.
\item \esp{la dirección} : l'adresse ; \esp{el código postal} : le code postal.
\item \esp{el correo} : le courrier, la poste ; \esp{enviar / mandar una carta} : envoyer une lettre.
\item \esp{el buzón} : la boîte aux lettres ; \esp{el cartero} : le facteur.
\end{itemize}
\end{definition}

\begin{definition}[Les mots du courriel et du message]
\begin{itemize}
\item \esp{el correo electrónico} : le courriel ; \esp{la dirección de correo} : l'adresse électronique.
\item \esp{el asunto} : l'objet (du courriel) ; \esp{el archivo adjunto} : la pièce jointe.
\item \esp{enviar / recibir / responder un correo} : envoyer / recevoir / répondre à un courriel.
\item \esp{el mensaje (de texto)} : le message (SMS) ; \esp{la bandeja de entrada} : la boîte de réception.
\item \esp{adjuntar un documento} : joindre un document ; \esp{descargar} : télécharger.
\end{itemize}
\end{definition}

\begin{definition}[Les formules de la lettre amicale]
\begin{itemize}
\item Ouvertures : \esp{Querido Andrés:} « Cher Andrés » ; \esp{Querida Mariam:} « Chère Mariam » ; \esp{¡Hola, Mariam!} « Salut, Mariam ! »
\item Premières phrases : \esp{¿Qué tal? ¿Cómo estás?} « Comment vas-tu ? » ; \esp{Gracias por tu carta.} « Merci pour ta lettre. » ; \esp{Te escribo para...} « Je t'écris pour... »
\item Clôtures : \esp{Un abrazo (fuerte)} « Je t'embrasse (fort) » ; \esp{Un beso} « Un baiser » ; \esp{Con cariño} « Affectueusement » ; \esp{Escríbeme pronto} « Écris-moi vite » ; \esp{Hasta pronto} « À bientôt ».
\end{itemize}
\end{definition}

\begin{definition}[Les formules de la lettre formelle]
\begin{itemize}
\item Ouvertures : \esp{Estimado señor:} « Monsieur » ; \esp{Estimada señora:} « Madame » ; \esp{Muy señor mío:} « Monsieur » (très formel) ; \esp{A quien corresponda:} « À qui de droit ».
\item Premières phrases : \esp{Le escribo para solicitar...} « Je vous écris pour demander... » ; \esp{Me dirijo a usted para...} « Je m'adresse à vous pour... » ; \esp{En respuesta a su carta...} « En réponse à votre lettre... »
\item Demandes polies : \esp{Le ruego que + subjonctif} « Je vous prie de... » ; \esp{Le agradecería que + subjonctif} « Je vous serais reconnaissant(e) de... »
\item Clôtures : \esp{Le saluda atentamente} « Veuillez agréer mes salutations distinguées » ; \esp{Atentamente} « Cordialement » ; \esp{Cordialmente} « Cordialement » ; \esp{Sin otro particular} « Sans autre objet » (formule administrative).
\end{itemize}
\end{definition}

\begin{attention}[Faux amis et pièges de vocabulaire]
\begin{itemize}
\item \esp{la carta} signifie « la lettre » mais aussi « la carte » (au restaurant, \esp{la carta} = le menu !). « La carte » géographique se dit \esp{el mapa}, et la carte à jouer \esp{la carta} ou \esp{el naipe}.
\item \esp{el correo} : c'est le courrier ou la poste, pas « le courrier » au sens de « facteur » : le facteur est \esp{el cartero}.
\item \esp{asunto} signifie « objet, affaire, sujet », pas « assunto » : \esp{el asunto del correo} = l'objet du courriel.
\item Attention au genre : \esp{el destinatario / la destinataria}, \esp{el remitente / la remitente}.
\end{itemize}
\end{attention}

\begin{exercice}[Vocabulaire : à chaque mot sa place]
Complète les phrases avec le mot qui convient : \esp{asunto, sello, remitente, adjunto, buzón, destinatario}.
\begin{enumerate}
\item Para enviar una carta, necesitas un \ldots{} en el sobre.
\item El \ldots{} es la persona que recibe la carta.
\item Escribe tu nombre y tu dirección: tú eres el \ldots{}.
\item En un correo electrónico, el \ldots{} resume el tema del mensaje.
\item Te envío mi certificado de estudios como archivo \ldots{}.
\item Echa la carta en el \ldots{} de la esquina.
\end{enumerate}
\end{exercice}

\begin{corrige}[Vocabulaire : à chaque mot sa place]
1) \esp{sello} (timbre) ; 2) \esp{destinatario} (destinataire) ; 3) \esp{remitente} (expéditeur) ; 4) \esp{asunto} (objet) ; 5) \esp{adjunto} (joint) ; 6) \esp{buzón} (boîte aux lettres).
\end{corrige}

\coursec{Grammaire appliquée : tú, usted et l'impératif de politesse}

\courssub{Tú ou usted : le choix du traitement}

\begin{propriete}[Tú ou usted : le choix du traitement]
En espagnol, le choix entre \esp{tú} (tu) et \esp{usted} (vous) est un marqueur de registre essentiel.
\begin{itemize}
\item \esp{Tú} + verbe à la 2\up{e} personne du singulier : amis, famille, jeunes entre eux. Exemple : \esp{¿Cómo estás?} « Comment vas-tu ? »
\item \esp{Usted} + verbe à la 3\up{e} personne du singulier : inconnus, administrations, personnes âgées ou supérieurs hiérarchiques. Exemple : \esp{Usted debe enviar los documentos.} « Vous devez envoyer les documents. »
\item Les pronoms objets changent aussi : \esp{te} (familier) devient \esp{le} ou \esp{lo/la} (formel). \esp{Te espero} « je t'attends » ; \esp{Le saluda atentamente} « il vous salue cordialement ».
\item Les possessifs suivent le même mouvement : \esp{tu carta} « ta lettre » devient \esp{su carta} « votre lettre ».
\end{itemize}
\end{propriete}

\begin{center}
\begin{tabularx}{\linewidth}{|X|X|X|}
\hline
\rowcolor{popblueL} Français & Español (tú) & Español (usted) \\
\hline
Comment vas-tu / allez-vous ? & \esp{¿Cómo estás?} & \esp{¿Cómo está usted?} \\
\hline
Tu peux / vous pouvez venir. & \esp{Puedes venir.} & \esp{Usted puede venir.} \\
\hline
Je t'écris / je vous écris. & \esp{Te escribo.} & \esp{Le escribo.} \\
\hline
Ta / votre lettre & \esp{tu carta} & \esp{su carta} \\
\hline
Dis-moi / dites-moi & \esp{Dime.} & \esp{Dígame.} \\
\hline
\end{tabularx}
\end{center}

\begin{attention}[Usted : la 3\up{e} personne qui veut dire « vous »]
Le piège classique du francophone : conjuguer le verbe avec \esp{usted} comme avec \esp{tú}. Non ! \esp{Usted} se construit toujours avec la 3\up{e} personne du singulier : on dit \esp{usted tiene} « vous avez » (et jamais \esp{usted tienes}), \esp{usted puede} « vous pouvez » (et jamais \esp{usted puedes}). Astuce : imagine que tu parles de « il/elle » tout en t'adressant à la personne.
\end{attention}

\courssub{L'impératif familier : donner des conseils à un ami}

Dans sa lettre, Andrés écrit : \esp{trae un abrigo} « apporte un manteau », \esp{Dime pronto} « dis-moi vite », \esp{Pregunta a tus padres} « demande à tes parents ». Ce sont des \cle{impératifs affirmatifs} à la personne \esp{tú}.

\begin{propriete}[L'impératif affirmatif avec tú]
\begin{itemize}
\item \textbf{Formation régulière :} on prend la 3\up{e} personne du singulier du présent de l'indicatif. \esp{hablar} $\rightarrow$ \esp{habla} « parle » ; \esp{escribir} $\rightarrow$ \esp{escribe} « écris » ; \esp{venir} $\rightarrow$ \esp{ven} (irrégulier, voir ci-dessous).
\item \textbf{Huit impératifs irréguliers à connaître par cœur :}
\end{itemize}
\begin{center}
\begin{tabular}{|l|l|l|}
\hline
\rowcolor{popgreenL} Infinitivo & Imperativo (tú) & Français \\
\hline
\esp{decir} & \esp{di} & dis \\
\hline
\esp{hacer} & \esp{haz} & fais \\
\hline
\esp{ir} & \esp{ve} & va \\
\hline
\esp{poner} & \esp{pon} & mets \\
\hline
\esp{salir} & \esp{sal} & sors \\
\hline
\esp{ser} & \esp{sé} & sois \\
\hline
\esp{tener} & \esp{ten} & aie / tiens \\
\hline
\esp{venir} & \esp{ven} & viens \\
\hline
\end{tabular}
\end{center}
\begin{itemize}
\item \textbf{Pronoms :} ils se placent après l'impératif affirmatif et s'y soudent, avec un accent pour garder la prononciation : \esp{dime} « dis-moi », \esp{escríbeme} « écris-moi », \esp{mándame la fecha} « envoie-moi la date ».
\end{itemize}
\end{propriete}

\begin{propriete}[L'impératif négatif et l'impératif avec usted : le subjonctif]
\begin{itemize}
\item \textbf{Impératif négatif (tú) :} \esp{no} + subjonctif présent. \esp{No escribas} « n'écris pas » ; \esp{No vengas tarde} « ne viens pas en retard » ; \esp{No me digas} « ne me dis pas » (le pronom se place avant).
\item \textbf{Impératif de politesse (usted) :} subjonctif présent, affirmatif ou négatif. \esp{Escriba} « écrivez » ; \esp{Adjunte los documentos} « joignez les documents » ; \esp{No dude en llamarnos} « n'hésitez pas à nous appeler » ; \esp{Envíe la carta antes del día quince} « envoyez la lettre avant le quinze ».
\item Rappel de formation du subjonctif présent : radical de la 1\up{re} personne du présent + terminaisons inversées (\esp{-e} pour \esp{-ar}, \esp{-a} pour \esp{-er/-ir}) : \esp{hablar} $\rightarrow$ \esp{hable} ; \esp{tener} $\rightarrow$ \esp{tenga} ; \esp{escribir} $\rightarrow$ \esp{escriba}.
\end{itemize}
\end{propriete}

\begin{exemplebox}[Du familier au formel : la même demande, deux registres]
\begin{itemize}
\item \esp{Escríbeme pronto.} « Écris-moi vite. » (à un ami) $\rightarrow$ \esp{Escríbame antes del día quince.} « Écrivez-moi avant le quinze. » (à une administration)
\item \esp{Mándame tu dirección.} « Envoie-moi ton adresse. » $\rightarrow$ \esp{Envíe su dirección, por favor.} « Envoyez votre adresse, s'il vous plaît. »
\item \esp{No olvides tu abrigo.} « N'oublie pas ton manteau. » $\rightarrow$ \esp{No olvide adjuntar los documentos.} « N'oubliez pas de joindre les documents. »
\end{itemize}
\end{exemplebox}

\begin{exoresolu}[Thème guidé : transformer un message]
\textbf{Énoncé :} Transforme ces phrases adressées à un ami en phrases formelles adressées à un consul. 1) \esp{Mándame los documentos.} 2) \esp{No escribas muy tarde.} 3) \esp{Dime tu fecha de llegada.}
\tcblower
1) \esp{Envíeme los documentos, por favor.} « Envoyez-moi les documents, s'il vous plaît. » (impératif de politesse : subjonctif \esp{envíe} + pronom \esp{me} soudé, d'où l'accent). 2) \esp{No escriba muy tarde.} « N'écrivez pas trop tard. » (négation + subjonctif \esp{escriba}). 3) \esp{Dígame su fecha de llegada.} « Dites-moi votre date d'arrivée. » (\esp{decir} : subjonctif \esp{diga} + \esp{me} $\rightarrow$ \esp{dígame} ; possessif formel \esp{su}).
\end{exoresolu}

\begin{exercice}[Grammaire : à toi de conjuguer]
\begin{enumerate}
\item Conjugue à l'impératif familier (\esp{tú}) : \esp{venir} (viens), \esp{escribir} (écris), \esp{hacer} (fais), \esp{poner} (mets).
\item Mets les mêmes verbes à l'impératif de politesse (\esp{usted}).
\item Traduis en espagnol : « Écris-moi vite. » (à un ami) ; « Joignez votre passeport. » (formel) ; « Ne viens pas en retard. » (à un ami).
\end{enumerate}
\end{exercice}

\begin{corrige}[Grammaire : à toi de conjuguer]
1) \esp{ven, escribe, haz, pon}. 2) \esp{venga, escriba, haga, ponga}. 3) \esp{Escríbeme pronto.} ; \esp{Adjunte su pasaporte.} ; \esp{No vengas tarde.}
\end{corrige}

\coursec{Structure de la lettre et du courriel}

\courssub{La lettre : cinq parties obligatoires}

Toute lettre espagnole, amicale ou formelle, suit la même architecture en cinq parties. Apprends-la comme un plan fixe : c'est la grille de lecture et de rédaction.

\begin{propriete}[Les cinq parties de la lettre]
\begin{enumerate}
\item \textbf{Lugar y fecha} (lieu et date), en haut à droite : \esp{Yaundé, a 12 de mayo de 2025}. En espagnol, la date s'écrit : ville, virgule, \esp{a} + article + jour + \esp{de} + mois + \esp{de} + année. Le mois ne prend pas de majuscule.
\item \textbf{Saludo} (salut d'ouverture), suivi de \textbf{deux-points} : \esp{Querido Andrés:} ou \esp{Estimado señor:}. Attention : en espagnol, après les deux-points, on passe à la ligne et on commence par une majuscule.
\item \textbf{Cuerpo} (corps de la lettre) : deux ou trois paragraphes. Premier paragraphe : salutations et raison d'écrire (\esp{¿Cómo estás? Te escribo para...}). Paragraphes suivants : les nouvelles, la demande, l'invitation.
\item \textbf{Despedida} (formule de clôture) : \esp{Un abrazo} ou \esp{Le saluda atentamente}, suivie d'une virgule.
\item \textbf{Firma} (signature) : le prénom (lettre amicale) ou le nom complet et la fonction (lettre formelle).
\end{enumerate}
\end{propriete}

\begin{popfigure}
\begin{tikzpicture}[font=\small]
\node[draw=popblue, fill=popblueL, rounded corners, align=left, text width=12.5cm] (l1) at (0,0) {\textbf{1. Lugar y fecha} : \esp{Yaundé, a 12 de mayo de 2025}};
\node[draw=popgreen, fill=popgreenL, rounded corners, align=left, text width=12.5cm, below=0.15cm of l1] (l2) {\textbf{2. Saludo} : \esp{Querido Andrés:} / \esp{Estimado señor:}};
\node[draw=poporange, fill=poporangeL, rounded corners, align=left, text width=12.5cm, below=0.15cm of l2] (l3) {\textbf{3. Cuerpo} : saludos + razón de escribir ; noticias ; petición o invitación};
\node[draw=poppurple, fill=poppurpleL, rounded corners, align=left, text width=12.5cm, below=0.15cm of l3] (l4) {\textbf{4. Despedida} : \esp{Un abrazo,} / \esp{Le saluda atentamente,}};
\node[draw=poppink, fill=poppinkL, rounded corners, align=left, text width=12.5cm, below=0.15cm of l4] (l5) {\textbf{5. Firma} : \esp{Mariam} / \esp{Mariam Ngo Bell, alumna de Primera A}};
\end{tikzpicture}
\legende{Schéma de la lettre : cinq parties dans un ordre fixe, du lieu-date à la signature.}
\end{popfigure}

\courssub{Le courriel : la même structure, plus l'objet}

Le courriel reprend la structure de la lettre, avec une case supplémentaire : l'\esp{asunto} (l'objet), une ligne courte qui résume le sujet, par exemple \esp{Asunto: Solicitud de información sobre el intercambio escolar} « Objet : demande d'informations sur l'échange scolaire ». Le message sur téléphone portable, lui, est beaucoup plus libre : phrases courtes, abréviations, ton direct — mais il reste réservé aux amis.

\begin{attention}[Les erreurs de structure à éviter]
\begin{itemize}
\item N'oublie pas les \textbf{deux-points} après le salut : \esp{Querido Andrés:} (jamais de virgule comme en français).
\item Après les deux-points, la première phrase commence par une \textbf{majuscule} sur une nouvelle ligne.
\item La date espagnole : \esp{a 12 de mayo de 2025}, avec \esp{de} entre les éléments et sans majuscule au mois. Ne calque pas « le 12 mai 2025 ».
\item Dans une lettre formelle, ne mélange jamais les registres : pas de \esp{¡Hola!} ni de \esp{Un abrazo} après un \esp{Estimado señor}.
\end{itemize}
\end{attention}

\coursec{Production guidée : escribo a mi correspondiente}

\courssub{Méthode : rédiger une lettre amicale en 80 mots}

\begin{methode}[Rédiger une lettre amicale (environ 80 mots)]
\begin{enumerate}
\item \textbf{Prépare le plan.} Note les cinq parties : lieu et date, salut, corps (2 ou 3 paragraphes), clôture, signature.
\item \textbf{Rédige le corps.} Paragraphe 1 : salutations et raison d'écrire (\esp{¿Cómo estás? Te escribo para...}). Paragraphe 2 : deux ou trois informations ou une invitation, avec un ou deux impératifs (\esp{ven, trae, dime}). Paragraphe 3 : une question et une phrase d'attente (\esp{Te espero con muchas ganas.}).
\item \textbf{Vérifie le registre.} Tutoiement partout (\esp{tú, te, tu}), formules \esp{Querido/a} et \esp{Un abrazo}, ponctuation espagnole ¿ ¡.
\item \textbf{Compte et corrige.} Environ 80 mots ; vérifie les accents, la concordance des genres et la date espagnole.
\end{enumerate}
\end{methode}

\begin{textebox}[Modelo : la respuesta de Mariam a Andrés]
Yaundé, a 20 de mayo de 2025

Querido Andrés:

¡Hola! ¿Qué tal? Estoy muy bien, gracias por tu carta. Me hace mucha ilusión tu invitación: ¡quiero conocer Bogotá y a tu familia!

Mis padres dicen que sí. Puedo viajar en agosto, durante las vacaciones. Dime qué ropa debo llevar: aquí siempre hace calor, así que no tengo muchos abrigos. También quiero probar el ajiaco y bailar salsa con tus primos.

Mándame la dirección de tu casa y el número de teléfono de tu madre. Te espero... ¡no, tú me esperas a mí, jaja!

Un abrazo fuerte,

Mariam
\end{textebox}

\textbf{Traduction du modèle :} « Yaoundé, le 20 mai 2025. Cher Andrés : Salut ! Comment ça va ? Je vais très bien, merci pour ta lettre. Ton invitation me fait très plaisir : je veux connaître Bogotá et ta famille ! Mes parents disent oui. Je peux voyager en août, pendant les vacances. Dis-moi quels vêtements je dois apporter : ici il fait toujours chaud, alors je n'ai pas beaucoup de manteaux. Je veux aussi goûter l'ajiaco et danser la salsa avec tes cousins. Envoie-moi l'adresse de ta maison et le numéro de téléphone de ta mère. Je t'attends... non, c'est toi qui m'attends, haha ! Je t'embrasse fort, Mariam. »

\begin{definition}[Glossaire du modèle]
\begin{itemize}
\item \esp{me hace mucha ilusión} : cela me fait très plaisir (tournure très courante en Espagne).
\item \esp{la ropa} : les vêtements ; \esp{llevar} : porter, apporter.
\item \esp{así que} : donc, alors ; \esp{también} : aussi.
\end{itemize}
\end{definition}

\begin{exoresolu}[Version guidée : traduire une lettre]
\textbf{Énoncé :} Traduis en espagnol : « Cher Pablo : Salut ! Comment vas-tu ? Je t'écris pour t'inviter à Yaoundé pendant les vacances. Tu peux visiter le marché central et jouer au football avec mes amis. Apporte un appareil photo ! Dis-moi vite ta date d'arrivée. Je t'embrasse, Carlos. »
\tcblower
\esp{Querido Pablo: ¡Hola! ¿Cómo estás? Te escribo para invitarte a Yaundé durante las vacaciones. Puedes visitar el mercado central y jugar al fútbol con mis amigos. ¡Trae una cámara de fotos! Dime pronto tu fecha de llegada. Un abrazo, Carlos.} Points de vigilance : deux-points après le salut ; \esp{invitarte a} + nom de ville ; \esp{jugar al fútbol} (et non \esp{jugar fútbol}) ; impératifs \esp{trae} et \esp{dime} ; clôture \esp{Un abrazo}.
\end{exoresolu}

\begin{experience}[Tarea final : tu carta a un corresponsal hispanohablante]
À ton tour ! Rédige une lettre amicale d'environ 80 mots à un correspondant espagnol ou latino-américain (imaginaire ou réel). Tu dois : 1) écrire le lieu et la date en espagnol ; 2) ouvrir avec \esp{Querido/a} ; 3) te présenter brièvement et donner une nouvelle de ton lycée ou de ta ville ; 4) inviter ton correspondant au Cameroun avec au moins un impératif (\esp{ven, trae, dime}) ; 5) poser une question ; 6) clore avec \esp{Un abrazo} et signer. Variante pour les plus rapides : rédige ensuite le courriel formel (5 lignes) que tu enverrais au consulat pour demander des informations sur un échange scolaire, avec \esp{Estimado señor:}, \esp{Le escribo para solicitar...} et \esp{Le saluda atentamente}.
\end{experience}

\begin{aretenir}[L'essentiel de la leçon]
\begin{itemize}
\item La correspondance espagnole repose sur deux registres : familier (\esp{Querido/a}, \esp{tú}, \esp{Un abrazo}) et formel (\esp{Estimado señor/a}, \esp{usted}, \esp{Le saluda atentamente}).
\item \esp{Usted} se conjugue à la 3\up{e} personne du singulier ; les pronoms et possessifs changent aussi (\esp{te} $\rightarrow$ \esp{le}, \esp{tu} $\rightarrow$ \esp{su}).
\item Impératif familier : 3\up{e} personne du présent (\esp{escribe}) et huit irréguliers (\esp{di, haz, ve, pon, sal, sé, ten, ven}) ; impératif de politesse et impératif négatif : subjonctif (\esp{escriba, no escribas}).
\item Structure fixe de la lettre : lieu et date, salut suivi de deux-points, corps en paragraphes, clôture, signature ; le courriel ajoute l'\esp{asunto}.
\item Les formules de politesse se traduisent par leur équivalent français, jamais mot à mot.
\end{itemize}
\end{aretenir}

\coursec{Lexique : el mundo de la correspondencia}

Après avoir analysé dans la section précédente deux messages aux registres radicalement différents — un message \emph{informel} entre amis et un courriel \emph{formel} adressé à une administration —, nous disposons maintenant du matériel brut pour bâtir notre répertoire lexical. Cette section a pour but de fixer, classer et enrichir le vocabulaire indispensable pour rédiger toute correspondance en espagnol, qu'il s'agisse d'une lettre manuscrite déposée dans un \esp{buzón} à Yaoundé ou d'un \esp{correo electrónico} envoyé depuis un smartphone à Buenos Aires.

\courssub{Le vocabulaire de la lettre manuscrite (la carta tradicional)}

Commençons par observer les éléments matériels qui composent une lettre physique. Dans le texte support, l'ami écrivait : \esp{« Te mando esta carta por correo ordinario »} (Je t'envoie cette lettre par la poste ordinaire). Décomposons la chaîne logistique.

\begin{definition}[El registro : le niveau de langue]
Le \cle{registro} (registre de langue) est l'ensemble des choix lexicaux, grammaticaux et stylistiques que l'on adapte en fonction du destinataire, du contexte et de l'intention communicative. On distingue schématiquement :
\begin{itemize}
    \item \textbf{Registro informal / familiar :} usado con amigos, familia, compañeros de clase (tú, abreviaturas, coloquialismos).
    \item \textbf{Registro formal :} usado con desconocidos, autoridades, empresas, profesores (usted/ustedes, fórmulas fijas, vocabulario culto).
\end{itemize}
Chaque mot de vocabulaire ci-dessous porte une « étiquette » de registre.
\end{definition}

\begin{center}
\begin{tabularx}{\linewidth}{|X|X|X|}
\hline
\rowcolor{popblueL}
\textbf{Español} & \textbf{Français} & \textbf{Remarque / Registre} \\
\hline
\esp{la carta} & la lettre (missive) & Neutre. Attention au faux ami : \esp{una letra} = une lettre de l'alphabet. \\
\hline
\esp{el sobre} & l'enveloppe & Neutre. Verbe associé : \esp{sobrescribir} (adresser l'enveloppe). \\
\hline
\esp{el sello} / \esp{la estampilla} & le timbre-poste & \esp{Sello} (Espagne), \esp{estampilla} (Amérique latine, Guinée équatoriale). \\
\hline
\esp{el remitente} & l'expéditeur & Qui \emph{envoie} (verbe \esp{remitir}). Se place au dos de l'enveloppe. \\
\hline
\esp{el destinatario} & le destinataire & Qui \emph{reçoit} (verbe \esp{destinar}). Sur la face avant. \\
\hline
\esp{la dirección} & l'adresse & \esp{Calle, número, piso, puerta}. \\
\hline
\esp{el código postal} & le code postal & 5 chiffres en Espagne (ex. 28001 Madrid), 5 ou 4 selon pays. \\
\hline
\esp{el buzón} & la boîte aux lettres & Public (jaune en Espagne, rouge au Cameroun) ou privé. \\
\hline
\esp{el cartero} / \esp{la cartera} & le facteur / la factrice & Métier. \esp{Repartidor} pour les colis. \\
\hline
\esp{el correo ordinario} & le courrier ordinaire & Sans suivi, délai standard. \\
\hline
\esp{el correo certificado} & le courrier recommandé & Avec accusé de réception (\esp{acuse de recibo}). \\
\hline
\end{tabularx}
\end{center}

\begin{exemplebox}[Exemples traduits : les acteurs de la correspondance]
\begin{itemize}
    \item \esp{El remitente escribe la dirección en el sobre.} $\rightarrow$ L'expéditeur écrit l'adresse sur l'enveloppe.
    \item \esp{El destinatario recibe la carta tres días después.} $\rightarrow$ Le destinataire reçoit la lettre trois jours plus tard.
    \item \esp{Olvidé poner el sello y la carta volvió a casa.} $\rightarrow$ J'ai oublié de mettre le timbre et la lettre est revenue à la maison.
    \item \esp{El código postal de Douala es 00237 (fictif pour l'exemple).} $\rightarrow$ Le code postal de Douala est...
\end{itemize}
\end{exemplebox}

\begin{attention}[Faux ami classique : \texorpdfstring{\esp{carta}}{carta} vs \texorpdfstring{\esp{letra}}{letra}]
C'est l'erreur n°1 des francophones !
\begin{itemize}
    \item \esp{Una carta} = une lettre (que l'on met dans une enveloppe). \textbf{Synonyme :} \esp{una misiva}.
    \item \esp{Una letra} = une lettre de l'alphabet (a, b, c...). Aussi : la parole d'une chanson (\esp{la letra de la canción}) ou un effet de commerce (\esp{letra de cambio}).
\end{itemize}
\textbf{Mnémotechnique :} Dans \esp{ca\underline{r}ta}, il y a un \textbf{R} comme dans Envelope\underline{r} / Envoye\underline{r}. Dans \esp{le\underline{t}ra}, il y a un \textbf{T} comme dans Alpha\underline{t} / Alphabe\underline{t}.
\end{attention}

\textbf{Le parcours d'une lettre : frise chronologique.}
Voici le cycle de vie d'une lettre physique, de la plume à la main du destinataire.

\begin{popfigure}
\begin{tikzpicture}[
    node distance=1.5cm and 0.8cm,
    box/.style={draw=popblue, thick, rounded corners, fill=popblueL, minimum height=1.6cm, minimum width=2.6cm, align=center, font=\small},
    arrow/.style={-Stealth, thick, popdark}
]
% Nodes
\node[box, align=center] (e1){\textbf{1. Escribir}\\\esp{La carta}\\(con boli/pluma)};
\node[box, right=of e1, align=center] (e2){\textbf{2. Meter en el sobre}\\\esp{El sobre}\\(\esp{doblar} la carta)};
\node[box, right=of e2, align=center] (e3){\textbf{3. Poner el sello}\\\esp{El sello / estampilla}\\(\esp{Franquear})};
\node[box, right=of e3, align=center] (e4){\textbf{4. Echar al buzón}\\\esp{El buzón}\\(\esp{Correos})};
\node[box, right=of e4, align=center] (e5){\textbf{5. El cartero}\\\esp{Recoge / Reparte}\\(\esp{Repartidor})};
\node[box, right=of e5, align=center] (e6){\textbf{6. El destinatario}\\\esp{Recibe la carta}\\(\esp{Abre el sobre})};

% Arrows
\draw[arrow] (e1) -- (e2);
\draw[arrow] (e2) -- (e3);
\draw[arrow] (e3) -- (e4);
\draw[arrow] (e4) -- (e5);
\draw[arrow] (e5) -- (e6);
\end{tikzpicture}
\legende{Le parcours d'une lettre manuscrite : de la rédaction à la réception. Verbes clés : \esp{escribir, doblar, franquear, echar, recoger, repartir, recibir, abrir}.}
\end{popfigure}

\courssub{Le vocabulaire du courriel (el correo electrónico)}

Le texte support montrait un échange par \esp{email}. Le vocabulaire technique est largement transparent pour un francophone (anglicismes partagés), mais le genre des noms et les verbes associés piègent souvent.

\begin{center}
\begin{tabularx}{\linewidth}{|X|X|X|}
\hline
\rowcolor{poppurpleL}
\textbf{Español} & \textbf{Français} & \textbf{Remarque / Verbe associé} \\
\hline
\esp{el correo electrónico} / \esp{el email} / \esp{el e-mail} & le courriel / l'email & Masculin. \esp{Correo-e} (abrév. admin). \\
\hline
\esp{la dirección de correo} / \esp{el email} & l'adresse email & \esp{Mi email es juan@ejemplo.com}. \\
\hline
\esp{el asunto} & l'objet (du mail) & \textbf{Obligatoire} en formel. Court et clair. \\
\hline
\esp{el mensaje} / \esp{el cuerpo del mensaje} & le message / le corps du texte & Ne pas confondre avec \esp{el mensajero} (le messager). \\
\hline
\esp{el archivo adjunto} / \esp{el adjunto} & la pièce jointe & \esp{Adjunto} = adjectif $\rightarrow$ nom masculin. \textbf{Pas de « la pièce jointe » !} \\
\hline
\esp{enviar} / \esp{mandar} & envoyer & \esp{Enviar un correo}. \esp{Mandar} plus familier. \\
\hline
\esp{recibir} & recevoir & \esp{He recibido tu email}. \\
\hline
\esp{responder} / \esp{contestar} & répondre & \esp{Responder a alguien}. \textbf{Préposition \esp{a} obligatoire.} \\
\hline
\esp{reenviar} / \esp{reenvío} & transférer (forward) & \esp{Reenvío el correo a mi jefe}. \\
\hline
\esp{adjuntar} & joindre (un fichier) & \esp{Adjunto mi CV}. (Verbe transitif direct). \\
\hline
\esp{la bandeja de entrada} & la boîte de réception & \esp{Entrada} = entrée. \esp{Bandeja de salida} = éléments envoyés. \\
\hline
\esp{el spam} / \esp{el correo basura} & le courrier indésirable & \esp{Carpeta de spam}. \\
\hline
\esp{la firma} & la signature (bloc final) & Automatique : \esp{firma automática}. \\
\hline
\end{tabularx}
\end{center}

\begin{exoresolu}[Choisir le bon verbe et la bonne préposition]
\textbf{Énoncé :} Complétez avec le verbe adapté (\esp{enviar, adjuntar, responder, reenviar, recibir}) à la bonne forme (infinitif ou conjugué présent) et la préposition si nécessaire.
\begin{enumerate}
    \item Ayer \_\_\_\_\_\_ un correo a mi profesor con los deberes.
    \item Tengo que \_\_\_\_\_\_ el archivo PDF \_\_\_\_\_\_ el correo.
    \item ¿Ya has \_\_\_\_\_\_ \_\_\_\_\_\_ ese mensaje tan importante?
    \item Por favor, \_\_\_\_\_\_\_\_\_\_ este correo al director.
    \item No olvides \_\_\_\_\_\_\_\_\_\_ \_\_\_\_\_\_ mi pregunta.
\end{enumerate}

\tcblower
\textbf{Solution détaillée :}
\begin{enumerate}
    \item \esp{envié} (ou \esp{mandé}). \emph{Verbe transitif direct : envoyer quelque chose à quelqu'un $\rightarrow$ \esp{enviar algo a alguien}.}
    \item \esp{adjuntar / al} (ou \esp{en}). \emph{\esp{Adjuntar algo a/en algo}. Attention : \esp{adjuntar} est transitif direct pour le fichier, prépositionnel pour le support.}
    \item \esp{respondido / a}. \emph{\esp{Responder A alguien / A algo}. Erreur fréquente : *\esp{responder alguien} (calque du français « répondre à quelqu'un » sans « à » mais en espagnol la préposition \esp{a} reste).}
    \item \esp{reenviar} (ou \esp{reenvíe} impératif formel / \esp{reenvía} informel). \emph{Transitif direct : \esp{reenviar algo a alguien}.}
    \item \esp{contestar / a} (ou \esp{responder / a}). \emph{\esp{Contestar a alguien}.}
\end{enumerate}
\end{exoresolu}

\courssub{Les formules d'ouverture (las fórmulas de saludo / apertura)}

C'est le premier marqueur de registre. Dans le texte support, l'ami écrivait \esp{« Hola, Patri: »} alors que l'administration utilisait \esp{« Estimado Sr. Director: »}. Observons la règle de ponctuation cruciale.

\begin{savaistu}[Les deux-points \texorpdfstring{\esp{:}}{:} après la formule d'ouverture]
En espagnol (Espagne, Amérique latine, Guinée équatoriale), la norme typographique standard impose \textbf{les deux-points} après la formule d'appel, et non la virgule comme en français.
\begin{itemize}
    \item \textbf{Correct :} \esp{Querida Ana:} \quad \esp{Estimado Sr. López:} \quad \esp{Hola, Pedro:}
    \item \textbf{Incorrect (calque français) :} \esp{Querida Ana,} \quad \esp{Estimado Sr. López,}
\end{itemize}
Après les deux-points, on passe à la ligne et on commence le corps du message \textbf{avec une majuscule}.
\esp{Estimado Sr. López:} \\
\esp{Le escribo para informarle...}
\end{savaistu}

\paragraph{A. Formules amicales / informelles (Registro informal)}
Utilisées avec \esp{tú} (ou \esp{vos} en zone \esp{voseante} : Argentine, Uruguay, partie de l'Amérique centrale). Le prénom est obligatoire.

\begin{center}
\begin{tabularx}{\linewidth}{|X|X|X|}
\hline
\rowcolor{popgreenL}
\textbf{Fórmula} & \textbf{Traduction / Contexte} & \textbf{Nuance} \\
\hline
\esp{Querido [nombre] :} / \esp{Querida [nombre] :} & Cher / Chère [prénom] & Standard affectueux. Accord en genre. Pluriel : \esp{Queridos amigos:}. \\
\hline
\esp{Hola [nombre] :} & Salut [prénom] & Très courant, neutre-amical. SMS, emails courts. \\
\hline
\esp{¿Qué tal [nombre]? :} / \esp{¿Cómo estás? :} & Comment vas-tu ? & Ouverture dialoguée, attend une réponse (rhétorique ou non). \\
\hline
\esp{¡Hola, [nombre]!} & Salut [prénom] ! & Avec points d'exclamation (ton enthousiaste). Pas de deux-points si point d'excl. \\
\hline
\esp{Buenos días / tardes, [nombre] :} & Bonjour / Bonsoir [prénom] & Politesse de base, légèrement plus distant que \esp{Hola}. \\
\hline
\end{tabularx}
\end{center}

\paragraph{B. Formules formelles (Registro formal)}
Utilisées avec \esp{usted} / \esp{ustedes}. On utilise le titre + nom de famille, ou formules génériques.

\begin{center}
\begin{tabularx}{\linewidth}{|X|X|X|}
\hline
\rowcolor{poporangeL}
\textbf{Fórmula} & \textbf{Traduction / Contexte} & \textbf{Nuance / Grammaire} \\
\hline
\esp{Estimado Sr. [Apellido] :} / \esp{Estimada Sra. [Apellido] :} & Monsieur / Madame [Nom] & \textbf{Standard professionnel.} Accord : \esp{Estimado} (masc), \esp{Estimada} (fém). Abrév. : \esp{Sr.}, \esp{Sra.}, \esp{Dr.}, \esp{Dra.}, \esp{Prof.} \\
\hline
\esp{Muy señor mío :} / \esp{Muy señora mía :} & Monsieur / Madame & Très formel, administratif, juridique. Vieilli mais correct. Sans nom. \\
\hline
\esp{Distinguido Sr. [Apellido] :} / \esp{Distinguida Sra. [Apellido] :} & Distinguished Mr/Mrs & Très cérémonieux, invitations officielles, lettres de recommandation. \\
\hline
\esp{A quien corresponda :} & À qui de droit & Impersonnel. Réclamations, références, lettre de motivation sans nom de contact. \textbf{Jamais de nom après.} \\
\hline
\esp{Estimados señores :} / \esp{Estimadas señoras :} & Messieurs / Mesdames & Pluriel. Adresse à une entreprise, un service, un jury. \\
\hline
\end{tabularx}
\end{center}

\begin{attention}[Pièges grammaticaux des formules formelles]
\begin{enumerate}
    \item \textbf{Accord de l'adjectif :} \esp{Estimad\underline{o} Sr. García} (masc) vs \esp{Estimad\underline{a} Sra. García} (fém). Erreur : *\esp{Estimado Sra. García}.
    \item \textbf{Majuscules de politesse :} \esp{Sr.}, \esp{Sra.}, \esp{Dr.} prennent la majuscule. \esp{usted} / \esp{ustedes} s'écrivent \textbf{en minuscule} dans le corps du texte, mais souvent abrégés \esp{Vd.} / \esp{Vds.} (majuscule) dans la correspondance très formelle ancienne.
    \item \textbf{\esp{A quien corresponda} :} \esp{quien} sans accent (pronom relatif), \esp{corresponda} (subjonctif présent !) car antécédent indéfini/négatif. Ne pas écrire *\esp{A quién corresponde} (indicatif = question directe).
\end{enumerate}
\end{attention}

\courssub{Les formules de clôture (las fórmulas de despedida / cierre)}

Elles scellent le registre. Une lettre formelle signée \esp{Un abrazo} est un non-sens sociolinguistique majeur.

\paragraph{A. Formules amicales (Registro informal)}

\begin{center}
\begin{tabularx}{\linewidth}{|X|X|X|}
\hline
\rowcolor{popgreenL}
\textbf{Fórmula} & \textbf{Traduction} & \textbf{Usage} \\
\hline
\esp{Un abrazo,} / \esp{Un fuerte abrazo,} & Une accolade / Une grosse accolade & Standard universel (amitié, famille proche). \\
\hline
\esp{Un beso,} / \esp{Besos,} / \esp{Muchos besos,} & Un bisou / Bisous / Plein de bisous & Très affectueux (famille, couple, amies très proches). \textbf{Éviter en contexte purement amical masculin ou mixte formel.} \\
\hline
\esp{Tu amigo,} / \esp{Tu amiga,} / \esp{Tus amigos,} & Ton ami / Ta amie / Vos amis & Signature identitaire. \esp{Tu} = possessif (ton), pas \esp{Tú} (pronom). \\
\hline
\esp{Hasta pronto,} / \esp{Hasta luego,} / \esp{Nos vemos,} & À bientôt / À plus tard / On se voit & Ouverture sur la suite. \esp{Nos vemos} très oral. \\
\hline
\esp{Escríbeme pronto,} / \esp{Contéstame,} & Écris-moi vite / Réponds-moi & Incitation à la réponse. Impératif \esp{tú} (accent sur \esp{é} pour \esp{escríbeme}). \\
\hline
\esp{Cuídate,} / \esp{Cuídense,} & Prends soin de toi / Prenez soin de vous & Chaleureux, fin de lettre longue ou période difficile. \\
\hline
\end{tabularx}
\end{center}

\paragraph{B. Formules formelles (Registro formal)}

\begin{center}
\begin{tabularx}{\linewidth}{|X|X|X|}
\hline
\rowcolor{poporangeL}
\textbf{Fórmula} & \textbf{Traduction} & \textbf{Usage / Registre} \\
\hline
\esp{Atentamente,} & Cordialement / Bien à vous & \textbf{La plus utilisée.} Neutre, professionnelle, administrative. Équivalent de « Cordialement ». \\
\hline
\esp{Le saluda atentamente,} / \esp{Les saluda atentamente,} & Je vous salue respectueusement & \textbf{Très formel.} \esp{Le} (COI \esp{usted}) / \esp{Les} (COI \esp{ustedes}). Verbe \esp{saludar} à la 3e pers. sg. Sujet implicite \esp{yo}. \\
\hline
\esp{Cordialmente,} & Cordialement & Légèrement plus chaleureux qu'\esp{Atentamente}. Relations professionnelles suivies. \\
\hline
\esp{Reciba un cordial saludo,} / \esp{Reciban un cordial saludo,} & Recevez mes salutations cordiales & Impératif formel \esp{usted} (\esp{reciba}) / \esp{ustedes} (\esp{reciban}). Très poli. \\
\hline
\esp{Quedo a su disposición,} / \esp{Quedo a la espera de su respuesta,} & Je reste à votre disposition / J'attends votre réponse & Formules de « transition » avant la signature finale. \esp{Quedo} (1e pers. sg). \\
\hline
\esp{Sin otro particular, le saluda atentamente,} & Sans autre sujet, je vous salue... & Style administratif très codifié (fin de lettre officielle). \\
\hline
\end{tabularx}
\end{center}

\begin{methode}[Comment choisir sa formule de clôture ?]
\begin{enumerate}
    \item \textbf{Identifiez le registre :} \esp{tú} (ami) $\rightarrow$ colonne verte. \esp{usted} (inconnu, supérieur, admin) $\rightarrow$ colonne orange.
    \item \textbf{Vérifiez la cohérence :} Ouverture \esp{Estimado Sr.} $\rightarrow$ Fermeture \esp{Atentamente}. Ouverture \esp{Hola Ana:} $\rightarrow$ Fermeture \esp{Un abrazo}.
    \item \textbf{Ponctuation :} En espagnol, la formule de clôture est suivie d'une \textbf{virgule} (pas de point, pas de deux-points).
    \item \textbf{Signature :} Saute une ligne après la virgule, écris ton nom/prénom. En formel : Prénom + Nom. En amical : Prénom (ou surnom).
\end{enumerate}
\end{methode}

\courssub{Verbes utiles pour l'action épistolaire}

Ces verbes structurent le contenu de votre message (exposer un motif, demander, remercier, joindre). Maîtrisez leur construction (transitivité, prépositions).

\begin{center}
\begin{tabularx}{\linewidth}{|X|X|X|X|}
\hline
\rowcolor{poppinkL}
\textbf{Verbo} & \textbf{Signification} & \textbf{Construction / Exemple} & \textbf{Registre} \\
\hline
\esp{escribir} & écrire & \esp{Escribo \textbf{a} mi amigo.} (écrire \textbf{à} qqn) \newline \esp{Escribo \textbf{una carta}.} (écrire qqch) & Neutre \\
\hline
\esp{contestar} / \esp{responder} & répondre & \esp{Contesto \textbf{a} tu correo.} (\textbf{A} qqn/qqch) \newline \esp{Respondo \textbf{a} tu pregunta.} & Neutre \\
\hline
\esp{agradecer} & remercier & \esp{Agradezco \textbf{tu ayuda}.} (TD) \newline \esp{Le agradezco \textbf{su tiempo}.} (COI \esp{le} = \esp{a usted}) & Formel / Poli \\
\hline
\esp{invitar} & inviter & \esp{Te invito \textbf{a} mi cumpleaños.} (Inviter qqn \textbf{à} + inf / événement) & Amical / Formel \\
\hline
\esp{solicitar} & solliciter, demander formellement & \esp{Solicito \textbf{información sobre}...} (TD + \textbf{sobre} / \textbf{de}) \newline \esp{Le solicito \textbf{una cita}.} & Formel (admin, pro) \\
\hline
\esp{informar} & informer & \esp{Le informo \textbf{de} que...} (Informer qqn \textbf{de} qqch) \newline \esp{Informamos \textbf{sobre} el cambio.} & Formel \\
\hline
\esp{adjuntar} & joindre (fichier) & \esp{Adjunto \textbf{mi CV}.} (TD direct, pas de préposition) \newline \esp{Adjunto \textbf{el archivo} \textbf{al} correo.} & Technique / Formel \\
\hline
\esp{esperar} & attendre / espérer & \esp{Espero \textbf{su respuesta}.} (Attendre qqch - TD) \newline \esp{Espero \textbf{que vengas}.} (Espérer que + \textbf{subjonctif}) & Neutre / Formel \\
\hline
\end{tabularx}
\end{center}

\begin{exemplebox}[Ces verbes en action (phrases types pour le Bac)]
\begin{itemize}
    \item \esp{\textbf{Le escribo} para \textbf{solicitar} información sobre los cursos de verano.} (Je vous écris pour demander des infos sur les cours d'été.)
    \item \esp{\textbf{Adjunto} mi currículum y \textbf{espero} su respuesta.} (Je joins mon CV et j'attends votre réponse.)
    \item \esp{\textbf{Agradezco} de antemano su atención.} (Je vous remercie par avance de votre attention.) \emph{Formule magique formelle.}
    \item \esp{\textbf{Te invito} a que vengas a Yaoundé en diciembre.} (Je t'invite à venir à Yaoundé en décembre.) \emph{\esp{invitar a que} + subjonctif.}
    \item \esp{\textbf{Quedo a la espera} de tus noticias.} (J'attends de tes nouvelles.) \emph{Clôture semi-formelle / amicale soutenue.}
\end{itemize}
\end{exemplebox}

\courssub{Mini-liste à mémoriser pour le Bac (Fiche de révision)}

\begin{textebox}[5 formules AMICALES + 5 formules FORMELLES « Prêtes à l'emploi »]
\textbf{AMICALES (Registro informal / \esp{tú})} \\
1. \esp{Querido/a [Nombre]:} \hfill (Ouverture standard) \\
2. \esp{Hola [Nombre]:} \hfill (Ouverture rapide / email) \\
3. \esp{Un abrazo,} \hfill (Clôture standard) \\
4. \esp{Tu amigo/a,} \hfill (Signature identitaire) \\
5. \esp{Escríbeme pronto,} \hfill (Incitation réponse) \\[0.5em]
\textbf{FORMELLES (Registro formal / \esp{usted})} \\
1. \esp{Estimado Sr. / Estimada Sra. [Apellido]:} \hfill (Ouverture standard pro/admin) \\
2. \esp{A quien corresponda:} \hfill (Ouverture sans nom de contact) \\
3. \esp{Atentamente,} \hfill (Clôture standard universelle) \\
4. \esp{Le saluda atentamente,} \hfill (Clôture très respectueuse) \\
5. \esp{Quedo a la espera de su respuesta,} \hfill (Clôture « action » / motivation)
\end{textebox}

\begin{exercice}[Entraînement : Registre et Vocabulaire]
\textbf{1. Classer :} Rangez ces mots dans le bon tableau : \esp{sello, adjunto, buzón, asunto, remitente, spam, sobre, archivo adjunto, cartero, bandeja de entrada}.
\begin{center}
\begin{tabular}{|l|l|}
\hline
\textbf{Carta manuscrita} & \textbf{Correo electrónico} \\
\hline
... & ... \\
\hline
\end{tabular}
\end{center}
\textbf{2. Corriger les erreurs (registre / grammaire / vocabulaire) :}
\begin{enumerate}
    \item \esp{Estimado Sr. Pérez, \textbf{Un abrazo,} Juan}
    \item \esp{Hola Ana: \textbf{Le saluda atentamente,} Pedro}
    \item \esp{El \textbf{remitente} recibe la carta.}
    \item \esp{Adjunto \textbf{a} el correo mi foto.}
    \item \esp{Espero \textbf{que usted contesta} pronto.}
\end{enumerate}
\textbf{3. Micro-production (30 mots) :} Écris un email formel à \esp{Prof. García} pour lui envoyer un devoir en pièce jointe et demander confirmation de réception. Objet : \esp{Entrega de tarea}.
\end{exercice}

\begin{corrige}[Exercice : Registre et Vocabulaire]
\textbf{1. Classification :}
\begin{center}
\begin{tabular}{|l|l|}
\hline
\textbf{Carta manuscrita} & \textbf{Correo electrónico} \\
\hline
el sello, el buzón, el remitente, el sobre, el cartero & el asunto, el adjunto / archivo adjunto, el spam, la bandeja de entrada \\
\hline
\end{tabular}
\end{center}
\textbf{2. Corrections :}
\begin{enumerate}
    \item \textbf{Incohérence registre.} \esp{Estimado Sr. Pérez:} (formel) $\neq$ \esp{Un abrazo} (informel). \textbf{Correct :} \esp{Estimado Sr. Pérez: ... Atentamente, Juan.}
    \item \textbf{Incohérence registre.} \esp{Hola Ana:} (informel) $\neq$ \esp{Le saluda atentamente} (formel). \textbf{Correct :} \esp{Hola Ana: ... Un abrazo, Pedro.}
    \item \textbf{Erreur sémantique.} \esp{El \textbf{destinatario} recibe la carta.} (Le destinataire reçoit). L'expéditeur (\esp{remitente}) \emph{envoie}.
    \item \textbf{Erreur construction verbe.} \esp{Adjunto mi foto \textbf{al} correo.} (\esp{Adjuntar} qqch \textbf{a/en} qqch. \esp{Al} = \esp{a el}). Ou simplement : \esp{Adjunto mi foto.}
    \item \textbf{Erreur mode (subjonctif).} \esp{Espero \textbf{que usted conteste} pronto.} (\esp{Esperar que} + subjonctif présent : \esp{conteste} pour \esp{usted}).
\end{enumerate}
\textbf{3. Proposition de production :}
\begin{textebox}[Correction production]
\esp{Asunto: Entrega de tarea - [Votre Nom]} \\
\esp{Estimado Prof. García:} \\
\esp{Le escribo para entregarle la tarea solicitada. Adjunto el archivo en formato PDF.} \\
\esp{Agradezco de antemano su tiempo y quedo a la espera de su confirmación de recepción.} \\
\esp{Atentamente,} \\
\esp{[Votre Prénom Nom]}
\end{textebox}
\end{corrige}

\begin{aretenir}[Check-list lexicale « Correspondance »]
\begin{itemize}
    \item \textbf{Matériel lettre :} \esp{carta, sobre, sello/estampilla, remitente, destinatario, dirección, código postal, buzón, cartero}.
    \item \textbf{Matériel email :} \esp{correo electrónico, asunto, mensaje, archivo adjunto, enviar, recibir, responder (a), reenviar, adjuntar, bandeja de entrada}.
    \item \textbf{Ouverture :} \textbf{Deux-points (:)} obligatoires. \esp{Querido/a...:} (inf.) vs \esp{Estimado/a Sr./Sra....:} (form.) vs \esp{A quien corresponda:} (anonyme).
    \item \textbf{Clôture :} \textbf{Virgule (,)}. \esp{Un abrazo / Tu amigo/a} (inf.) vs \esp{Atentamente / Le saluda atentamente} (form.).
    \item \textbf{Verbes clés :} \esp{escribir a, contestar a, agradecer, solicitar, informar de, adjuntar, esperar (+ subj. si changement de sujet)}.
    \item \textbf{Faux amis :} \esp{Carta} $\neq$ \esp{Letra}. \esp{Adjunto} (pièce jointe, masc.) $\neq$ \esp{Adjetivo}.
\end{itemize}
\end{aretenir}

\coursec{Grammaire appliquée : tú, usted et l'impératif de politesse}

Après avoir exploré le lexique spécifique de la correspondance dans la section précédente, nous allons maintenant analyser le cœur grammatical qui permet de \textbf{choisir le bon registre} : l'opposition entre \cle{tú} (tutoiement) et \cle{usted} (vouvoiement), ainsi que l'usage de l'\cle{impératif} (affirmatif et négatif) pour formuler des demandes, des invitations, des instructions ou des interdictions. En espagnol, ce choix n'est pas anodin : il structure toute la phrase, du pronom au verbe en passant par les possessifs et les pronoms objets.

\courssub{Observation à partir du texte support}

Reprenons deux phrases extraites du dialogue \emph{Dos mensajes, dos registros} (section 1) pour observer la différence de structure.

\begin{exemplebox}[Extrait du texte support]
\textbf{Message 1 (amical) :} \esp{« Hola Ana: \dots \textbf{Escríbeme} pronto, ¿vale? \textbf{Tienes} mi número. Un abrazo, Pedro »}

\textbf{Message 2 (formel) :} \esp{« Estimado Sr. Director: \dots Le ruego que \textbf{me escriba} a la mayor brevedad. \textbf{Tiene} usted mis datos adjuntos. Atentamente, Pedro Martínez »}
\end{exemplebox}

\begin{itemize}
    \item Dans le premier message, Pedro s'adresse à son amie Ana. Il utilise la \textbf{2e personne du singulier} : \esp{escríbeme} (impératif tú + pronom enclitique), \esp{tienes} (présent indicatif tú). Le ton est proche, familier.
    \item Dans le second, il s'adresse à un directeur qu'il ne connaît pas. Il utilise la \textbf{3e personne du singulier} : \esp{escriba} (subjonctif présent / impératif usted), \esp{tiene} (présent indicatif usted). Le pronom \esp{usted} apparaît explicitement dans \esp{tiene usted} pour lever toute ambiguïté, mais il est souvent omis. On note aussi la structure de politesse \esp{Le ruego que} + subjonctif.
\end{itemize}

\courssub{La règle fondamentale : tú vs. usted}

\begin{propriete}[Le vouvoiement espagnol : \textit{usted} + 3e personne]
En espagnol, le vouvoiement ne se fait \textbf{pas} avec la 2e personne du pluriel (\emph{vosotros}) comme en français ancien, ni avec la 2e personne du singulier. On utilise le pronom \cle{usted} (abrégé \emph{Ud.} ou \emph{Vd.}) qui \textbf{exige le verbe à la 3e personne du singulier} (mêmes formes que \emph{él/ella}).

\begin{center}
\begin{tabularx}{\linewidth}{|X|X|X|}
\hline
\rowcolor{popblueL} \textbf{Français} & \textbf{Espagnol (tuteo)} & \textbf{Espagnol (vouvoiement)} \\ \hline
Tu écris & \esp{Tú escribes} / \esp{Escribes} & \esp{Usted escribe} / \esp{Escribe} \\ \hline
Tu as & \esp{Tú tienes} / \esp{Tienes} & \esp{Usted tiene} / \esp{Tiene} \\ \hline
Tu peux & \esp{Tú puedes} / \esp{Puedes} & \esp{Usted puede} / \esp{Puede} \\ \hline
Écris ! (impératif affirmatif) & \esp{¡Escribe!} & \esp{¡Escriba!} \\ \hline
N'écris pas ! (impératif négatif) & \esp{¡No escribas!} & \esp{¡No escriba!} \\ \hline
\end{tabularx}
\end{center}

\textbf{Points clés :}
\begin{enumerate}
    \item \textbf{Omission du pronom :} En espagnol, les pronoms sujets (\esp{tú}, \esp{usted}) sont généralement \textbf{omis} car la terminaison verbale suffit (sauf ambiguïté à la 3e personne : \esp{escribe} = tú ou usted). Le contexte ou la forme de l'impératif lèvent le doute.
    \item \textbf{Accord des possessifs :} Le choix de \esp{tú} ou \esp{usted} impose les adjectifs possessifs correspondants :
    \begin{itemize}
        \item \esp{tu amigo} (ton ami) vs \esp{su amigo} (votre ami).
        \item \esp{tu carta} vs \esp{su carta}.
    \end{itemize}
    \item \textbf{Accord des pronoms objets :} \esp{te} (t' / à toi) vs \esp{le / la / se} (vous / à vous / le / la).
    \item \textbf{Cohérence absolue :} On \textbf{ne mélange jamais} les deux registres dans un même message. Pas de \esp{*Estimado Juan, le escribo para decirte...*} (mélange \emph{usted} / \emph{tú}).
\end{enumerate}
\end{propriete}

\begin{attention}[Piège fréquent : \textit{vosotros} n'est pas « vous » formel]
En Espagne, \esp{vosotros} est le « vous » \textbf{familier} (pluriel de \emph{tú}). En Amérique latine et en Guinée équatoriale, on utilise \emph{ustedes} (3e pers. pl.) pour tout « vous » (formel ou familier). Au Cameroun, l'espagnol enseigné suit souvent la norme péninsulaire pour la grammaire, mais la norme internationale pour la communication : \textbf{usted / ustedes} pour le formel. N'utilisez jamais \emph{vosotros} dans une lettre formelle.
\end{attention}

\courssub{L'impératif affirmatif dans la correspondance}

L'impératif affirmatif est le mode de la \textbf{demande directe}, de l'instruction ou de l'invitation. En correspondance, il est omniprésent : \esp{Responda}, \esp{Envíeme}, \esp{Dime}, \esp{Espera}.

\begin{propriete}[Formation de l'impératif affirmatif (réguliers)]
\begin{itemize}
    \item \textbf{Imperativo \textit{tú} :} = \textbf{3e personne du singulier du présent indicatif}. (Prononciation : accent sur l'avant-dernière syllabe).
    \begin{center}
    \begin{tabular}{|l|l|l|l|}
    \hline
    \rowcolor{popgreenL} \textbf{Infinitif} & \textbf{Présent 3e p. sg.} & \textbf{Imp. tú} & \textbf{Traduction} \\ \hline
    \esp{escribir} & \esp{escribe} & \esp{¡Escribe!} & Écris ! \\ \hline
    \esp{hablar} & \esp{habla} & \esp{¡Habla!} & Parle ! \\ \hline
    \esp{comer} & \esp{come} & \esp{¡Come!} & Mange ! \\ \hline
    \esp{esperar} & \esp{espera} & \esp{¡Espera!} & Attends ! \\ \hline
    \end{tabular}
    \end{center}
    \item \textbf{Imperativo \textit{usted} :} = \textbf{1e personne du singulier du présent de subjonctif} (racine du présent indicatif 1e sg + voyelle opposée : -ar $\to$ -e, -er/-ir $\to$ -a).
    \begin{center}
    \begin{tabular}{|l|l|l|l|}
    \hline
    \rowcolor{popblueL} \textbf{Infinitif} & \textbf{Présent 1e sg.} & \textbf{Subj. 1e sg.} & \textbf{Imp. usted} \\ \hline
    \esp{escribir} & \esp{escribo} & \esp{escriba} & \esp{¡Escriba!} \\ \hline
    \esp{hablar} & \esp{hablo} & \esp{hable} & \esp{¡Hable!} \\ \hline
    \esp{comer} & \esp{como} & \esp{coma} & \esp{¡Coma!} \\ \hline
    \esp{esperar} & \esp{espero} & \esp{espere} & \esp{¡Espere!} \\ \hline
    \end{tabular}
    \end{center}
\end{itemize}
\end{propriete}

\courssub{Les verbes irréguliers à l'impératif affirmatif (très fréquents en lettres)}

Certains verbes de haute fréquence ont un impératif \textbf{irrégulier pour \textit{tú}} (souvent monosyllabique) mais régulier pour \textit{usted} (subjonctif). Il est \textbf{impératif} de les connaître par cœur.

\begin{popfigure}
\legende{Tableau des impératifs affirmatifs irréguliers essentiels pour la correspondance}
\begin{center}
\begin{tikzpicture}[
    cell/.style={rectangle, draw=popline, minimum height=0.9cm, align=center, font=\small},
    header/.style={cell, fill=popblue, font=\bfseries, text=white},
    rowt/.style={cell, fill=popblueL},
    rowu/.style={cell, fill=poporangeL},
    rowf/.style={cell, fill=popgreenL}
]
% En-têtes
\node[header, minimum width=2.2cm] (v) at (0,0) {Verbo};
\node[header, minimum width=1.8cm] (tu) at (2.2,0) {Imperativo \textit{tú}};
\node[header, minimum width=1.8cm] (ud) at (4.0,0) {Imperativo \textit{usted}};
\node[header, minimum width=2.5cm] (tr) at (5.8,0) {Traduction (tú / Ud.)};

% Lignes
\foreach \i/\verb/\imt/\imu/\tr in {
    1/decir/di/diga/{Dis / Dites},
    2/hacer/haz/haga/{Fais / Faites},
    3/ir/ve/vaya/{Va / Allez},
    4/poner/pon/ponga/{Mets / Mettez},
    5/salir/sal/salga/{Sors / Sortez},
    6/tener/ten/tenga/{Aie / Ayez},
    7/venir/ven/venga/{Viens / Venez},
    8/ser/sé/sea/{Sois / Soyez}
} {
    \pgfmathsetmacro{\y}{-\i*0.95};
    \node[rowt] at (0,\y) {\verb};
    \node[rowu] at (2.2,\y) {\textbf{\esp{\imt}}};
    \node[rowf] at (4.0,\y) {\textbf{\esp{\imu}}};
    \node[cell] at (5.8,\y) {\tr};
}
\end{tikzpicture}
\end{center}
\end{popfigure}

\begin{exemplebox}[Exemples en contexte épistolaire]
\begin{itemize}
    \item \textbf{Amical (tú) :} \esp{« \textbf{Dime} la verdad. \textbf{Hazme} el favor de \textbf{venir} mañana. \textbf{Ten} paciencia. »} \\
    (Dis-moi la vérité. Fais-moi le plaisir de venir demain. Sois patient.)
    \item \textbf{Formel (usted) :} \esp{« \textbf{Dígame} la verdad, por favor. \textbf{Haga} el favor de \textbf{venir} mañana. \textbf{Tenga} paciencia. »} \\
    (Dites-moi la vérité, s'il vous plaît. Faites-moi le plaisir de venir demain. Ayez patience.)
    \item \textbf{Avec pronoms enclitiques (accrochés au verbe) :}
    \begin{itemize}
        \item \esp{Escríbeme} (tú) / \esp{Escríbame} (usted) — \emph{Écris-moi / Écrivez-moi}.
        \item \esp{Dímelo} (tú) / \esp{Dígamelo} (usted) — \emph{Dis-le-moi / Dites-le-moi}.
        \item \esp{Envíamelo} (tú) / \esp{Envíemelo} (usted) — \emph{Envoie-le-moi / Envoyez-le-moi}.
    \end{itemize}
    \textbf{Rappel accentuation :} L'ajout du ou des pronoms crée une syllabe supplémentaire $\to$ accent écrit sur la voyelle tonique d'origine (\esp{escríbe} $\to$ \esp{escríbeme}, \esp{dí} $\to$ \esp{dímelo}, \esp{ven} $\to$ \esp{véngame}, \esp{pon} $\to$ \esp{póngamelo}).
\end{itemize}
\end{exemplebox}

\courssub{L'impératif négatif : interdire ou déconseiller}

Tout aussi fréquent dans les courriers formels (consignes, mises en garde) ou amicaux (conseils), l'impératif négatif se forme \textbf{toujours avec le subjonctif présent} pour les deux registres.

\begin{propriete}[Formation de l'impératif négatif]
\begin{itemize}
    \item \textbf{Structure :} \esp{No} + \textbf{subjonctif présent} (2e p. sg pour \textit{tú}, 3e p. sg pour \textit{usted}).
    \item \textbf{Pronoms :} Ils se placent \textbf{avant} le verbe (proclitiques) : \esp{No me escribas}, \esp{No me escriba}.
\end{itemize}

\begin{center}
\begin{tabular}{|l|l|l|l|l|}
\hline
\rowcolor{poppurpleL} \textbf{Infinitif} & \textbf{No + tú (subj)} & \textbf{No + usted (subj)} & \textbf{Traduction} \\ \hline
\esp{escribir} & \esp{No escribas} & \esp{No escriba} & N'écris pas / N'écrivez pas \\ \hline
\esp{hablar} & \esp{No hables} & \esp{No hable} & Ne parle pas / Ne parlez pas \\ \hline
\esp{comer} & \esp{No comas} & \esp{No coma} & Ne mange pas / Ne mangez pas \\ \hline
\esp{ir} & \esp{No vayas} & \esp{No vaya} & Ne va pas / N'allez pas \\ \hline
\esp{ser} & \esp{No seas} & \esp{No sea} & Ne sois pas / Ne soyez pas \\ \hline
\esp{preocuparse} & \esp{No te preocupes} & \esp{No se preocupe} & Ne t'inquiète pas / Ne vous inquiétez pas \\ \hline
\end{tabular}
\end{center}
\end{propriete}

\begin{exemplebox}[Impératif négatif en contexte]
\begin{itemize}
    \item \textbf{Amical :} \esp{« \textbf{No te preocupes} por el examen, \textbf{no olvides} llamarme. »} (Ne t'inquiète pas pour l'examen, n'oublie pas de m'appeler.)
    \item \textbf{Formel :} \esp{« \textbf{No se preocupe} por las molestias. \textbf{No olvide} adjuntar el documento. »} (Ne vous inquiétez pas pour le dérangement. N'oubliez pas de joindre le document.)
    \item \textbf{Avertissement / Consigne :} \esp{« \textbf{No responda} a este correo automático. \textbf{No envíe} datos personales. »}
\end{itemize}
\end{exemplebox}

\begin{aretenir}[Synthèse : Choisir la bonne forme pour écrire]
\begin{center}
\begin{tabularx}{\linewidth}{|X|X|X|}
\hline
\rowcolor{popgoldL} \textbf{Situation} & \textbf{Registre} & \textbf{Formes clés (Écrire / Ne pas écrire)} \\ \hline
Ami, famille, camarade & \textbf{Tuteo (tú)} & \esp{Escribe / No escribas} \newline \esp{Dime / No me digas} \newline \esp{Tu carta, tu amigo} \\ \hline
Supérieur, administration, inconnu, client & \textbf{Vouvoiement (usted)} & \esp{Escriba / No escriba} \newline \esp{Dígame / No me diga} \newline \esp{Su carta, su amigo} \\ \hline
\end{tabularx}
\end{center}
\textbf{Règle d'or :} Une fois le registre choisi (tú \textbf{ou} usted), on maintient \textbf{tous} les marqueurs (verbes, possessifs, pronoms objets, formules de politesse) cohérents jusqu'à la signature.
\end{aretenir}

\coursec{Structure de la lettre et du courriel}

Après avoir distingué les registres \cle{tú} et \cle{usted} et révisé l'impératif de politesse (\esp{por favor, tenga la bondad de…}), nous abordons maintenant l'architecture même du message écrit. En espagnol, comme en français, la lettre et le courriel obéissent à des codes stricts : la \cle{structure} garantit la clarté et marque le respect dû au destinataire. Que vous écriviez à un ami à Madrid ou à une université à Buenos Aires, le squelette reste identique ; seul le \cle{registre} (vocabulaire, formules, grammaire) change.

\courssub{Les cinq parties de la lettre formelle et amicale}

Toute lettre hispanophone — qu'elle soit manuscrite, dactylographiée ou envoyée en pièce jointe — se compose de \textbf{cinq zones obligatoires}, toujours dans cet ordre :

\begin{propriete}[Les 5 parties canoniques de la carta]
\begin{enumerate}
    \item \textbf{Lugar y fecha} (Lieu et date) : en haut à droite. Ex. \esp{Yaoundé, a 15 de marzo de 2025}.
    \item \textbf{Saludo} (Formule d'appel) : à gauche, suivi de deux-points \esp{:} (jamais de virgule). Ex. \esp{Querido Juan:} / \esp{Estimado Sr. Director:}.
    \item \textbf{Cuerpo} (Corps du message) : organisé en trois paragraphes logiques :
    \begin{itemize}
        \item \emph{Introducción} : motif de la lettre.
        \item \emph{Desarrollo} : détails, arguments, narration.
        \item \emph{Conclusión} : appel à l'action, vœu de réponse.
    \end{itemize}
    \item \textbf{Despedida} (Formule de clôture) : à gauche, avant la signature. Ex. \esp{Un abrazo,} / \esp{Atentamente,}.
    \item \textbf{Firma} (Signature) : prénom seul (amical) ou Prénom Nom (formel).
\end{enumerate}
\end{propriete}

\begin{popfigure}
\begin{tikzpicture}[
    box/.style={draw=popblue, thick, fill=popblueL, rounded corners=4pt, minimum height=1.8cm, text width=5.5cm, align=center, font=\small},
    arrow/.style={->, >=Stealth, thick, popdark},
    label/.style={font=\tiny\bfseries, text=popblue, anchor=west}
]
% Boxes
\node[box, align=center] (fecha){\textbf{1. LUGAR Y FECHA}\\\esp{Yaoundé, a 15 de marzo de 2025}};
\node[box, below=1.2cm of fecha, align=center] (saludo){\textbf{2. SALUDO}\\\esp{Querido Pablo:}\\(ou \esp{Estimado Sr. Gómez:})};
\node[box, below=1.2cm of saludo, minimum height=3.5cm, align=center] (cuerpo){\textbf{3. CUERPO}\\\emph{Intro:} \esp{Te escribo para contarte…}\\\emph{Desarrollo:} \esp{El viaje fue increíble…}\\\emph{Conclusión:} \esp{Espero tus noticias pronto.}};
\node[box, below=1.2cm of cuerpo, align=center] (despedida){\textbf{4. DESPEDIDA}\\\esp{Un fuerte abrazo,}};
\node[box, below=1.2cm of despedida, align=center] (firma){\textbf{5. FIRMA}\\\esp{Tu amigo,}\\\esp{Samuel}};

% Arrows
\draw[arrow] (fecha.south) -- (saludo.north);
\draw[arrow] (saludo.south) -- (cuerpo.north);
\draw[arrow] (cuerpo.south) -- (despedida.north);
\draw[arrow] (despedida.south) -- (firma.north);

% Side labels
\node[label] at ($(fecha.east)+(0.3,0)$) {À droite};
\node[label] at ($(saludo.east)+(0.3,0)$) {À gauche + \esp{:}};
\node[label] at ($(cuerpo.east)+(0.3,0)$) {3 paragraphes};
\node[label] at ($(despedida.east)+(0.3,0)$) {À gauche + \esp{,}};
\node[label] at ($(firma.east)+(0.3,0)$) {Sous la despedida};
\end{tikzpicture}
\legende{Maquette de la lettre hispanophone : les 5 zones obligatoires et leur emplacement}
\end{popfigure}

\begin{exemplebox}[Lettre amicale (registro informal)]
\esp{
Yaoundé, a 20 de mayo de 2025

Querido Carlos:

Te escribo porque quiero contarte mis vacaciones. El mes pasado viajé a la costa con mi familia. Visitamos Limbe y Kribi; las playas son hermosas y comimos pescado fresco todos los días. Además, jugué al fútbol con unos amigos locales en la arena.

El tiempo pasó muy rápido. Mañana vuelvo a las clases en el liceo. Espero que tú también lo pases bien este verano. Escríbeme cuando puedas.

Un fuerte abrazo,

Samuel
}
\textbf{Traduction :} Yaoundé, le 20 mai 2025. Cher Carlos, Je t'écris parce que je veux te raconter mes vacances. Le mois dernier, j'ai voyagé sur la côte avec ma famille. Nous avons visité Limbe et Kribi ; les plages sont magnifiques et nous avons mangé du poisson frais tous les jours. De plus, j'ai joué au foot avec des amis locaux sur le sable. Le temps a passé très vite. Demain, je reprends les cours au lycée. J'espère que toi aussi tu passes un bel été. Écris-moi quand tu peux. Gros bisous, Samuel.
\end{exemplebox}

\begin{exemplebox}[Lettre formelle (registro formal)]
\esp{
Douala, a 10 de octubre de 2025

Estimado Sr. Director:

Me dirijo a usted para solicitar información sobre el proceso de admisión en su universidad para el curso 2026. Soy estudiante de Terminale en el Liceo Bilingüe de Douala y deseo cursar Estudios Internacionales.

Me gustaría conocer los requisitos de acceso, las fechas límite de inscripción y la posibilidad de obtener una beca para estudiantes extranjeros. Adjunto mi expediente académico y una carta de motivación.

Quedo a la espera de su amable respuesta y le saluda atentamente,

Samuel Mbarga
}
\textbf{Traduction :} Douala, le 10 octobre 2025. Monsieur le Directeur, Je m'adresse à vous pour solliciter des informations sur le processus d'admission dans votre université pour l'année 2026. Je suis élève de Terminale au Lycée Bilingue de Douala et je souhaite suivre des Études Internationales. Je voudrais connaître les conditions d'accès, les dates limites d'inscription et la possibilité d'obtenir une bourse pour étudiants étrangers. Je joins mon dossier scolaire et une lettre de motivation. Dans l'attente de votre aimable réponse, je vous adresse mes salutations distinguées. Samuel Mbarga.
\end{exemplebox}

\begin{attention}[Mois et jours : pas de majuscule !]
En espagnol, les noms des mois (\esp{enero, febrero, marzo…}) et des jours de la semaine (\esp{lunes, martes, miércoles…}) s'écrivent \textbf{toujours en minuscules}, même dans la date en tête de lettre.
\begin{center}
\esp{Yaoundé, a 15 de marzo de 2025} \quad (correct) \\
\esp{Yaoundé, a 15 de Marzo de 2025} \quad (incorrect)
\end{center}
\emph{Note pour francophones :} en français aussi, les mois et jours sont en minuscules (\emph{mars, lundi}), mais l'erreur fréquente vient de l'influence de l'anglais (\emph{March, Monday}), très présent dans l'environnement scolaire camerounais. Soyez vigilants !
\end{attention}

\begin{methode}[Rédiger une lettre en 5 étapes]
\begin{enumerate}
    \item \textbf{Lieu et date} : en haut à droite. \esp{Yaoundé, a [jour] de [mois] de [année]}. La préposition \esp{a} est obligatoire devant le jour.
    \item \textbf{Salutation} : à gauche. \esp{Querido/a [Prénom]:} (ami) \textbf{ou} \esp{Estimado/a Sr./Sra. [Nom]:} (formel). \textbf{Toujours deux-points \esp{:}}.
    \item \textbf{Corps en 3 paragraphes} : sauter une ligne entre chaque.
    \begin{itemize}
        \item \emph{Intro} : \esp{Te escribo para… / Me dirijo a usted para…}
        \item \emph{Développement} : connecteurs \esp{primero, además, por eso, pero…}
        \item \emph{Conclusion} : \esp{Espero tu respuesta. / Quedo a la espera de su respuesta.}
    \end{itemize}
    \item \textbf{Formule de clôture} : à gauche, suivie d'une virgule \esp{,}.
    \begin{itemize}
        \item Amical : \esp{Un abrazo,} / \esp{Un fuerte abrazo,} / \esp{Besos,} / \esp{Hasta pronto,}
        \item Formel : \esp{Atentamente,} / \esp{Le saluda atentamente,} / \esp{Cordialmente,}
    \end{itemize}
    \item \textbf{Signature} : prénom (ami) ou Prénom Nom (formel). Pas de point final.
\end{enumerate}
\end{methode}

\courssub{La date en espagnol : règles et usage}

La formule \esp{lugar, a día de mes de año} est immuable. La préposition \esp{a} (sans accent) introduit le jour.

\begin{center}
\begin{tabularx}{\linewidth}{|X|X|X|}
\hline
\rowcolor{popblueL} \textbf{Français} & \textbf{Espagnol} & \textbf{Remarque} \\ \hline
Yaoundé, le 15 mars 2025 & \esp{Yaoundé, a 15 de marzo de 2025} & \esp{a} + jour + \esp{de} + mois + \esp{de} + année \\ \hline
Le 1er janvier & \esp{a 1 de enero} / \esp{a primero de enero} & \esp{primero} possible seulement pour le 1er \\ \hline
Lundi 15 mars & \esp{lunes 15 de marzo} & jour en minuscule, pas de virgule \\ \hline
\emph{Date courte (chiffres)} & \esp{15/03/2025} ou \esp{15-03-2025} & Jour/Mois/Année (comme en France) \\ \hline
\end{tabularx}
\end{center}

\begin{aretenir}[La préposition \esp{a} dans la date]
Dans l'en-tête d'une lettre, on écrit \esp{a 15 de marzo} (pas \esp{en}, pas \esp{el}). Attention à la différence de fonction : dans une phrase, quand la date est sujet ou complément, on utilise l'article : \esp{El 15 de marzo es mi cumpleaños.} « Le 15 mars, c'est mon anniversaire. » Mais \textbf{jamais} d'article dans l'en-tête de la lettre.
\end{aretenir}

\courssub{Structure du courriel (correo electrónico)}

Le courriel reprend la logique de la lettre, mais avec des champs techniques imposés par l'interface de messagerie. Le registre (tú/usted) dicte le ton du \esp{saludo}, du \esp{cuerpo} et de la \esp{despedida}.

\begin{propriete}[Anatomie del correo electrónico]
\begin{enumerate}
    \item \textbf{Para :} adresse(s) du/des destinataire(s).
    \item \textbf{CC / CCO :} copie (con copia) / copie cachée (con copia oculta).
    \item \textbf{Asunto :} \cle{objet}. Court, précis, informatif. Jamais vide !
    \item \textbf{Saludo :} identique à la lettre (\esp{Querido…:} / \esp{Estimado…:}).
    \item \textbf{Cuerpo :} même structure en 3 paragraphes. Phrases complètes, pas de style « SMS ».
    \item \textbf{Despedida :} identique à la lettre.
    \item \textbf{Firma :} nom + coordonnées (téléphone, établissement) en registre formel.
\end{enumerate}
\end{propriete}

\begin{popfigure}
\begin{tikzpicture}[
    mailbox/.style={draw=poppurple, thick, fill=poppurpleL, rounded corners=4pt, minimum height=1.1cm, text width=8cm, align=left, font=\small, inner xsep=6pt, inner ysep=5pt}
]
\node[mailbox] (para) {\textbf{Para:} \esp{director@universidad.edu}};
\node[mailbox, below=0.3cm of para] (asunto) {\textbf{Asunto:} \esp{Solicitud de información — Admisión 2026}};
\node[mailbox, below=0.3cm of asunto, minimum height=4.2cm, align=center] (cuerpo){\textbf{Saludo:} \esp{Estimado Sr. Director:}\\[2pt]
\textbf{Cuerpo:}\\
\quad \esp{Me dirijo a usted para solicitar…}\\
\quad \esp{Adjunto mi expediente…}\\
\quad \esp{Quedo a la espera de su respuesta…}\\[2pt]
\textbf{Despedida:} \esp{Atentamente,}\\[2pt]
\textbf{Firma:} \esp{Samuel Mbarga}\\
\quad \esp{Tel: +237 6XX XX XX XX}};
\node[above=3pt of para.north, font=\bfseries\small, text=poppurple] {INTERFACE MESSAGERIE};
\end{tikzpicture}
\legende{Maquette d'un courriel formel : champs techniques (Para, Asunto) et zones rédactionnelles}
\end{popfigure}

\begin{propriete}[L'objet (Asunto) : règles d'or]
\begin{itemize}
    \item \textbf{Court :} 5 à 8 mots maximum.
    \item \textbf{Précis :} mots-clés : \esp{Solicitud, Información, Invitación, Reclamación, Confirmación, Cancelación}.
    \item \textbf{Pas de phrase complète :} pas de verbe conjugué, pas de point final.
    \item \textbf{Exemples :}
    \begin{center}
    \begin{tabular}{|l|l|}
    \hline
    \rowcolor{poporangeL} \textbf{\esp{Asunto}} & \textbf{Contexte} \\ \hline
    \esp{Solicitud de beca MAEC-AECID} & demande de bourse \\ \hline
    \esp{Invitación: conferencia 15 de mayo} & invitation à un événement \\ \hline
    \esp{Reclamación: pedido n.\textsuperscript{o} 452} & réclamation sur une commande \\ \hline
    \esp{Confirmación de matrícula} & confirmation administrative \\ \hline
    \end{tabular}
    \end{center}
\end{itemize}
\end{propriete}

\courssub{Connecteurs et formules utiles pour le corps de la lettre}

Le \esp{cuerpo} doit être cohérent. Voici les connecteurs (\esp{conectores}) essentiels du niveau A2+/B1, classés par fonction.

\begin{center}
\begin{tabularx}{\linewidth}{|X|X|X|}
\hline
\rowcolor{popgreenL} \textbf{Fonction} & \textbf{Conectores (Español)} & \textbf{Français} \\ \hline
Ordonner / énumérer & \esp{primero, en primer lugar, segundo, luego, finalmente} & premièrement, en premier lieu, deuxièmement, ensuite, enfin \\ \hline
Ajouter & \esp{además, también, igualmente, asimismo} & de plus, aussi, de même, par ailleurs \\ \hline
Contraster & \esp{pero, sin embargo, no obstante, aunque} & mais, cependant, néanmoins, bien que \\ \hline
Cause / conséquence & \esp{porque, por eso, por tanto, por consiguiente} & parce que, c'est pourquoi, donc, par conséquent \\ \hline
But & \esp{para + inf., con el fin de + inf., para que + subj.} & pour, afin de, pour que \\ \hline
Exprimer l'attente & \esp{espero que + subj., quedo a la espera de} & j'espère que + subj., je reste dans l'attente de \\ \hline
\end{tabularx}
\end{center}

Les formules de demande varient radicalement selon le registre.

\begin{definition}[Formules de demande (pedir / solicitar)]
\begin{itemize}
    \item \textbf{Amical (tú) :} \esp{¿Me puedes…? / ¿Me das…? / Quiero que… / Avísame…}
    \item \textbf{Standard poli (usted) :} \esp{Quisiera… / Me gustaría… / Le agradecería que… + subjonctif}
    \item \textbf{Très formel (administratif) :} \esp{Solicito… / Tengo el honor de solicitar… / Ruego a usted que… + subjonctif}
\end{itemize}
\end{definition}

\begin{exemplebox}[Exemples traduits : du souhait à la demande formelle]
\begin{itemize}
    \item \esp{Quisiera solicitar información sobre los cursos.} \\
    « Je souhaiterais demander des informations sur les cours. » (Politesse standard : \cle{quisiera} est l'imparfait du subjonctif de \esp{querer}, employé comme conditionnel de politesse.)
    \item \esp{Me gustaría que me enviaran el programa.} \\
    « J'aimerais qu'ils m'envoient le programme. » (\esp{me gustaría} + \cle{subjonctif imparfait} \esp{enviaran}.)
    \item \esp{Le agradecería que confirmara la cita.} \\
    « Je vous serais reconnaissant de bien vouloir confirmer le rendez-vous. » (\esp{agradecer} au conditionnel + subjonctif imparfait \esp{confirmara}.)
    \item \esp{Espero tu respuesta.} / \esp{Espero su respuesta.} \\
    « J'attends ta réponse. » / « J'attends votre réponse. » (Nom après \esp{esperar} : pas de subjonctif.)
    \item \esp{Quedo a la espera de su respuesta y le saluda atentamente.} \\
    « Dans l'attente de votre réponse, je vous adresse mes salutations distinguées. » (Formule figée de clôture formelle.)
\end{itemize}
\end{exemplebox}

\begin{exoresolu}[Transformer une demande amicale en demande formelle]
\textbf{Énoncé :} Réécrivez le paragraphe suivant en registre formel (\cle{usted}), en adaptant les formules de politesse, les connecteurs et la structure.

\esp{
Querido Juan:

Te escribo porque quiero pedirte un favor. ¿Me puedes mandar los apuntes de historia? Además, necesito que me expliques el tema de la Revolución. Por eso, avísame cuando puedas.

Un abrazo,

Pedro
}

\tcblower
\textbf{Solution détaillée :}

\esp{
Estimado Sr. Juan:

Me dirijo a usted para pedirle un favor. \textbf{Quisiera} que me \textbf{enviara} los apuntes de historia. \textbf{Asimismo}, \textbf{me gustaría que me explicara} el tema de la Revolución. \textbf{Por consiguiente}, \textbf{le agradecería que me avisara} cuando le sea posible.

Atentamente,

Pedro
}

\textbf{Analyse des transformations :}
\begin{enumerate}
    \item \esp{Querido Juan:} $\rightarrow$ \esp{Estimado Sr. Juan:} (salutation formelle).
    \item \esp{Te escribo porque quiero pedirte} $\rightarrow$ \esp{Me dirijo a usted para pedirle} (style formel ; \esp{le} = COI pour \esp{usted}).
    \item \esp{¿Me puedes mandar?} $\rightarrow$ \esp{Quisiera que me enviara} (\cle{quisiera} + subjonctif imparfait \esp{enviara}).
    \item \esp{Además} $\rightarrow$ \esp{Asimismo} (connecteur plus soutenu).
    \item \esp{Necesito que me expliques} $\rightarrow$ \esp{Me gustaría que me explicara} (conditionnel + subjonctif imparfait \esp{explicara}).
    \item \esp{Por eso} $\rightarrow$ \esp{Por consiguiente} (conséquence formelle).
    \item \esp{Avísame} (impératif \esp{tú}) $\rightarrow$ \esp{Le agradecería que me avisara} (conditionnel + subjonctif).
    \item \esp{Un abrazo} $\rightarrow$ \esp{Atentamente} (clôture formelle).
\end{enumerate}
\end{exoresolu}

\begin{attention}[Pièges fréquents des francophones]
\begin{itemize}
    \item \textbf{Deux-points, pas virgule :} \esp{Estimado Sr. Gómez:} (deux-points obligatoires après la salutation). En français : « Monsieur Gómez, » (virgule). \textbf{Ne pas écrire} \esp{Estimado Sr. Gómez,}.
    \item \textbf{Préposition \esp{a} dans la date :} \esp{Yaoundé, a 15 de marzo}. Pas \esp{en}, pas \esp{el}.
    \item \textbf{Objet vide ou phrase complète :} \esp{Asunto: Hola} (mauvais). \esp{Asunto: Solicitud de información} (bon).
    \item \textbf{Subjonctif après \esp{quisiera que / me gustaría que / le agradecería que} :} \esp{Quisiera que me \textbf{enviara} los apuntes} (pas \esp{envía}).
    \item \textbf{Faux ami :} \esp{asunto} signifie « objet / affaire », pas « assunto » au sens de « sujet d'examen » ; pour « sujet » d'un devoir, on dit \esp{tema}.
    \item \textbf{Signature :} pas de point après le nom : \esp{Samuel Mbarga} (pas \esp{Samuel Mbarga.}).
\end{itemize}
\end{attention}

\courssub{Complément Bac : la carte postale et le message d'invitation courts (30--50 mots)}

L'épreuve d'expression écrite du Baccalauréat camerounais (LV2) propose souvent un sujet court : carte postale (\esp{tarjeta postal}), message d'invitation (\esp{nota de invitación}) ou petit message. La structure est compressée : \textbf{Saludo + 2 ou 3 phrases (motif + information clé) + Despedida + Firma}. La date peut être abrégée (\esp{Madrid, 20 de julio}).

\begin{textebox}[Modèle de lettre formelle à trous — Activité de classe]
Complétez la lettre ci-dessous avec les éléments fournis (formules, connecteurs, verbes au bon mode). Travail en binôme, puis correction collective au tableau.

\esp{
Douala, a \underline{\hspace{3cm}} de 2025

\underline{\hspace{4cm}} Sr. Director:

Me dirijo a usted \underline{\hspace{2cm}} (but) solicitar información sobre los másteres en Derecho Internacional. Soy estudiante de último año en la Universidad de Yaoundé II. \underline{\hspace{2cm}} (connecteur d'addition), \underline{\hspace{2cm}} (verbe gustar au conditionnel) conocer los requisitos de admisión y las fechas de matrícula.

\underline{\hspace{3cm}} (formule d'attente) su respuesta y le saluda \underline{\hspace{2cm}} (despedida formelle),

\underline{\hspace{3cm}} (Firma)
}

\textbf{Mots à placer :} \esp{20 de octubre / Estimado / con el fin de / Además / Me gustaría / Quedo a la espera de / atentamente / [votre nom]}.

\textbf{Glossaire :} \esp{matrícula} (inscription), \esp{requisitos} (conditions), \esp{último año} (dernière année).
\end{textebox}

\begin{savaistu}[Au Bac, la lettre amicale de 80 mots : la structure, ce sont des points faciles]
L'épreuve d'\textbf{expression écrite} (coefficient fort en Première A comme au Bac) contient très régulièrement un sujet de type « \emph{Carta a un amigo} » (environ 80 mots). Les correcteurs notent sur une grille précise :
\begin{itemize}
    \item \textbf{Respect de la structure :} lieu et date à droite, salutation avec \esp{:}, paragraphes visibles, \esp{despedida} + virgule, signature.
    \item \textbf{Registre :} \cle{tú} constant du début à la fin, vocabulaire affectif (\esp{querido, un abrazo, escríbeme}), aucune intrusion de \esp{usted}.
    \item \textbf{Contenu et cohérence :} réponse exacte au sujet posé, connecteurs (\esp{pero, también, por eso}).
    \item \textbf{Correction linguistique :} accents, \cle{ser/estar}, passé composé et imparfait, subjonctif après \esp{espero que} et \esp{quiero que}.
\end{itemize}
\textbf{Astuce :} avant de compter les mots, vérifiez que les \textbf{5 zones} sont présentes. Une lettre sans date ou sans \esp{despedida} perd des points de structure, même si l'espagnol est parfait !
\end{savaistu}

\begin{exercice}[Entraînement : du courriel à la carte postale]
\begin{enumerate}
    \item Rédigez un \textbf{courriel formel} (objet + corps complet) à l'adresse \esp{becas@educacion.gob.es} pour demander le dossier de candidature aux bourses MAEC-AECID (80--100 mots).
    \item Transformez-le en \textbf{carte postale amicale} (environ 35 mots) envoyée à votre correspondant \esp{Pablo} à Madrid pour lui annoncer que vous avez obtenu la bourse.
    \item Soulignez dans les deux textes : \textcolor{popblue}{la date}, \textcolor{poporange}{la salutation}, \textcolor{popgreen}{les connecteurs}, \textcolor{poppink}{la despedida}.
\end{enumerate}
\end{exercice}

\begin{corrige}[Exercice n°1 -- Proposition de correction]
\textbf{1. Courriel formel :}

\esp{
Para: becas@educacion.gob.es

Asunto: Solicitud de expediente — Becas MAEC-AECID 2026

Estimados señores:

Me dirijo a ustedes para solicitar el expediente de candidatura para las becas MAEC-AECID del curso 2026. Soy estudiante de Terminale en el Liceo Bilingüe de Yaoundé y deseo cursar estudios universitarios en España.

Me gustaría recibir información sobre los plazos, la documentación requerida y los criterios de selección. Adjunto mi certificado de notas y una carta de motivación.

Quedo a la espera de su amable respuesta y les saluda atentamente,

Samuel Mbarga

Tel: +237 6XX XX XX XX
}

\textbf{2. Carte postale amicale :}

\esp{
Madrid, 20 de julio de 2025

Querido Pablo:

¡Notición! Me han concedido la beca MAEC para estudiar en Madrid el próximo curso. Además, voy a vivir en el barrio de Lavapiés. Por eso, nos vemos en septiembre para celebrarlo.

Un abrazo fuerte,

Samuel
}

\textbf{Remarques de correction :}
\begin{itemize}
    \item \esp{¡Notición!} : exclamation familière (« super nouvelle ! »), typique du registre amical ; notez le signe \esp{¡} d'ouverture obligatoire.
    \item \esp{Me han concedido la beca} : passé composé (action récente avec conséquence présente), correct en Espagne.
    \item \esp{Estimados señores:} et \esp{les saluda atentamente} : le pluriel est maintenu du début à la fin (cohérence du registre).
    \item Dans le courriel, l'objet (\esp{Asunto}) ne contient ni verbe conjugué ni point final.
\end{itemize}
\end{corrige}

\coursec{Production guidée : escribo a mi correspondiente}

Après avoir étudié la structure de la lettre et du courriel dans la section précédente, nous passons maintenant à l'acte d'écriture lui-même. Savoir identifier les parties d'une lettre ne suffit pas ; il faut être capable de les articuler avec cohérence, de respecter un registre donné et de maîtriser la contrainte de longueur, exercice classique de l'examen de fin d'année et du Baccalauréat. Cette section vous guide pas à pas, de l'analyse du sujet à l'auto-évaluation, en passant par la rédaction d'un modèle et sa transformation en courriel formel.

\courssub{Analyse de la consigne type examen}

\begin{textebox}[Sujet d'examen type]
\textbf{Consigne :} \esp{Escribe una carta a tu amigo hispanohablante para invitarlo a pasar las vacaciones en Camerún (80 palabras aproximadamente).}

\emph{Traduction :} Écris une lettre à ton ami hispanophone pour l'inviter à passer les vacances au Cameroun (environ 80 mots).
\end{textebox}

\begin{methode}[Analyser un sujet de correspondance en 3 étapes]
\begin{enumerate}
    \item \textbf{Identifier le destinataire et le registre :} \esp{tu amigo hispanohablante} $\rightarrow$ registre \textbf{familier / amical} (tutoiement \cle{tú}, formules affectueuses).
    \item \textbf{Identifier la finalité communicative :} \esp{invitarlo a pasar las vacaciones en Camerún} $\rightarrow$ fonction \textbf{invitative} + \textbf{descriptive} (il faut « vendre » la destination).
    \item \textbf{Identifier les contraintes formelles :} \esp{80 palabras aproximadamente} $\rightarrow$ ni trop court (pénalité), ni trop long (risque d'hors-sujet). Tolérance habituelle : $\pm 10\%$, soit entre 72 et 88 mots.
\end{enumerate}
\end{methode}

\begin{attention}[Piège fréquent : le comptage des mots]
En espagnol, on compte les \textbf{mots graphiques} (séparés par des espaces). Les contractions \esp{del} (\esp{de + el}) et \esp{al} (\esp{a + el}) comptent pour \textbf{un seul mot}. Les formules de politesse (\esp{Querido Pedro:}, \esp{Un abrazo,}) comptent aussi. N'oubliez pas de les inclure dans votre total.
\end{attention}

\courssub{Brainstorming collectif et plan imposé}

Avant d'écrire, il faut générer des idées. Voici le fruit d'un brainstorming typique en classe, organisé selon le \textbf{plan imposé en 4 blocs} (structure standardisée pour garantir la cohérence et faciliter l'évaluation).

\begin{popfigure}
\centering
\begin{tikzpicture}[
    node distance=0.9cm and 0.6cm,
    block/.style={rectangle, draw=popblue, thick, fill=popblueL, rounded corners, minimum width=3.2cm, minimum height=2.2cm, align=center, font=\small},
    arrow/.style={-Stealth, thick, popdark},
    title/.style={font=\bfseries\small, text=popdark}
]
% Nodes
\node[block, align=center] (saludo){\textbf{1. SALUDO}\\\emph{(1 ligne)}\\\esp{Querido Pedro:}};
\node[block, right=of saludo, align=center] (noticias){\textbf{2. NOTICIAS}\\\emph{(2 lignes)}\\\esp{¿Qué tal? Yo estoy bien...}\\nouvelles, climat, études};
\node[block, right=of noticias, align=center] (invitacion){\textbf{3. INVITACIÓN}\\\textbf{+ PROGRAMA}\\\emph{(4 lignes)}\\\esp{Te invito a Camerún...}\\plages, faune, culture, dates};
\node[block, right=of invitacion, align=center] (despedida){\textbf{4. DESPEDIDA}\\\emph{(1 ligne)}\\\esp{Espero tu respuesta.}\\\esp{Un abrazo,}\\\esp{Tu amigo}};

% Arrows
\draw[arrow] (saludo.east) -- (noticias.west);
\draw[arrow] (noticias.east) -- (invitacion.west);
\draw[arrow] (invitacion.east) -- (despedida.west);

% Annotations
\node[title, above=0.15cm of saludo.north] {Bloc A};
\node[title, above=0.15cm of noticias.north] {Bloc B};
\node[title, above=0.15cm of invitacion.north] {Bloc C};
\node[title, above=0.15cm of despedida.north] {Bloc D};
\end{tikzpicture}
\legende{Organigramme du plan de la lettre amicale en 4 blocs (cible : 80 mots)}
\end{popfigure}

\begin{center}
\begin{tabularx}{\linewidth}{|X|X|X|}
\hline
\rowcolor{popblueL} \textbf{Bloc} & \textbf{Idées / Contenu} & \textbf{Connecteurs utiles} \\
\hline
\textbf{1. Saludo} & \esp{Querido/a} + prénom (ou \esp{Hola} + prénom) & -- \\
\hline
\textbf{2. Noticias} & Nouvelles personnelles (examens, famille, météo à Yaoundé ou Douala), question au destinataire. & \esp{Por cierto,} \esp{A propósito,} \esp{Además,} \\
\hline
\textbf{3. Invitación y programa} & \textbf{Cœur de la lettre.} Invitation explicite + 2 ou 3 arguments (plages de Kribi ou Limbé, parc de Waza, gastronomie : \esp{ndolé}, \esp{suya}, ambiance \esp{bendskin}) + dates proposées. & \esp{Por eso,} \esp{Por tanto,} \esp{Además,} \esp{También,} \esp{En concreto,} \\
\hline
\textbf{4. Despedida} & Appel à réponse + formule affective + signature (prénom seulement). & \esp{Espero tu carta,} \esp{Avísame pronto,} \\
\hline
\end{tabularx}
\end{center}

\courssub{Banque de phrases toutes faites (répertoire lexical)}

Cette boîte est votre « trousse à outils ». Mémorisez deux ou trois phrases par bloc pour gagner du temps le jour de l'examen.

\begin{textebox}[Banque de phrases — Lettre amicale (invitation)]
\textbf{Saludo / Apertura}

\esp{Querido Carlos:} / \esp{Querida Ana:} (Cher Carlos / Chère Ana)

\esp{Hola, Miguel:} (Salut Miguel)

\esp{¿Cómo estás? Espero que todo vaya bien.} (Comment vas-tu ? J'espère que tout va bien.)

\textbf{Noticias (dar noticias / preguntar)}

\esp{Hace mucho tiempo que no sé nada de ti.} (Cela fait longtemps que je n'ai pas de tes nouvelles.)

\esp{Por aquí todo bien, aunque con muchos exámenes.} (Ici tout va bien, mais avec beaucoup d'examens.)

\esp{El clima en Yaoundé es caluroso y húmedo.} (Le climat à Yaoundé est chaud et humide.)

\esp{¿Qué tal tus estudios / tu familia / el verano?} (Comment vont tes études / ta famille / l'été ?)

\textbf{Invitación y programa (función invitativa + descriptiva)}

\esp{Te escribo para invitarte a pasar las vacaciones en Camerún.} (Je t'écris pour t'inviter à passer les vacances au Cameroun.)

\esp{Me encantaría que vinieras en julio o agosto.} (J'aimerais beaucoup que tu viennes en juillet ou en août.) \textcolor{poppurple}{\emph{[Subjonctif imparfait !]}}

\esp{Podríamos visitar las playas de Kribi y el parque nacional de Waza.} (Nous pourrions visiter les plages de Kribi et le parc national de Waza.)

\esp{También probarás la comida típica: ndolé, poulet DG, suya.} (Tu goûteras aussi la cuisine typique.)

\esp{La gente es muy acogedora y hay buen ambiente.} (Les gens sont très accueillants et l'ambiance est bonne.)

\textbf{Despedida / Cierre}

\esp{Espero tu respuesta pronto.} (J'attends ta réponse bientôt.)

\esp{Avísame si puedes venir.} (Préviens-moi si tu peux venir.)

\esp{Un fuerte abrazo,} / \esp{Un abrazo fuerte,} (Une forte accolade)

\esp{Tu amigo,} / \esp{Tu amiga,} + prénom (Ton ami / Ton amie)
\end{textebox}

\courssub{Rédaction individuelle : méthode pour gérer les 80 mots}

\begin{methode}[Réussir la contrainte des 80 mots en 4 étapes]
\begin{enumerate}
    \item \textbf{Rédigez au brouillon} en suivant le plan en 4 blocs. Ne comptez pas encore.
    \item \textbf{Comptez les mots} (mots graphiques). Si moins de 72 : développez le bloc 3 (ajoutez une activité, un détail concret). Si plus de 88 : supprimez les adjectifs superflus, fusionnez deux phrases courtes avec un connecteur (\esp{y, pero, porque}).
    \item \textbf{Vérifiez la cohérence du registre :} \cle{tú} partout ? Pas de \cle{usted} ? Formules amicales ? Subjonctif après \esp{Me gustaría que} / \esp{Quiero que} ?
    \item \textbf{Recopiez proprement} sur la copie finale en respectant la mise en page (alinéas, formules sur des lignes séparées).
\end{enumerate}
\end{methode}

\begin{aretenir}[Astuces pour atteindre 80 mots]
\begin{itemize}
    \item \textbf{Deux connecteurs minimum} sont obligatoires (ex. \esp{Además, Por tanto}) pour lier les blocs.
    \item \textbf{Évitez les répétitions :} variez le sujet (\esp{Yo, Nosotros, Tú, La gente, El país}).
    \item \textbf{Détaillez le programme :} au lieu de \esp{Visitaremos parques}, écrivez \esp{Visitaremos el parque de Waza para ver elefantes}.
    \item \textbf{La signature} (le prénom seul) ne compte généralement pas dans les 80 mots, mais la formule de politesse (\esp{Un abrazo,}) si.
\end{itemize}
\end{aretenir}

\courssub{Exercice résolu : lettre modèle commentée (80 mots)}

\begin{exoresolu}[Rédaction de la lettre modèle]
\textbf{Énoncé :} Rédigez la lettre type (80 mots) en utilisant la banque de phrases et le plan imposé. Présentez-la avec sa traduction française en regard.

\tcblower

\textbf{Solution détaillée (lettre modèle)}

\begin{center}
\begin{tabularx}{\linewidth}{|X|X|}
\hline
\rowcolor{popblueL} \textbf{Español (81 palabras)} & \textbf{Français} \\
\hline
\esp{Querido Pedro:} & \emph{Cher Pierre,} \\
\hline
\esp{¿Cómo estás? Hace mucho que no sé de ti. Por aquí todo bien, aunque el calor en Yaoundé es fuerte.} & \emph{Comment vas-tu ? Cela fait longtemps que je n'ai pas de tes nouvelles. Ici tout va bien, bien que la chaleur à Yaoundé soit forte.} \\
\hline
\esp{Te escribo para invitarte a pasar las vacaciones en Camerún. Me encantaría que vinieras en agosto.} & \emph{Je t'écris pour t'inviter à passer les vacances au Cameroun. J'aimerais beaucoup que tu viennes en août.} \\
\hline
\esp{Además, podríamos ir a la playa de Kribi y visitar el parque de Waza para ver elefantes.} & \emph{De plus, nous pourrions aller à la plage de Kribi et visiter le parc de Waza pour voir des éléphants.} \\
\hline
\esp{La comida es deliciosa: probarás el ndolé. La gente es muy acogedora.} & \emph{La nourriture est délicieuse : tu goûteras le ndolé. Les gens sont très accueillants.} \\
\hline
\esp{Espero tu respuesta pronto.} & \emph{J'attends ta réponse bientôt.} \\
\hline
\esp{Un abrazo fuerte,} & \emph{Une forte accolade,} \\
\hline
\esp{Tu amigo,} & \emph{Ton ami,} \\
\hline
\esp{Paul} & \emph{Paul} \\
\hline
\end{tabularx}
\end{center}

\textbf{Comptage et analyse :}
\begin{itemize}
    \item \textbf{Blocs 1 + 2 :} \esp{Querido Pedro: ¿Cómo estás? Hace mucho que no sé de ti. Por aquí todo bien, aunque el calor en Yaoundé es fuerte.} = 21 mots.
    \item \textbf{Invitation :} \esp{Te escribo para invitarte a pasar las vacaciones en Camerún. Me encantaría que vinieras en agosto.} = 17 mots.
    \item \textbf{Programme :} \esp{Además, podríamos ir a la playa de Kribi y visitar el parque de Waza para ver elefantes.} = 18 mots.
    \item \textbf{Arguments :} \esp{La comida es deliciosa: probarás el ndolé. La gente es muy acogedora.} = 12 mots.
    \item \textbf{Bloc 4 :} \esp{Espero tu respuesta pronto. Un abrazo fuerte, Tu amigo, Paul} = 11 mots.
    \item \textbf{Total :} 21 + 17 + 18 + 12 + 11 = \textbf{79 mots} (81 avec la signature comptée). Dans la tolérance $\pm 10\%$ : parfait.
\end{itemize}

\textbf{Points grammaticaux validés :}
\begin{itemize}
    \item \esp{Hace mucho que no sé de ti} : structure temporelle \cle{hace + tiempo + que} + présent.
    \item \esp{Me encantaría que vinieras} : \cle{conditionnel + subjonctif imparfait}, structure très valorisée.
    \item \esp{Podríamos ir / visitar} : conditionnel pour exprimer une proposition hypothétique.
    \item \esp{Además, La comida..., La gente...} : connecteurs + variation des sujets.
\end{itemize}
\end{exoresolu}

\courssub{Auto-évaluation avec grille de critères}

Avant de rendre votre copie, utilisez cette grille. Cochez \textbf{Oui} ou \textbf{Non}. Si vous avez un \textbf{Non}, corrigez avant la version finale.

\begin{center}
\begin{tabularx}{\linewidth}{|X|c|c|X|}
\hline
\rowcolor{popblueL} \textbf{Critère} & \textbf{Oui} & \textbf{Non} & \textbf{Action corrective si Non} \\
\hline
Structure : les 4 blocs sont respectés ? & & & Réorganiser : Saludo / Noticias / Invitación / Despedida. \\
\hline
Registre : \esp{tú} et formules amicales ? & & & Remplacer \esp{usted} par \esp{tú} ; \esp{Atentamente} $\rightarrow$ \esp{Un abrazo}. \\
\hline
Vocabulaire spécifique au Cameroun ? & & & Ajouter : Kribi, Waza, \esp{ndolé}, \esp{bendskin}, \esp{suya}. \\
\hline
Grammaire : subjonctif si nécessaire ? & & & Vérifier : \esp{quiero que vengas}, \esp{me encantaría que vinieras}. \\
\hline
Longueur : 80 mots ($\pm 10\%$) ? & & & Compter. Couper ou développer le bloc 3. \\
\hline
\end{tabularx}
\end{center}

\courssub{Correction collective au tableau : analyse de la lettre modèle}

Le professeur projette (ou réécrit) la lettre modèle de l'exercice résolu. La classe, guidée par le professeur, valide chaque critère de la grille à l'oral.

\begin{experience}[Mise en commun / Correction collective]
\textbf{Déroulement (15 min) :}
\begin{enumerate}
    \item \textbf{Lecture à voix haute} par un élève (travail phonétique : accent tonique, intonation interrogative montante dans \esp{¿Cómo estás?}).
    \item \textbf{Comptage collectif} au tableau : un élève compte, un autre note au tableau. Validation du total.
    \item \textbf{Chasse aux trésors grammaticaux :} les élèves soulignent au tableau avec des couleurs différentes :
        \begin{itemize}
            \item \textcolor{popblue}{Bleu} : les connecteurs (\esp{Además, aunque, para}).
            \item \textcolor{popgreen}{Vert} : les verbes au conditionnel (\esp{encantaría, podríamos}) et au subjonctif (\esp{vinieras}).
            \item \textcolor{poporange}{Orange} : le vocabulaire culturel camerounais (\esp{Yaoundé, Kribi, Waza, ndolé}).
        \end{itemize}
    \item \textbf{Validation de la grille :} toute la classe répond « \esp{¡Sí!} » ou « \esp{¡No!} » pour chaque critère.
\end{enumerate}
\end{experience}

\begin{attention}[Erreur fatale : la majuscule après les deux-points]
En français, on écrit : \emph{Cher Pierre,} (virgule, puis minuscule à la ligne).
En espagnol, après la formule d'appel suivie de \textbf{deux-points}, on \textbf{change de ligne} et on commence par une \textbf{MAJUSCULE}.
\begin{center}
\begin{tabular}{|l|l|}
\hline
\rowcolor{popblueL} \textbf{Correct (espagnol)} & \textbf{Incorrect (calque du français)} \\
\hline
\esp{Querido Pedro:} & \esp{Querido Pedro,} \\
\esp{¿Cómo estás?} & \esp{¿cómo estás?} \\
\hline
\esp{Estimado señor:} & \esp{Estimado señor,} \\
\esp{Le escribo...} & \esp{le escribo...} \\
\hline
\end{tabular}
\end{center}
\emph{Règle :} les deux-points (\esp{:}) isolent la formule d'appel. La phrase qui suit est une phrase autonome $\rightarrow$ majuscule obligatoire.
\end{attention}

\courssub{Variante formelle : transformation en courriel professionnel}

Compétence visée : \textbf{adapter le registre au destinataire}. Même situation de communication (une invitation), mais le destinataire est \esp{el director de una escuela asociada} (le directeur d'un établissement partenaire). Registre : \textbf{formel} (\cle{usted}).

\begin{methode}[Transformer une lettre amicale en courriel formel (5 opérations)]
\begin{enumerate}
    \item \textbf{Objet (\esp{Asunto}) :} obligatoire, clair, concis. \esp{Asunto: Invitación a un programa de intercambio en Camerún}.
    \item \textbf{Salutation :} \esp{Estimado señor Director:} / \esp{Estimada señora Directora:} (pas de prénom, \cle{usted} implicite).
    \item \textbf{Corps :} vouvoiement systématique (\esp{usted, su, le}). Style impersonnel ou \esp{nosotros} institutionnel. Supprimer les détails trop personnels (chaleur, examens). Valoriser l'institution.
    \item \textbf{Formule de clôture :} \esp{Atentamente,} / \esp{Le saluda atentamente,} / \esp{Cordiales saludos,}.
    \item \textbf{Signature complète :} prénom + nom + fonction + établissement + téléphone ou courriel.
\end{enumerate}
\end{methode}

\begin{textebox}[Modèle de courriel formel (version transformée)]
\esp{Asunto: Invitación a un programa de intercambio educativo en Camerún}

\esp{Estimado señor Director:}

\esp{Le escribo en mi calidad de profesor de español del Lycée Général Leclerc de Yaoundé para invitar a su centro a participar en un programa de intercambio durante el mes de agosto.}

\esp{Nos gustaría recibir a un grupo de estudiantes para visitar nuestra institución, descubrir la cultura camerunesa (la reserva de Waza, las playas de Kribi) y fortalecer los lazos de cooperación.}

\esp{Quedo a la espera de su respuesta para concretar los detalles.}

\esp{Atentamente,}

\esp{Paul Mbarga}\\
\esp{Profesor de Español}\\
\esp{Lycée Général Leclerc, Yaoundé}\\
\esp{Tel.: +237 6XX XX XX XX}

\emph{Traduction :} Objet : Invitation à un programme d'échange éducatif au Cameroun. Monsieur le Directeur, je vous écris en ma qualité de professeur d'espagnol du Lycée Général Leclerc de Yaoundé pour inviter votre établissement à participer à un programme d'échange pendant le mois d'août. Nous aimerions accueillir un groupe d'élèves pour visiter notre institution, découvrir la culture camerounaise (la réserve de Waza, les plages de Kribi) et renforcer les liens de coopération. Je reste dans l'attente de votre réponse pour concrétiser les détails. Cordialement, Paul Mbarga, professeur d'espagnol.
\end{textebox}

\begin{center}
\begin{tabularx}{\linewidth}{|X|X|X|}
\hline
\rowcolor{popblueL} \textbf{Élément} & \textbf{Amical (\esp{tú})} & \textbf{Formel (\esp{usted})} \\
\hline
Destinataire & \esp{Querido Pedro:} & \esp{Estimado señor Director:} \\
\hline
Verbe principal & \esp{Te invito / Me encantaría que vinieras} & \esp{Le invito / Nos gustaría recibirles} \\
\hline
Vocabulaire & \esp{playa, ndolé, abrazo, amigo} & \esp{programa, cooperación, institución, atentamente} \\
\hline
Grammaire & \esp{tú, te, tu, vinieras} (subj.) & \esp{usted, le, su, participen} (subj.) \\
\hline
Connecteurs & \esp{Además, Por cierto} & \esp{Asimismo, En consecuencia, Por tanto} \\
\hline
Clôture & \esp{Un abrazo fuerte, Tu amigo} & \esp{Atentamente,} + nom et fonction \\
\hline
\end{tabularx}
\end{center}

\courssub{Exercice d'application : à vous de jouer !}

\begin{exercice}[Entraînement individuel — devoir maison ou en classe]
\begin{enumerate}
    \item \textbf{Version amicale (80 mots) :} vous êtes \esp{Ana} (Yaoundé). Écrivez à \esp{Luis} (Madrid) pour l'inviter à la fête de fin d'année du lycée. Incluez : la date, les activités (danse \esp{bendskin}, repas), la tenue vestimentaire.
    \item \textbf{Version formelle (90-100 mots) :} vous êtes le président du Club d'espagnol de votre lycée. Écrivez un courriel à \esp{la Directora del Instituto de Cultura Hispánica de Yaoundé} pour solliciter un partenariat (don de livres, venue d'un conférencier).
\end{enumerate}
\emph{Utilisez la grille d'auto-évaluation pour vous corriger avant de rendre.}
\end{exercice}

\begin{corrige}[Exercice 1 — Proposition de correction (amical)]
\esp{Querido Luis:}

\esp{¿Qué tal? Aquí todo bien, preparando la fiesta de fin de curso. Te escribo para invitarte el sábado 15 de junio en el patio del liceo.}

\esp{Me encantaría que vinieras. Habrá música bendskin, baile y comida rica: poulet DG y plátanos. La gente vendrá con ropa tradicional colorida.}

\esp{Además, podrás conocer a mis amigos y practicar tu francés. Avísame si puedes venir.}

\esp{Un abrazo fuerte,}

\esp{Ana}

\emph{Analyse : environ 66 mots dans le corps + formules ; registre amical correct ; subjonctif \esp{vinieras} ; vocabulaire culturel ; connecteurs \esp{Además, y}. Pour atteindre 80 mots, on peut ajouter : \esp{La fiesta empezará a las cinco y terminará a medianoche.} (La fête commencera à cinq heures et finira à minuit.)}
\end{corrige}

\begin{corrige}[Exercice 2 — Proposition de correction (formel)]
\esp{Asunto: Solicitud de colaboración — Club de Español}

\esp{Estimada señora Directora:}

\esp{Le escribo en nombre del Club de Español del Lycée Général Leclerc. Queremos solicitar su colaboración para el próximo curso académico.}

\esp{Nos gustaría recibir una donación de libros de literatura juvenil y, si fuera posible, organizar una charla cultural con un ponente de su Instituto.}

\esp{Quedamos a la espera de su favorable respuesta para concertar una cita.}

\esp{Atentamente,}

\esp{María Ngono}\\
\esp{Presidenta del Club de Español}\\
\esp{Lycée Général Leclerc, Yaoundé}

\emph{Analyse : registre soutenu (\esp{usted} implicite, \esp{nosotros} institutionnel), subjonctif imparfait \esp{fuera} après \esp{si} (hypothèse), formule de clôture \esp{Atentamente}, signature complète.}
\end{corrige}

\begin{savaistu}[Le subjonctif imparfait : la « star » du Bac]
Dans la phrase modèle \esp{« Me encantaría que \textbf{vinieras} »}, le verbe \esp{vinieras} est à l'\textbf{imparfait du subjonctif}.
\begin{itemize}
    \item \textbf{Formation :} radical de la 3\textsuperscript{e} personne du pluriel du passé simple (\esp{vinieron} $\rightarrow$ \esp{vinie-}) + terminaisons \esp{-ra, -ras, -ra, -ramos, -rais, -ran}. (Une forme équivalente en \esp{-se} existe : \esp{viniese}, plus littéraire.)
    \item \textbf{Emploi :} après une expression de volonté, de souhait ou d'émotion au \textbf{conditionnel} (\esp{me gustaría, quisiera, sería bueno}) ou au \textbf{passé} (\esp{quería que vinieras}), et dans les hypothèses avec \esp{si} : \esp{Si vinieras, sería fantástico.} (Si tu venais, ce serait fantastique.)
    \item \textbf{Valeur à l'examen :} l'utilisation correcte et spontanée du subjonctif imparfait est un \textbf{marqueur de niveau B1/B2} très valorisé par les correcteurs : elle prouve la maîtrise de l'hypothèse.
    \item \textbf{Attention :} ne pas confondre avec l'imparfait de l'indicatif \esp{venías} (incorrect après \esp{querer que} + volonté).
\end{itemize}
\emph{Entraînez-vous avec les verbes irréguliers fréquents :} \esp{Quisiera que \textbf{vinieras} (venir), \textbf{tuvieras} (tener), \textbf{supieras} (saber), \textbf{hicieras} (hacer), \textbf{pudieras} (poder), \textbf{dijeras} (decir).}
\end{savaistu}

\newpage

\coursec{Activité d'intégration}

\courssub{Objectifs de l'activité}
Cette activité d'intégration vise à \cle{mobiliser l'ensemble des savoir-faire} travaillés lors de la leçon EA19 dans une tâche de communication écrite authentique, complexe et contextualisée. À l'issue de cette séance, l'élève doit être capable de :
\begin{itemize}
    \item \cle{Identifier le registre} approprié (familier / formel) en fonction du destinataire et de l'intention communicationnelle.
    \item \cle{Structurer correctement} une lettre amicale et un courriel formel (en-tête, lieu, date, formules d'appel, corps du texte, formules de clôture, signature).
    \item \cle{Appliquer les règles grammaticales} du choix entre \esp{tú} et \esp{usted} (possessifs, verbes, pronoms) et de l'impératif de politesse (\esp{usted / ustedes}).
    \item \cle{Rédiger} un texte cohérent d'environ 80 mots en espagnol correct (niveau A2+/B1), en utilisant le lexique de la correspondance et des échanges culturels.
    \item \cle{S'auto-évaluer et évaluer un pair} à l'aide d'une grille de critères précise (structure, registre, correction linguistique, richesse lexicale).
    \item \cle{Traduire} (version) un courriel formel de l'espagnol vers le français en respectant les conventions épistolaires françaises.
\end{itemize}

\courssub{Situation}
Dans le cadre d'un \cle{projet de jumelage scolaire} entre le \textbf{Lycée Leclerc de Yaoundé} (Cameroun) et le \textbf{Colegio San Carlos de Bogotá} (Colombie), et avec l'appui de la \textbf{Section consulaire de l'Ambassade d'Espagne à Yaoundé}, vous participez à une simulation de correspondance croisée. L'objectif réel est de préparer un éventuel voyage d'études et d'échange culturel pour l'été prochain.

\begin{textebox}[Annonce du projet sur l'ENT de l'établissement]
\textbf{¡PROYECTO INTERCAMBIO CAMERÚN -- COLOMBIA -- ESPAÑA!}

Estimados alumnos de 1.ª A:

El Lycée Leclerc y el Colegio San Carlos de Bogotá lanzan un programa de correspondencia y de futuro intercambio cultural para julio de 2025.

\textbf{FASE 1 (AHORA): Correspondencia cruzada.}

\textbf{Grupo A (Cartas amistosas):} escribís a vuestro correspondiente colombiano (16--17 años). Objetivo: presentaros, describir vuestro entorno (Yaoundé, la cultura, la gastronomía, el fútbol, los bendskins) e \textbf{invitarlo a descubrir Camerún}. Registro: \emph{tú}. Cierre: \emph{Un abrazo}.

\textbf{Grupo B (Correos formales):} escribís a la \textbf{Sección Consular de la Embajada de España en Yaundé}. Objetivo: solicitar información oficial sobre los \textbf{visados de estudiante}, las \textbf{becas de la AECID} y los \textbf{programas de intercambio escolar} homologados para menores de 18 años. Registro: \emph{usted}. Asunto obligatorio. Cierre: \emph{Le saluda atentamente}.

\textbf{FASE 2:} corrección entre pares, lectura de las mejores producciones y traducción (versión) del correo formal.

\emph{¡Vuestras palabras construyen puentes entre el Golfo de Guinea y los Andes!}
\end{textebox}

\emph{Traduction / Glossaire :} \esp{correspondencia cruzada} = correspondance croisée ; \esp{Sección Consular} = Section consulaire ; \esp{AECID} = Agence espagnole de coopération internationale pour le développement ; \esp{visados de estudiante} = visas d'étudiant ; \esp{homologados} = homologués, reconnus officiellement ; \esp{menores de 18 años} = mineurs de moins de 18 ans ; \esp{Golfo de Guinea} = golfe de Guinée ; \esp{puentes} = ponts.

\courssub{Consignes}
Les consignes sont conçues pour être exécutées \textbf{par binômes} (un élève « A » et un élève « B ») pendant la séance de 2 heures. Le professeur circule, guide et note les productions pour l'évaluation formative.

\begin{enumerate}
    \item \textbf{Analyse de la situation (10 min, en binôme).}
    Lisez l'annonce ci-dessus. Répondez par écrit (dans votre cahier) aux questions suivantes pour \cle{clarifier la tâche} :
    \begin{itemize}
        \item Qui est le destinataire de la lettre du Groupe A ? Quel registre (\esp{tú} / \esp{usted}) ? Quelle formule de fin ?
        \item Qui est le destinataire du courriel du Groupe B ? Quel registre ? Quelle formule de fin ? Quel élément est \textbf{obligatoire} dans l'en-tête d'un courriel alors qu'il est absent d'une lettre papier ?
        \item Citez 3 informations culturelles camerounaises que vous pourriez mentionner pour « inviter à découvrir le Cameroun ».
        \item Citez 3 informations précises que vous devez demander au Consulat (visa, bourses, programmes, hébergement, assurance...).
    \end{itemize}

    \item \textbf{Rédaction individuelle (35 min).}
    \begin{itemize}
        \item \textbf{Élève A (Lettre amicale) :} rédigez la lettre à votre correspondant colombien (prénom inventé : \esp{Juan}, \esp{María}, \esp{Camilo}...). \cle{Contrainte de longueur : environ 80 mots ($\pm$ 10 mots).} Utilisez \esp{tú}, le présent de l'indicatif, quelques verbes pronominaux (\esp{me llamo, me gusta, vivo en}) et l'impératif affirmatif à \esp{tú} pour l'invitation (\esp{¡Ven!, ¡Prueba!, ¡Visita!}). Respectez la structure : \emph{Lieu et date -- Querido/a... : -- corps -- Un abrazo -- ton prénom}.
        \item \textbf{Élève B (Courriel formel) :} rédigez le courriel adressé à \esp{consulado.yaunde@maec.es}. \cle{Contrainte de longueur : 90--110 mots.} Objet précis (ex. : \esp{Asunto: Solicitud de información -- Visado de estudiante / Intercambio escolar 2025}). Utilisez \esp{usted} (verbes à la 3\up{e} personne du singulier : \esp{quiere, necesita, puede, le escribo}), les possessifs \esp{su / sus} et le pronom \esp{le}. Employez l'impératif de politesse à \esp{usted} pour les demandes (\esp{Hágame saber, Envíeme, Indíqueme}). Structure : \emph{De / Para / Asunto -- Estimado señor / Estimada señora : -- corps -- Le saluda atentamente -- nom complet + école + téléphone}.
    \end{itemize}

    \item \textbf{Échange et correction croisée (25 min).}
    Échangez vos productions. À l'aide de la \textbf{grille d'évaluation} ci-dessous, corrigez le travail de votre camarade (surlignez les erreurs, écrivez des suggestions en marge). Attribuez une note sur 20.

    \begin{center}
    \begin{tabularx}{\linewidth}{|X|c|c|c|}\hline
    \rowcolor{popblueL} \textbf{Critère} & \textbf{4 pts} & \textbf{2 pts} & \textbf{0 pt} \\ \hline
    \textbf{Structure} & Tous les éléments obligatoires présents et bien placés. & 1--2 éléments manquants ou mal placés. & Structure incohérente ou absente. \\ \hline
    \textbf{Registre et formules} & \esp{tú / usted} respecté partout ; formules d'appel et de clôture parfaites. & 1--2 erreurs de registre (tutoiement dans le formel, vouvoiement dans l'amical) ou formules approximatives. & Registre non respecté ; formules françaises calquées. \\ \hline
    \textbf{Grammaire (verbes, accords)} & 0--1 erreur mineure (accent, accord). & 2--4 erreurs (mauvais temps, confusion \esp{ser/estar}, impératif faux). & Plus de 4 erreurs bloquantes. \\ \hline
    \textbf{Lexique et cohérence} & Vocabulaire varié (leçon + personnel), connecteurs (\esp{además, por eso, pero}). & Vocabulaire correct mais répétitif ; phrases simples juxtaposées. & Lexique très pauvre, anglicismes, calques français incompréhensibles. \\ \hline
    \textbf{Respect des consignes} & Longueur respectée ; contenu pertinent (invitation / demande précise). & Longueur hors normes ou contenu partiel. & Hors sujet. \\ \hline
    \end{tabularx}
    \end{center}

    \item \textbf{Mise en commun et oral (20 min).}
    Le professeur sélectionne 2--3 des meilleures lettres amicales et 2--3 des meilleurs courriels formels. Les auteurs (ou le binôme) les lisent à voix haute. La classe commente : \emph{« Pourquoi ce registre est-il bien tenu ? », « Quelle formule est élégante ? », « Avez-vous repéré un subjonctif ? »} (ex. : \esp{Espero que puedas venir.} « J'espère que tu pourras venir. »).

    \item \textbf{Prolongement : version (traduction) -- devoirs / prochaine séance (10 min + maison).}
    Chaque élève B remet son courriel formel (corrigé) à son binôme A. L'élève A traduit ce courriel \textbf{en français formel} (version), en respectant les conventions françaises (Objet, Madame/Monsieur, formules de politesse françaises développées : \emph{Je vous prie d'agréer, Madame, Monsieur, l'expression de mes salutations distinguées}). Remise au professeur pour notation sommative.
\end{enumerate}

\courssub{Ressources à mobiliser}
Voici un aide-mémoire synthétique à garder sous les yeux pendant la rédaction. Il reprend les points clés des sections 2, 3 et 4 de la leçon.

\begin{definition}[Formules types -- Lettre amicale (Grupo A)]
\begin{itemize}
    \item \textbf{Lieu et date :} \esp{Yaoundé, 15 de octubre de 2024} (Ville, le jour de mois de année ; mois sans majuscule).
    \item \textbf{Appel :} \esp{Querido Juan:} (masculin) / \esp{Querida María:} (féminin), suivi de \textbf{deux-points}. Jamais \esp{Señor} dans une lettre amicale.
    \item \textbf{Corps :} \esp{tú} (verbes à la 2\up{e} personne du singulier : \esp{eres, vives, te gusta, quieres}). Possessifs : \esp{tu, tus}. Pronoms : \esp{te}.
    \item \textbf{Invitation (impératif affirmatif \esp{tú}) :} \esp{¡Ven a Yaoundé!} « Viens à Yaoundé ! » ; \esp{¡Prueba el ndolé!} « Goûte le ndolé ! » ; \esp{¡Visita el mercado de Mokolo!} « Visite le marché de Mokolo ! »
    \item \textbf{Clôture :} \esp{Un abrazo,} / \esp{Un fuerte abrazo,} / \esp{Besos,} (entre filles ou famille), puis le prénom seul à la ligne.
\end{itemize}
\end{definition}

\begin{definition}[Formules types -- Courriel formel (Grupo B)]
\begin{itemize}
    \item \textbf{En-tête :} \esp{De:} (expéditeur) ; \esp{Para:} (destinataire) ; \esp{Asunto:} (objet, \textbf{obligatoire} et bref).
    \item \textbf{Appel :} \esp{Estimado señor:} / \esp{Estimada señora:} / \esp{Muy señor mío:} (très formel), suivi de deux-points.
    \item \textbf{Corps :} \esp{usted} (verbes à la 3\up{e} personne du singulier : \esp{quiere, puede, necesita}). Possessifs : \esp{su, sus}. Pronom : \esp{le}. Première phrase type : \esp{Le escribo para solicitar información sobre...} « Je vous écris pour demander des informations sur... ».
    \item \textbf{Demandes polies (impératif \esp{usted}) :} \esp{Envíeme el dossier, por favor.} « Envoyez-moi le dossier, s'il vous plaît. » ; \esp{Indíqueme los requisitos.} « Indiquez-moi les conditions requises. » ; \esp{Hágame saber las fechas.} « Faites-moi savoir les dates. »
    \item \textbf{Clôture :} \esp{Le saluda atentamente,} / \esp{Atentamente,} / \esp{Cordialmente,} puis nom complet, établissement et coordonnées.
\end{itemize}
\end{definition}

\begin{propriete}[Rappel de grammaire : \esp{tú} ou \esp{usted} ?]
\begin{itemize}
    \item \esp{tú} : verbe à la 2\up{e} personne du singulier (\esp{hablas, eres, vives}) ; possessifs \esp{tu / tus} ; pronom \esp{te}. Emploi : ami, correspondant, membre de la famille.
    \item \esp{usted} : verbe à la 3\up{e} personne du singulier (\esp{habla, es, vive}) ; possessifs \esp{su / sus} ; pronom \esp{le}. Emploi : autorité, administration, personne inconnue ou âgée.
    \item \textbf{Impératif de politesse (\esp{usted})} : on prend le subjonctif présent : \esp{hacer} $\rightarrow$ \esp{haga} ; \esp{enviar} $\rightarrow$ \esp{envíe} ; \esp{indicar} $\rightarrow$ \esp{indique} ; \esp{escribir} $\rightarrow$ \esp{escriba}. Avec un pronom, l'accent se maintient : \esp{envíe} $+$ \esp{me} $\rightarrow$ \esp{envíeme}.
    \item \textbf{Impératif affirmatif \esp{tú}} : 3\up{e} personne du singulier du présent (\esp{venir} $\rightarrow$ \esp{ven} ; \esp{probar} $\rightarrow$ \esp{prueba} ; \esp{visitar} $\rightarrow$ \esp{visita}). Irréguliers : \esp{decir} $\rightarrow$ \esp{di}, \esp{hacer} $\rightarrow$ \esp{haz}, \esp{ir} $\rightarrow$ \esp{ve}, \esp{poner} $\rightarrow$ \esp{pon}, \esp{salir} $\rightarrow$ \esp{sal}, \esp{tener} $\rightarrow$ \esp{ten}.
\end{itemize}
\end{propriete}

\begin{attention}[Pièges fréquents des francophones]
\begin{itemize}
    \item Ne calquez pas « Cher Monsieur » par \esp{Querido señor} dans une lettre formelle : on écrit \esp{Estimado señor:}. \esp{Querido} est réservé aux proches.
    \item Après l'appel, l'espagnol met des \textbf{deux-points} (\esp{Querida María:}) et la phrase suivante commence par une \textbf{majuscule} sur une nouvelle ligne.
    \item « S'il vous plaît » se traduit \esp{por favor} (jamais \esp{si usted quiere}).
    \item Attention au registre : ne mélangez jamais \esp{tú} et \esp{usted} dans le même texte (\esp{*Querido Juan, ¿cómo está usted?} est incohérent).
    \item La date : \esp{15 de octubre de 2024}, sans majuscule au mois et sans \esp{el} avant le jour dans l'en-tête de lettre.
    \item \esp{Un abrazo} ne se traduit pas mot à mot : c'est « Je t'embrasse / Amitiés ».
\end{itemize}
\end{attention}

\begin{exoresolu}[Modèle de lettre amicale (Grupo A) -- environ 85 mots]
Rédigez la lettre à votre correspondant colombien en respectant la structure et le registre \esp{tú}.
\tcblower
\begin{center}
\begin{tabularx}{\linewidth}{X}
\esp{Yaoundé, 15 de octubre de 2024}\\[4pt]
\esp{Querido Juan:}\\[4pt]
\esp{Me llamo Amina y tengo dieciséis años. Vivo en Yaoundé, la capital de Camerún, con mi familia. Me gusta mucho el fútbol y los fines de semana voy al mercado de Mokolo con mi madre.}\\
\esp{Aquí la comida es deliciosa: tienes que probar el ndolé y el poulet DG. Además, la gente es muy acogedora.}\\
\esp{¡Ven a descubrir mi país este verano! ¡Visita nuestras playas de Kribi! Espero que puedas venir pronto.}\\[4pt]
\esp{Un abrazo,}\\
\esp{Amina}\\
\end{tabularx}
\end{center}
\emph{Traduction :} Yaoundé, le 15 octobre 2024. Cher Juan : Je m'appelle Amina et j'ai seize ans. Je vis à Yaoundé, la capitale du Cameroun, avec ma famille. J'aime beaucoup le football et le week-end je vais au marché de Mokolo avec ma mère. Ici la cuisine est délicieuse : tu dois goûter le ndolé et le poulet DG. De plus, les gens sont très accueillants. Viens découvrir mon pays cet été ! Visite nos plages de Kribi ! J'espère que tu pourras venir bientôt. Je t'embrasse, Amina.
\par\smallskip
\emph{Points de langue :} \esp{tienes que probar} = « tu dois goûter » (\esp{tener que} $+$ infinitif) ; \esp{acogedora} = accueillante (accord féminin avec \esp{la gente}) ; \esp{Espero que puedas venir} : subjonctif après \esp{esperar que}.
\end{exoresolu}

\begin{exoresolu}[Modèle de courriel formel (Grupo B) -- environ 100 mots]
Rédigez le courriel au Consulat d'Espagne en respectant la structure et le registre \esp{usted}.
\tcblower
\begin{center}
\begin{tabularx}{\linewidth}{X}
\esp{De: jean.mbarga@lyceeleclerc.cm}\\
\esp{Para: consulado.yaunde@maec.es}\\
\esp{Asunto: Solicitud de información -- Visado de estudiante e intercambio escolar 2025}\\[4pt]
\esp{Estimado señor:}\\[4pt]
\esp{Soy alumno de primera del Lycée Leclerc de Yaundé. Le escribo para solicitar información sobre los visados de estudiante y las becas de la AECID para un intercambio escolar en julio de 2025.}\\
\esp{¿Podría indicarme los requisitos y los documentos necesarios? Envíeme, por favor, el formulario de solicitud. También hágame saber las fechas límite de inscripción.}\\
\esp{Muchas gracias de antemano por su ayuda.}\\[4pt]
\esp{Le saluda atentamente,}\\
\esp{Jean Mbarga -- Lycée Leclerc, Yaundé -- Tel.: 6 90 00 00 00}\\
\end{tabularx}
\end{center}
\emph{Traduction :} De : ... ; À : ... ; Objet : Demande d'information -- Visa d'étudiant et échange scolaire 2025. Monsieur, Je suis élève de première au Lycée Leclerc de Yaoundé. Je vous écris pour demander des informations sur les visas d'étudiant et les bourses de l'AECID pour un échange scolaire en juillet 2025. Pourriez-vous m'indiquer les conditions et les documents nécessaires ? Envoyez-moi, s'il vous plaît, le formulaire de demande. Faites-moi également savoir les dates limites d'inscription. Merci beaucoup d'avance pour votre aide. Je vous prie d'agréer, Monsieur, l'expression de mes salutations distinguées. Jean Mbarga...
\par\smallskip
\emph{Points de langue :} \esp{¿Podría indicarme...?} : conditionnel de politesse, plus élégant que \esp{¿Puede...?} ; \esp{de antemano} = d'avance ; \esp{las fechas límite} = les dates limites ; \esp{su ayuda} : possessif \esp{su} car on vouvoie.
\end{exoresolu}

\begin{methode}[Comment réussir la version (courriel espagnol $\rightarrow$ français formel)]
\begin{enumerate}
    \item \textbf{Lire} tout le courriel et repérer le registre (\esp{usted} $\rightarrow$ vouvoiement en français).
    \item \textbf{Traduire l'en-tête} : \esp{De} = De / Expéditeur ; \esp{Para} = À / Destinataire ; \esp{Asunto} = Objet.
    \item \textbf{Adapter les formules}, ne pas les calquer : \esp{Estimado señor:} $\rightarrow$ « Monsieur, » ; \esp{Le saluda atentamente} $\rightarrow$ « Je vous prie d'agréer, Monsieur, l'expression de mes salutations distinguées. »
    \item \textbf{Traduire les demandes polies} au conditionnel français : \esp{Envíeme} $\rightarrow$ « Veuillez m'envoyer » ou « Pourriez-vous m'envoyer ».
    \item \textbf{Relire} : vérifier les faux amis (\esp{solicitar} = demander, non « solliciter » au sens strict ; \esp{asunto} = objet/sujet, non « astuce »).
\end{enumerate}
\end{methode}

\begin{exercice}[À vous de jouer -- entraînement rapide]
\begin{enumerate}
    \item Remettez dans l'ordre les éléments d'une lettre amicale : \esp{Un abrazo, / Querida María: / Yaoundé, 3 de marzo de 2025 / Amina / ¿Cómo estás? Espero que estés bien.}
    \item Choisissez le bon registre : \esp{(¿Cómo estás? / ¿Cómo está usted?)} pour écrire au Consul ; \esp{(tu / su)} familia (dans une lettre formelle) ; \esp{(Ven / Venga)} a visitarnos (à un ami).
    \item Traduisez en espagnol formel : « Pourriez-vous m'envoyer la liste des documents nécessaires ? »
\end{enumerate}
\end{exercice}

\begin{corrige}[Exercice « À vous de jouer »]
\begin{enumerate}
    \item Ordre correct : \esp{Yaoundé, 3 de marzo de 2025} $\rightarrow$ \esp{Querida María:} $\rightarrow$ \esp{¿Cómo estás? Espero que estés bien.} $\rightarrow$ \esp{Un abrazo,} $\rightarrow$ \esp{Amina}.
    \item Au Consul : \esp{¿Cómo está usted?} ; lettre formelle : \esp{su} familia ; à un ami : \esp{Ven} a visitarnos.
    \item \esp{¿Podría enviarme la lista de los documentos necesarios?} (conditionnel de politesse $+$ pronom \esp{me} joint à l'infinitif).
\end{enumerate}
\end{corrige}

\begin{savaistu}[La correspondance, une tradition hispanique vivante]
Dans le monde hispanophone, la lettre formelle garde des formules très codifiées : \esp{Muy señor mío} ou \esp{Distinguida señora} ouvrent les courriers administratifs, et la clôture peut être longue : \esp{Quedo a su disposición y le saluda atentamente} « Je reste à votre disposition et vous salue cordialement ». En Guinée équatoriale, seul pays hispanophone d'Afrique subsaharienne, les élèves rédigent ce type de lettres en espagnol dès le collège : c'est le même savoir-faire que celui que vous venez de pratiquer entre Yaoundé et Bogotá !
\end{savaistu}

\newpage

\coursec{Méthodes et exercices résolus}

\coursec{La correspondencia y los intercambios escritos}

\courssub{Texte support : Dos mensajes, dos registros}

\begin{textebox}[Deux messages, deux registres -- Contexte : Cameroun / Espagne]
\emph{Message 1 (WhatsApp -- Registre familier / \cle{tú})}
Hola Amina:
¿Qué tal? Aquí todo bien, ya estoy de vacaciones. ¡Por fin! El examen de español fue difícil pero creo que aprobé. Mañana voy a la playa con mis primos. ¿Tú qué planes tienes? Cuéntame.
Un abrazo,
Diego

\emph{Message 2 (Courriel -- Registre formel / \cle{usted})}
Asunto: Solicitud de información sobre cursos de español

Estimado Sr. Director:
Le escribo para solicitar información sobre los cursos de español para extranjeros que ofrece su instituto para el próximo trimestre. Me interesa especialmente la modalidad intensiva.
Quedo a la espera de su respuesta.
Atentamente,
Paul Biya Mvondo
\end{textebox}

\begin{propriete}[Analyse contrastive des deux registres]
\begin{center}
\begin{tabularx}{\linewidth}{|X|X|X|}\hline
\rowcolor{popblueL} \textbf{Critère} & \textbf{Message 1 (Familier)} & \textbf{Message 2 (Formel)} \\ \hline
\textbf{Destinataire} & Ami (\esp{Diego}) & Inconnu / Supérieur (\esp{Sr. Director}) \\ \hline
\textbf{Pronom de politesse} & \cle{tú} (\esp{estás, tienes, aprobé}) & \cle{usted} (\esp{Le escribo, ofrece, interesa}) \\ \hline
\textbf{Verbes} & 2e pers. sg (\esp{estás, tienes, vas}) & 3e pers. sg (\esp{escribe, ofrece, interesa}) \\ \hline
\textbf{Ouverture} & \esp{Hola Amina:} / \esp{Querida Amina:} & \esp{Estimado Sr. Director:} \\ \hline
\textbf{Clôture} & \esp{Un abrazo,} / \esp{Besos,} & \esp{Atentamente,} / \esp{Le saluda atentamente,} \\ \hline
\textbf{Vocabulaire} & \esp{Por fin, primos, Cuéntame} & \esp{Solicitar, modalidad, Quedo a la espera} \\ \hline
\textbf{Ponctuation} & \esp{¿ ? ¡ !} obligatoires & \esp{¿ ? ¡ !} obligatoires \\ \hline
\end{tabularx}
\end{center}
\end{propriete}

\courssub{Lexique : El mundo de la correspondencia}

\begin{definition}[Vocabulaire essentiel -- Correspondance]
\begin{center}
\begin{tabularx}{\linewidth}{|l|X|X|}\hline
\rowcolor{popblueL} \textbf{Español} & \textbf{Français} & \textbf{Contexte / Exemple} \\ \hline
\esp{el remitente} & l'expéditeur & \esp{El remitente es Paul.} \\ \hline
\esp{el destinatario} & le destinataire & \esp{El destinatario es el director.} \\ \hline
\esp{el sobre} & l'enveloppe & \esp{Pon la carta en el sobre.} \\ \hline
\esp{el sello / la estampilla} & le timbre & \esp{Comprar sellos en correos.} \\ \hline
\esp{la dirección} & l'adresse & \esp{Escribe la dirección clara.} \\ \hline
\esp{el código postal} & le code postal & \esp{El código postal de Yaoundé.} \\ \hline
\esp{el asunto} & l'objet (mail) & \esp{Asunto: Solicitud de empleo.} \\ \hline
\esp{el cuerpo del mensaje} & le corps du message & \esp{En el cuerpo explico el motivo.} \\ \hline
\esp{la salutación} & la salutation & \esp{Salutación inicial: Estimado...} \\ \hline
\esp{la despedida} & la formule de politesse finale & \esp{Despedida: Atentamente.} \\ \hline
\esp{la firma} & la signature & \esp{Firma: Paul Biya.} \\ \hline
\esp{adjunto / adjunta} & pièce jointe (accordé) & \esp{Adjunto mi CV.} \\ \hline
\esp{el acuse de recibo} & l'accusé de réception & \esp{Solicito acuse de recibo.} \\ \hline
\esp{reenviar} & transférer (mail) & \esp{Reenvíame el correo, por favor.} \\ \hline
\esp{responder / contestar} & répondre & \esp{Contéstame pronto.} \\ \hline
\end{tabularx}
\end{center}
\end{definition}

\begin{savaistu}[Correspondance au Cameroun et en Guinée équatoriale]
La Guinée équatoriale (\esp{Guinea Ecuatorial}) est le seul pays africain où l'espagnol est langue officielle. À Malabo et Bata, la correspondance administrative suit les normes espagnoles (registre \esp{usted}, formules \esp{Excmo. Sr.}, \esp{Atentamente}). Au Cameroun, l'apprentissage de l'espagnol (LV2) ouvre vers la CEMAC et les échanges commerciaux (cacao, bois, pétrole). Rédiger un courriel formel en espagnol est une compétence professionnelle recherchée à Douala (port, entreprises) et Yaoundé (administrations, ambassades).
\end{savaistu}

\courssub{Grammaire appliquée : Tú, Usted et l'impératif de politesse}

\begin{propriete}[Choix du pronom : \cle{tú} vs \cle{usted}]
\begin{itemize}
    \item \textbf{\cle{tú} (familier) :} Amis, famille, camarades, enfants, animaux, Dieu. \rightarrow Verbe 2e pers. sg (\esp{tú hablas, tú eres, tú ve}).
    \item \textbf{\cle{usted} (formel, sg) / \cle{ustedes} (formel, pl -- Espagne \emph{et} AmLat pour pl formel) :} Inconnus, supérieurs, administration, professeurs, personnes âgées, contexte professionnel. \rightarrow Verbe 3e pers. sg/pl (\esp{usted habla, ustedes hablan}).
    \item \textbf{Accords :} Adjectifs et participes s'accordent avec le \emph{vrai} genre/nombre de la personne. \esp{Estimado Sr. Gómez (masc) $\rightarrow$ usted está cansado.} \esp{Estimada Sra. López (fem) $\rightarrow$ usted está cansada.}
\end{itemize}
\end{propriete}

\begin{propriete}[Impératif de politesse (\cle{usted} / \cle{ustedes}) = Subjonctif présent 3e pers.]
\begin{center}
\begin{tabular}{|l|l|l|l|}\hline
\rowcolor{popblueL} \textbf{Infinitif} & \textbf{Tú (familier)} & \textbf{Usted (formel)} & \textbf{Ustedes (pluriel)} \\ \hline
\esp{Hablar} & \esp{Habla} / \esp{No hables} & \esp{Hable} / \esp{No hable} & \esp{Hablen} / \esp{No hablen} \\ \hline
\esp{Comer} & \esp{Come} / \esp{No comas} & \esp{Coma} / \esp{No coma} & \esp{Coman} / \esp{No coman} \\ \hline
\esp{Escribir} & \esp{Escribe} / \esp{No escribas} & \esp{Escriba} / \esp{No escriba} & \esp{Escriban} / \esp{No escriban} \\ \hline
\esp{Ir (irrégulier)} & \esp{Ve} / \esp{No vayas} & \esp{Vaya} / \esp{No vaya} & \esp{Vayan} / \esp{No vayan} \\ \hline
\esp{Ser (irrégulier)} & \esp{Sé} / \esp{No seas} & \esp{Sea} / \esp{No sea} & \esp{Sean} / \esp{No sean} \\ \hline
\esp{Poner (irrégulier)} & \esp{Pon} / \esp{No pongas} & \esp{Ponga} / \esp{No ponga} & \esp{Pongan} / \esp{No pongan} \\ \hline
\esp{Salir (irrégulier)} & \esp{Sal} / \esp{No salgas} & \esp{Salga} / \esp{No salga} & \esp{Salgan} / \esp{No salgan} \\ \hline
\esp{Tener (irrégulier)} & \esp{Ten} / \esp{No tengas} & \esp{Tenga} / \esp{No tenga} & \esp{Tengan} / \esp{No tengan} \\ \hline
\esp{Venir (irrégulier)} & \esp{Ven} / \esp{No vengas} & \esp{Venga} / \esp{No venga} & \esp{Vengan} / \esp{No vengan} \\ \hline
\end{tabular}
\end{center}
\emph{Usage en correspondance :} \esp{Por favor, responda a la brevedad.} \esp{Adjunte el documento, por favor.} \esp{No se preocupe.}
\end{propriete}

\begin{attention}[Leísmo / Loísmo / Laísmo en correspondance formelle]
En Espagne (norme RAE), on utilise \cle{le} (au lieu de \emph{lo/la}) comme COD pour une personne masculine au registre formel : \esp{Le escribo a usted} (au lieu de \esp{Lo escribo}). C'est le \emph{leísmo de cortesía}, accepté au Bac. En Amérique latine, on préfère souvent \esp{Lo escribo} / \esp{La escribo}. \textbf{Conseil :} Utilise \esp{Le} pour \esp{usted} masculin (standard camerounais calqué sur l'Espagne), \esp{La} pour \esp{usted} féminin. Pour COI : toujours \esp{le} (\esp{Le envío el correo}).
\end{attention}

\courssub{Structure de la lettre et du courriel}

\begin{methode}[Rédiger une lettre amicale (correspondance personnelle)]
\textbf{Objectif :} Écrire une lettre ou un message à un ami ou un proche en respectant les codes du registre familier (\esp{tú}).
\begin{enumerate}
    \item \textbf{Analyser le destinataire et le contexte :} Ami d'enfance, correspondant, cousin ? Le ton est affectueux, spontané, moins structuré qu'un écrit formel.
    \item \textbf{Choisir les formules d'ouverture :} \esp{Querido [nombre] / Querida [nombre] :} (suivi de deux-points, saut de ligne). \emph{Variantes :} \esp{Hola [nombre] :, ¡Qué tal [nombre] !}
    \item \textbf{Structurer le corps :}
    \begin{itemize}
        \item \emph{Salut / Nouvelles :} \esp{Espero que estés bien. / ¿Qué tal todo?}
        \item \emph{Motif de la lettre :} Nouvelles personnelles, invitation, remerciements, questions.
        \item \emph{Marqueurs de lien :} \esp{Por cierto, / A propósito, / Además, / Bueno,}
    \end{itemize}
    \item \textbf{Utiliser le \cle{tú} et l'indicatif présent / passé composé / futur proche :} \esp{Te escribo porque... / He ido a la playa / Voy a ir a Madrid.}
    \item \textbf{Choisir les formules de clôture affectueuses :} \esp{Un abrazo, / Un beso, / Muchos besos, / Cuídate, / Nos vemos pronto,}
    \item \textbf{Signer :} Prénom seulement (\esp{Tu amigo, Pablo} ou simplement \esp{Pablo}).
    \item \textbf{Relire :} Vérifier les accents (\esp{estás, qué, tú}), l'accord des adjectifs (\esp{Querida Ana}), la ponctuation d'ouverture/fermeture (\esp{¿ ? ¡ !}).
\end{enumerate}
\emph{Repères lexicaux :} \esp{contar novedades, echar de menos, ilusión, ganas de, ánimo.}
\end{methode}

\begin{methode}[Rédiger une lettre ou un courriel formel (administration, entreprise, professeur)]
\textbf{Objectif :} Produire un écrit formel (candidature, réclamation, demande d'information) en respectant les conventions strictes du registre \esp{usted}.
\begin{enumerate}
    \item \textbf{Identifier le destinataire :} Nom connu ? \esp{Estimado Sr. Gómez:} / \esp{Estimada Sra. López:} Inconnu ? \esp{Estimados señores:} / \esp{A la atención de...} (objet).
    \item \textbf{Rédiger l'en-tête (lettre papier) :} Expéditeur (haut gauche), Destinataire (haut droite), Lieu et date (ex. \esp{Yaoundé, 15 de octubre de 2025}), Objet (\esp{Asunto: Solicitud de prácticas}).
    \item \textbf{Corps du message -- Structure en 3 paragraphes :}
    \begin{enumerate}
        \item \emph{Motif :} \esp{Le escribo para solicitar... / Me dirijo a usted para informarle... / Por la presente, solicito...}
        \item \emph{Arguments / Détails :} Phrases complètes, vocabulaire précis (\esp{adjunto mi currículum, quedo a su disposición}). Éviter les contractions orales.
        \item \emph{Appel à l'action / Disponibilité :} \esp{Quedo a la espera de su respuesta. / Agradezco de antemano su atención.}
    \end{enumerate}
    \item \textbf{Grammaire du registre formel :}
    \begin{itemize}
        \item \cle{Usted / Ustedes} (verbes à la 3e personne sg/pl).
        \item \cle{Infinitif de courtoisie} ou \cle{Subjonctif} pour les demandes : \esp{Le agradecería que me informara} (conditionnel + imparfait subjonctif), \esp{Le ruego que me conteste} (subjonctif présent), \esp{Por favor, envíeme...} (impératif \esp{usted}).
        \item Passif réfléchi (\esp{se ruega contestar}) ou impersonnel (\esp{es necesario}).
    \end{itemize}
    \item \textbf{Formules de clôture (obligatoires, codifiées) :}
    \begin{itemize}
        \item Standard : \esp{Atentamente,}
        \item Plus formelle : \esp{Le saluda atentamente,} / \esp{Reciba un cordial saludo,}
        \item Si réponse attendue : \esp{Quedo a la espera de sus noticias,}
    \end{itemize}
    \item \textbf{Signature :} Prénom + Nom complet. Pièces jointes : \esp{Adjunto: Currículum vitae.}
\end{enumerate}
\end{methode}

\begin{methode}[Rédiger un courriel ou un message instantané (WhatsApp / Email) : adapter le registre]
\textbf{Objectif :} Maîtriser les codes numériques : objet, concision, salutations simplifiées, signatures automatiques, distinction \esp{tú} / \esp{usted} selon le destinataire.
\begin{enumerate}
    \item \textbf{L'objet (Asunto) :} Obligatoire en email pro/académique. Court, nominal, informatif. \esp{Asunto: Duda sobre el examen / Asunto: Entrega trabajo}.
    \item \textbf{Salutation selon le registre :}
    \begin{center}
    \begin{tabularx}{\linewidth}{|X|X|X|}\hline
    \rowcolor{popblueL} \textbf{Contexte} & \textbf{Ouverture} & \textbf{Clôture} \\ \hline
    Ami / Pair (WhatsApp) & \esp{Hola / ¡Buenas! / Qué tal?} & \esp{Saludos / Chao / Nos vemos} \\ \hline
    Professeur / Adulte connu (Email) & \esp{Estimado Prof. García: / Buenas tardes, profe:} & \esp{Atentamente / Un saludo / Gracias} \\ \hline
    Administration / Inconnu (Email) & \esp{Estimados señores: / A la atención de...} & \esp{Atentamente / Cordialmente} \\ \hline
    \end{tabularx}
    \end{center}
    \item \textbf{Corps du message :}
    \begin{itemize}
        \item \emph{Email :} Phrases complètes, paragraphes aérés, pas de SMS language (\esp{xq, tkm, qtl} interdits en examen).
        \item \emph{WhatsApp / SMS :} Phrases courtes, énoncés télégraphiques acceptés, émojis contextuels (pas en examen écrit), vocatifs fréquents (\esp{Oye, / Mira,}).
    \end{itemize}
    \item \textbf{Signatures :} Auto-signature configurée (Nom, Tél, Lien LinkedIn) en pro. Prénom seul en perso.
    \item \textbf{Pièces jointes :} Mention explicite : \esp{Adjunto el documento. / Te envío la foto.}
\end{enumerate}
\emph{Astuce Bac :} Si le sujet impose un "message à un ami", utilise \esp{tú} + ouverture \esp{Hola / Querido}. Si "courriel au proviseur", \esp{usted} + \esp{Estimado Sr. Director}.
\end{methode}

\courssub{Méthodes et exercices résolus}

\begin{exoresolu}[Lettre amicale : répondre à un correspondant]
\textbf{Énoncé :} Vous recevez ce message de votre correspondant argentin, Santiago. Répondez-lui en espagnol (80--100 mots) en utilisant le registre approprié. Vous devez : le saluer, donner de vos nouvelles (lycée, famille), parler d'un projet pour les vacances (voyage au Nord du Cameroun), lui poser deux questions sur sa vie à Buenos Aires, et clore la lettre.

\begin{textebox}[Message de Santiago]
¡Hola Pablo!
¿Qué tal? Yo estoy muy bien, empezando el verano aquí en Buenos Aires. Hace mucho calor. Mis clases en la universidad terminaron la semana pasada y ahora tengo mucho tiempo libre. ¿Y tú? ¿Cómo va el instituto? ¿Tienes planes para las vacaciones? ¡Escríbeme pronto!
Un abrazo grande,
Santiago
\end{textebox}

\tcblower
\textbf{Solution détaillée :}

\textbf{1. Analyse de la situation :} Destinataire = ami (\esp{Santiago}) $\rightarrow$ registre familier (\esp{tú}). Contexte = réponse à une lettre reçue. Contraintes : 5 points obligatoires (salut, nouvelles, projet vacances, 2 questions, clôture). Longueur : ~90 mots.

\textbf{2. Choix des structures grammaticales :}
\begin{itemize}
    \item Ouverture : \esp{Querido Santiago:} (standard amical).
    \item Réaction aux nouvelles : \esp{Me alegro de que estés bien.} (Subjonctif présent \esp{estés} après \esp{alegrarse de que} -- niveau B1, mais \esp{Qué bien que estás bien} indicatif acceptable A2). On utilise l'indicatif pour la simplicité : \esp{¡Qué bien que estás bien!}
    \item Nouvelles personnelles : Présent (\esp{estudio, vivo}), Passé récent (\esp{he terminado}), Futur proche (\esp{voy a viajar}).
    \item Questions : Interrogation directe avec \esp{¿...?}.
    \item Clôture : \esp{Un abrazo fuerte,} + Prénom.
\end{itemize}

\textbf{3. Rédaction (brouillon mental) :}
\esp{Querido Santiago: ¡Hola! Qué bien que estás bien. Yo también estoy bien. El instituto va bien, pero estoy cansado porque los exámenes fueron difíciles. Mi familia está bien. Para las vacaciones, voy a viajar al norte de Camerún con mis padres para ver los parques nacionales. ¿Te gusta el verano en Buenos Aires? ¿Qué vas a hacer este verano? Escríbeme pronto. Un abrazo fuerte, Pablo.}

\textbf{4. Version finale soignée (92 mots) :}

\begin{quote}
Querido Santiago:
¡Hola! Qué bien que estás bien y que ya tienes vacaciones. Yo también estoy bien, aunque el instituto ha sido un poco agotador este trimestre. Mis exámenes terminaron ayer y ahora puedo descansar.
Para las vacaciones, tengo un plan muy bonito: voy a viajar al norte de Camerún con mis padres para visitar el parque nacional de Waza y ver elefantes y jirafas. ¡Tengo muchas ganas!
¿Qué tal el calor en Buenos Aires? ¿Vas a viajar a la playa o te quedas en la ciudad?
Escríbeme cuando puedas.
Un abrazo fuerte,
Pablo
\end{quote}

\textbf{5. Justifications des points clés :}
\begin{itemize}
    \item \esp{Querido Santiago:} (accord masculin, deux-points).
    \item \esp{esté bien / tienes vacaciones} : indicatif (faits réels).
    \item \esp{ha sido / terminaron} : passé composé / indéfini (alternance naturelle).
    \item \esp{voy a viajar} : futur proche (projet décidé).
    \item \esp{Tengo muchas ganas (de)} : structure idiomatique "avoir envie de".
    \item Questions fermées/ouvertes avec \esp{¿ ?}.
    \item \esp{Un abrazo fuerte,} : clôture standard amicale.
\end{itemize}
\end{exoresolu}

\begin{exoresolu}[Courriel formel : demande de stage]
\textbf{Énoncé :} Vous êtes élève en Première A au Lycée de Nkolbisson (Yaoundé). Vous souhaitez effectuer un stage d'observation de deux semaines en juillet dans l'entreprise \esp{« Café Camerún S.A. »} à Douala. Rédigez le courriel formel (objet, formules, corps, clôture) adressé au Directeur des Ressources Humaines, M. \esp{Rodríguez}. (100--120 mots).

\tcblower
\textbf{Solution détaillée :}

\textbf{1. Analyse :} Contexte professionnel $\rightarrow$ \esp{usted}. Destinataire connu (M. Rodríguez) $\rightarrow$ \esp{Estimado Sr. Rodríguez:}. Type : Demande de stage (\esp{solicitud de prácticas}). Ton : respectueux, clair, structuré.

\textbf{2. Structure du courriel :}
\begin{itemize}
    \item \textbf{Objet :} Clair et précis.
    \item \textbf{Salutation :} \esp{Estimado Sr. Rodríguez:}
    \item \textbf{Paragraphe 1 (Qui/Quoi) :} Présentation + Objet précis (dates, durée).
    \item \textbf{Paragraphe 2 (Pourquoi/Compétences) :} Motivation, lien avec études (espagnol, commerce), qualités.
    \item \textbf{Paragraphe 3 (Action) :} Disponibilité entretien, pièces jointes (CV, lettre motivation).
    \item \textbf{Clôture :} \esp{Atentamente,} + Nom complet.
\end{itemize}

\textbf{3. Rédaction finale :}

\begin{quote}
\textbf{Para:} rrhh@cafecamerun.com \\
\textbf{Asunto:} Solicitud de prácticas de observación -- Julio 2025

Estimado Sr. Rodríguez:

Le escribo para solicitar la posibilidad de realizar un stage de observación de dos semanas en su empresa, Café Camerún S.A., durante el mes de julio (del 1 al 15 de julio de 2025).

Soy alumno de Primera A en el Liceo de Nkolbisson (Yaoundé) y estudio español como segunda lengua desde hace tres años. Me interesa mucho el sector del comercio internacional del café y me gustaría conocer el funcionamiento de su departamento de exportación. Soy una persona seria, puntual y con muchas ganas de aprender.

Adjunto mi currículum vitae y una carta de motivación. Quedo a su entera disposición para concertar una entrevista en la fecha que usted considere oportuna.

Le saluda atentamente,
\textbf{Paul Biya Mvondo} \\
Tel: +237 6XX XX XX XX
\end{quote}

\textbf{4. Points de vigilance validés :}
\begin{itemize}
    \item \esp{Le escribo} (pas \esp{Je t'écris}), \esp{su empresa} (pas \esp{tu empresa}), \esp{usted considere} (subjonctif après \esp{fecha que} / \esp{considere} = présent subjonctif 3e pers. sg).
    \item \esp{Adjunto} (participe passé adjectivé accordé avec \emph{mi currículum} masc. sg -- correct).
    \item \esp{Quedo a su entera disposición} : formule figée très prisée.
    \item Date : \esp{1 al 15 de julio de 2025} (pas de majuscule au mois).
    \item \emph{Note sur le leísmo :} \esp{Le escribo} (COD personne masculine formelle) est la norme péninsulaire standard au Cameroun.
\end{itemize}
\end{exoresolu}

\begin{exoresolu}[Transformation de registre : du formel à l'informel]
\textbf{Énoncé :} Voici un courriel formel envoyé à un professeur. Réécrivez-le en message WhatsApp adressé à ton camarade de classe Carlos pour lui demander la même chose. Respecte les codes du SMS/WhatsApp (concision, tutoiement, ton amical) sans utiliser d'abréviations type "langage SMS" (\esp{xq, k, tkm}) interdites au Bac.

\begin{textebox}[Courriel original (Formel - Usted)]
Asunto: Duda sobre la tarea de historia

Estimado Sr. Martínez:

Le escribo para preguntarle cuál es la fecha exacta de entrega del trabajo de historia. En la plataforma no aparece el plazo y yo no anoté la fecha en clase. Le agradecería que me lo confirmara cuando pueda.

Atentamente,
Lucía
\end{textebox}

\tcblower
\textbf{Solution détaillée :}

\textbf{1. Analyse de la transformation :}
\begin{itemize}
    \item \textbf{Expéditeur/Destinataire :} Lucía $\rightarrow$ Carlos (camarade) $\Rightarrow$ \cle{Tú}.
    \item \textbf{Canal :} Email formel $\rightarrow$ WhatsApp informel $\Rightarrow$ Suppression "Asunto", "Estimado", "Atentamente".
    \item \textbf{Contenu :} Demande date rendu devoir histoire. Info manquante sur plateforme.
    \item \textbf{Contraintes :} Pas d'abréviations illisibles (\esp{xq, q}), phrases correctes mais courtes.
\end{itemize}

\textbf{2. Opérations grammaticales :}
\begin{itemize}
    \item \esp{Le escribo para preguntarle} $\rightarrow$ \esp{Te escribo para preguntarte} (ou suppression : \esp{¿Me dices...?}).
    \item \esp{cuál es la fecha} $\rightarrow$ \esp{cuál es la fecha} (inchangé).
    \item \esp{Le agradecería que me lo confirmara} (Conditionnel + Subjonctif imparfait -- très formel) $\rightarrow$ \esp{¿Me lo confirmas? / Avísame, por fa.} (Indicatif présent / Impératif \esp{tú}).
    \item \esp{cuando pueda} $\rightarrow$ \esp{cuando puedas / cuando tengas un rato}.
\end{itemize}

\textbf{3. Proposition de message WhatsApp (naturel, < 160 caractères) :}

\begin{quote}
Hola Carlos, ¿qué tal? ¿Me dices cuándo se entrega el trabajo de historia? En la plataforma no sale la fecha y no la apunté en clase. ¡Gracias! 😊
\end{quote}

\textbf{4. Commentaire :}
\begin{itemize}
    \item \esp{Hola Carlos, ¿qué tal?} : Ouverture phatique standard.
    \item \esp{¿Me dices...?} : Demande directe, polie (\esp{por fa} implicite ou explicite).
    \item \esp{no sale / no la apunté} : Vocabulaire naturel (\esp{salir} pour "apparaître", \esp{apuntar} pour "noter").
    \item \esp{¡Gracias!} : Clôture simple. L'émoji 😊 remplace la formule de politesse écrite.
\end{itemize}
\end{exoresolu}

\begin{methode}[Réussir l'épreuve de rédaction courte (80--100 mots) : « Écrire à un correspondant »]
\textbf{Objectif :} Produire un texte cohérent, structuré, linguistiquement correct, respectant le nombre de mots imposé.
\begin{enumerate}
    \item \textbf{Décoder le sujet (5 min) :} Souligner : \textbf{Destinataire} (ami/formel), \textbf{Thèmes obligatoires} (souvent 3 à 4 points : saluer, nouvelles, projet, question, clôture), \textbf{Nombre de mots}.
    \item \textbf{Planifier / Schématiser (5 min) :} Faire un mini-plan par paragraphes (1 idée = 1 paragraphe). Compter les mots visés par paragraphe (ex. 20+30+30+20).
    \item \textbf{Rédiger le brouillon (15 min) :}
    \begin{itemize}
        \item Écrire en espagnol directement (pas de traduction mot à mot depuis le français).
        \item Utiliser des \cle{connecteurs logiques} : \esp{Por eso, / Además, / Sin embargo, / Por cierto, / En cuanto a...}
        \item Varier les structures : \esp{Tengo ganas de + inf., Acabo de + inf., Voy a + inf., Me gustaría + inf.}
        \item \textbf{Compter les mots} (articles, prépositions, contractions \esp{del, al} comptent pour 1 mot chacun).
    \end{itemize}
    \item \textbf{Relire et corriger (10 min) -- Check-list « Zéro Faute » :}
    \begin{itemize}
        \item \textbf{Accents :} \esp{qué/que, sí/si, tú/tu, él/el, mí/mi, sé/se, dé/de, té/te.} (Mots monosyllabes toniques).
        \item \textbf{Accords :} Adjectifs (genre/nombre), Participes passés (si \esp{haber} auxiliaire $\rightarrow$ invariable ; si \esp{ser/estar} $\rightarrow$ accord).
        \item \textbf{Ser / Estar :} Identité/Caractère (\esp{ser}) vs État/Lieu/Action en cours (\esp{estar}).
        \item \textbf{Por / Para :} Cause/Moyen (\esp{por}) vs But/Destinataire (\esp{para}).
        \item \textbf{Subjonctif :} Après \esp{querer que, necesito que, espero que, cuando (futur), aunque (concession).}
        \item \textbf{Ponctuation :} \esp{¿ ? ¡ !} ouvrants/fermants. Majuscules début phrase / noms propres.
    \end{itemize}
    \item \textbf{Copier au net (5 min) :} Écriture lisible, paragraphes sautés.
\end{enumerate}
\end{methode}

\begin{exoresolu}[Rédaction guidée : Lettre à un correspondant (Sujet type Bac)]
\textbf{Énoncé :} Tu t'appelles Amina. Tu habites à Douala. Écris une lettre à ton correspondant mexicain, Diego (80--100 mots). Tu dois :
\begin{enumerate}
    \item Le saluer et lui demander de ses nouvelles.
    \item Lui parler de ton lycée (nom, matière préférée : l'espagnol).
    \item Lui décrire ton quartier (marché, bruit, ambiance).
    \item Lui proposer de venir te rendre visite en décembre (vacances, plage Kribi).
    \item Clore la lettre.
\end{enumerate}

\tcblower
\textbf{Solution détaillée :}

\textbf{1. Planification (Mots cibles : ~90) :}
\begin{itemize}
    \item P1 (Ouverture + Nouvelles) : ~20 mots.
    \item P2 (Lycée + Espagnol) : ~25 mots.
    \item P3 (Quartier Douala) : ~25 mots.
    \item P4 (Invitation Décembre/Kribi) : ~20 mots.
    \item P5 (Clôture) : ~5 mots.
\end{itemize}

\textbf{2. Rédaction au brouillon (contrôle structures) :}
\esp{Querido Diego: ¿Qué tal? Yo estoy bien. Estudio en el Liceo de Akwa. Mi asignatura favorita es el español porque me gusta la cultura latina. Mi barrio es ruidoso pero divertido. Hay un gran mercado y mucha gente. En diciembre hay vacaciones. ¿Quieres venir a Camerún? Podemos ir a la playa de Kribi. Es muy bonita. Escríbeme pronto. Un abrazo, Amina.}

\textbf{3. Vérification \& Ajustement (Comptage + Grammaire) :}
\begin{itemize}
    \item \esp{asignatura favorita es} (accord), \esp{me gusta} (pas \esp{me gusto}).
    \item \esp{ruidoso pero divertido} (accords masc. sg).
    \item \esp{hay vacaciones} (impersonnel).
    \item \esp{¿Quieres venir...? Podemos ir...} (Présent indicatif + Pouvoir présent).
    \item \esp{playa de Kribi} (toponyme).
\end{itemize}
Comptage : 94 mots. Parfait.

\textbf{4. Version finale :}

\begin{quote}
Querido Diego:
¿Qué tal? Espero que estés bien. Yo estoy muy bien. Estudio en el Liceo de Akwa aquí en Douala. Mi asignatura favorita es el español porque me encanta la lengua y la cultura latinoamericana.
Mi barrio es muy animado y un poco ruidoso. Hay un mercado grande donde se vende de todo: frutas, ropa, teléfonos... Siempre hay mucha gente y música por la calle.
En diciembre tenemos vacaciones. ¿Te gustaría venir a visitarnos? Podríamos ir a la playa de Kribi, el agua es cristalina y hace sol.
Escríbeme pronto.
Un abrazo fuerte,
Amina
\end{quote}

\textbf{5. Justifications pédagogiques :}
\begin{itemize}
    \item \textbf{Registre :} \esp{tú} constant (\esp{estás, estudias, encanta, tienes, podríamos}).
    \item \textbf{Grammaire :} \esp{me encanta} (verbe comme \esp{gustar}), \esp{hay} (il y a), \esp{se vende} (passif réfléchi), \esp{podríamos} (conditionnel politesse/proposition), \esp{es cristalina} (\esp{ser} caractéristique permanente).
    \item \textbf{Vocabulaire :} \esp{animado, cristalina, de todo, por la calle}.
    \item \textbf{Structure :} Paragraphes distincts, connecteurs (\esp{porque, y, también, en diciembre}).
\end{itemize}
\end{exoresolu}

\begin{methode}[Traduire un extrait de correspondance (Thème / Version) : Méthode pas à pas]
\textbf{Objectif :} Passer du français à l'espagnol (Thème) ou de l'espagnol au français (Version) sans calque, en respectant la grammaire contrastive.
\begin{enumerate}
    \item \textbf{Analyse globale :} Identifier le registre (formel/amical), le temps principal, les pièges connus (faux amis, prépositions, ser/estar).
    \item \textbf{Segmentation :} Découper en unités de sens (propositions).
    \item \textbf{Transposition unité par unité :}
    \begin{itemize}
        \item \textbf{Sujet :} Souvent omis en espagnol si clair (conjugaison marquée). \emph{Fr: « Je pense que... » $\rightarrow$ Esp: \esp{Pienso que...} (pas \esp{Yo pienso} sauf insistance).}
        \item \textbf{Verbes :} Choisir le bon mode/tense. \emph{Fr: « Il faut que tu viennes » $\rightarrow$ \esp{Tienes que venir} (obligation) ou \esp{Es necesario que vengas} (subjonctif).}
        \item \textbf{Prépositions :} \emph{Fr: « Penser à » $\rightarrow$ \esp{Pensar en}; Fr: « Décider de » $\rightarrow$ \esp{Decidir(se) a}; Fr: « Intéressé par » $\rightarrow$ \esp{Interesado en}.}
        \item \textbf{Faux amis :} \emph{Fr: « Actuellement » $\rightarrow$ \esp{Actualmente} (maintenant) vs \esp{Realmente} (vraiment). Fr: « Éventuellement » $\rightarrow$ \esp{Posiblemente / Quizás} (éventuellement) vs \esp{Eventualmente} (occasionnellement / par hasard).}
    \end{itemize}
    \item \textbf{Relecture contrastive :} Relire la phrase espagnole en se demandant : « Est-ce que ça sonne espagnol ? » (ordre mots : \esp{Adj + Nom} souvent, \esp{No} avant verbe, pronoms objets \emph{avant} verbe conjugué / \emph{après} infinitif/gérondif/impératif affirmatif).
    \item \textbf{Soigner la ponctuation finale :} Accents, \esp{¿ ¡}, majuscules.
\end{enumerate}
\end{methode}

\begin{exoresolu}[Thème / Version : Extrait lettre formelle]
\textbf{Énoncé :}
\begin{enumerate}
    \item \textbf{Thème :} Traduire en espagnol (registre formel \esp{usted}).
    « Monsieur le Directeur, Je vous écris pour vous informer que je ne pourrai pas assister à la réunion prévue mardi prochain. Je vous prie de m'excuser. Dans l'attente de votre réponse, je vous prie d'agréer, Monsieur le Directeur, mes salutations distinguées. »
    \item \textbf{Version :} Traduire en français.
    \esp{Estimado Sr. Gómez: Le agradezco su oferta de empleo. Sin embargo, he decidido aceptar otro puesto que se ajusta mejor a mis objetivos. Quedo a su disposición para cualquier duda. Atentamente, Carlos Ruiz.}
\end{enumerate}

\tcblower
\textbf{Solution détaillée :}

\textbf{1. Thème (Français $\rightarrow$ Espagnol)}

\textit{Analyse segment par segment :}
\begin{itemize}
    \item « Monsieur le Directeur » $\rightarrow$ \esp{Estimado Sr. Director:} (Formule standard, majuscule à \esp{Director} comme titre).
    \item « Je vous écris pour vous informer que... » $\rightarrow$ \esp{Me dirijo a usted para informarle de que...} (\esp{Me dirijo} évite le leísmo/loïsme ; \esp{informarle} COI). \emph{Variante :} \esp{Le escribo para informarle de que...} (Leísmo de cortesía accepté).
    \item « je ne pourrai pas assister » $\rightarrow$ \esp{no podré asistir} (Futur simple indicatif -- plus formel que futur proche).
    \item « à la réunion prévue mardi prochain » $\rightarrow$ \esp{a la reunión prevista para el próximo martes.} (\esp{prevista} accord avec \esp{reunión}, \esp{para} = but/échéance, \esp{próximo} avant nom).
    \item « Je vous prie de m'excuser » $\rightarrow$ \esp{Le ruego que me disculpe.} (Subjonctif présent \esp{disculpe} après \esp{rogar que}). Ou formule figée : \esp{Ruego me disculpe.}
    \item « Dans l'attente de votre réponse » $\rightarrow$ \esp{Quedo a la espera de su respuesta.} (Formule figée incontournable).
    \item « je vous prie d'agréer... salutations distinguées » $\rightarrow$ \esp{Le saluda atentamente,} (Standard moderne remplaçant la lourde formule française). \emph{Note :} On ne traduit pas littéralement « agréer mes salutations distinguées ».
\end{itemize}

\textbf{Version finale (Thème) :}

\begin{quote}
Estimado Sr. Director:
Me dirijo a usted para informarle de que no podré asistir a la reunión prevista para el próximo martes. Le ruego que me disculpe.
Quedo a la espera de su respuesta.
Le saluda atentamente,
[Votre Nom]
\end{quote}

\textbf{2. Version (Espagnol $\rightarrow$ Français)}

\textit{Analyse segment par segment :}
\begin{itemize}
    \item \esp{Estimado Sr. Gómez:} $\rightarrow$ Monsieur Gómez, (ou Monsieur le Directeur si titre connu, mais ici nom propre $\rightarrow$ Monsieur Gómez).
    \item \esp{Le agradezco su oferta de empleo.} $\rightarrow$ Je vous remercie de votre offre d'emploi. (\esp{Le} = COI \emph{usted}, \esp{agradezco} = je remercie).
    \item \esp{Sin embargo,} $\rightarrow$ Cependant, / Toutefois,
    \item \esp{he decidido aceptar otro puesto} $\rightarrow$ j'ai décidé d'accepter un autre poste. (\esp{he decidido} = Passé Composé $\rightarrow$ Passé Composé FR ; \esp{aceptar} infinitif après \esp{decidir}).
    \item \esp{que se ajusta mejor a mis objetivos.} $\rightarrow$ qui correspond mieux à mes objectifs. (\esp{ajustar} = s'adapter, correspondre ; \esp{se ajusta} passif réfléchi / voix moyenne $\rightarrow$ « qui correspond »).
    \item \esp{Quedo a su disposición para cualquier duda.} $\rightarrow$ Je reste à votre disposition pour toute question. (\esp{Quedo} = Je reste / Je me tiens ; \esp{para cualquier duda} = pour toute question/doute).
    \item \esp{Atentamente,} $\rightarrow$ Cordialement, / Bien à vous,
    \item \esp{Carlos Ruiz.} $\rightarrow$ Carlos Ruiz.
\end{itemize}

\textbf{Version finale (Version) :}

\begin{quote}
\textbf{Monsieur Gómez,}

Je vous remercie de votre offre d'emploi. Cependant, j'ai décidé d'accepter un autre poste qui correspond mieux à mes objectifs. Je reste à votre disposition pour toute question.

Cordialement,
\textbf{Carlos Ruiz}
\end{quote}
\end{exoresolu}

\begin{attention}[Les 5 erreurs qui coûtent des points au Bac]
\begin{enumerate}
    \item \textbf{Confusion \esp{tú} / \esp{usted} et accords qui en découlent :} Écrire \esp{Estimado Juan:} (formel) puis \esp{¿Cómo estás?} (familier) dans la même lettre. \textbf{Règle :} Un seul registre par écrit. \esp{Estimado Sr. López} $\rightarrow$ \esp{usted} (\esp{¿Cómo está?}). \esp{Querido Juan} $\rightarrow$ \esp{tú} (\esp{¿Cómo estás?}). \emph{Astuce :} Coche la case \esp{tú} ou \esp{usted} en haut de ta copie avant d'écrire.
    \item \textbf{Oubli des accents distinctifs (diacritiques) sur les monosyllabes :} \esp{tu} (ton) vs \esp{tú} (toi), \esp{mi} (mon) vs \esp{mí} (moi), \esp{si} (si) vs \esp{sí} (oui), \esp{el} (le) vs \esp{él} (lui), \esp{de} (de) vs \esp{dé} (donne), \esp{se} (se) vs \esp{sé} (je sais), \esp{te} (te) vs \esp{té} (thé). \textbf{Sanction :} 0,5 pt par accent manquant/erroné en thème/version/rédaction.
    \item \textbf{Calque structurel « \esp{Por} / \esp{Para} » :} Dire \esp{Gracias por el regalo} (Correct : Cause) mais \esp{Este regalo es por ti} (Incorrect : But/Destinataire $\rightarrow$ \esp{para ti}). Dire \esp{Voy por Madrid} (Incorrect : Destination $\rightarrow$ \esp{Voy a Madrid} ou \esp{Voy para Madrid}). \textbf{Mémo :} \esp{POR} = Cause / Moyen / Durée / Échange / Auteur (Passif). \esp{PARA} = But / Destinataire / Échéance / Opinion / Direction.
    \item \textbf{Mauvais choix \esp{Ser} / \esp{Estar} dans la description :} \esp{Mi barrio es ruidoso} (Caractère permanent = \cle{Ser}) vs \esp{Mi barrio está ruidoso hoy} (État passager = \cle{Estar}). \esp{La playa es bonita} (Caractère) vs \esp{La playa está bonita hoy} (Aspect actuel inhabituel). \textbf{Piège Bac :} « Il est mort » $\rightarrow$ \esp{Está muerto} (État résultatif = \esp{Estar} + Participe), PAS \esp{Es muerto}.
    \item \textbf{Faux amis « Éventuellement / Eventualmente » et formules de clôture « franco-espagnoles » inexistantes :}
    \begin{itemize}
        \item \esp{Eventualmente} = « occasionnellement / par hasard » (PAS « éventuellement / au cas où » $\rightarrow$ \esp{posiblemente, quizás}).
        \item Traduire littéralement « Je vous prie d'agréer... » par \esp{Le ruego que acepte mis saludos distinguidos.} \textbf{Correct :} \esp{Atentamente,} / \esp{Le saluda atentamente,} / \esp{Un cordial saludo.} Pour l'amical : \esp{Un abrazo,} / \esp{Besos,} / \esp{Nos vemos.} Pas de \esp{Amicalement,} (calque de « Amicalement »).
    \end{itemize}
\end{enumerate}
\end{attention}

\courssub{Production guidée : Escribo a mi correspondiente}

\begin{experience}[Jeu de rôle : « La entrevista de prácticas »]
\textbf{Objectif oral :} Simuler un entretien téléphonique ou en présentiel pour le stage de l'exercice résolu.
\begin{enumerate}
    \item \textbf{Binômes :} A = Étudiant (Paul), B = Responsable RH (M. Rodríguez).
    \item \textbf{Préparation (5 min) :} A prépare ses arguments (études, motivation, disponibilités). B prépare ses questions (dates, langues, compétences, pourquoi notre entreprise ?).
    \item \textbf{Jeu (3 min) :} B appelle A (ou inversement). Registre \esp{usted} obligatoire. Utiliser l'impératif de politesse (\esp{Dígame, Espere un momento, Tome nota}).
    \item \textbf{Retour (2 min) :} Échange sur la fluidité, le vocabulaire pro (\esp{prácticas, convenio, tutor, horario}), la correction \esp{usted}.
\end{enumerate}
\emph{Variante :} Écrire le compte-rendu de l'entretien par email (confirmation de stage).
\end{experience}

\begin{exercice}[Entraînement 1 : Le bon registre]
Complétez par la formule d'ouverture ou de clôture appropriée (choisir entre \esp{tú} / \esp{usted}).
\begin{enumerate}
    \item À ton cousin Carlos : \esp{\_\_\_\_\_\_\_\_\_\_ Carlos:} ... \esp{\_\_\_\_\_\_\_\_\_\_, Ana.}
    \item Au Directeur du lycée : \esp{\_\_\_\_\_\_\_\_\_\_ Sr. Director:} ... \esp{\_\_\_\_\_\_\_\_\_\_.}
    \item À ton professeur d'espagnol (email) : \esp{\_\_\_\_\_\_\_\_\_\_ Profe:} ... \esp{\_\_\_\_\_\_\_\_\_\_.}
    \item À une entreprise pour un stage (inconnu) : \esp{\_\_\_\_\_\_\_\_\_\_ señores:} ... \esp{\_\_\_\_\_\_\_\_\_\_.}
\end{enumerate}
\end{exercice}

\begin{corrige}[Exercice 1]
\begin{enumerate}
    \item \esp{Querido Carlos:} ... \esp{Un abrazo,} (ou \esp{Besos,})
    \item \esp{Estimado Sr. Director:} ... \esp{Atentamente,} (ou \esp{Le saluda atentamente,})
    \item \esp{Estimado Prof. [Apellido]:} / \esp{Buenas tardes, profe:} ... \esp{Atentamente,} / \esp{Un saludo,}
    \item \esp{Estimados señores:} ... \esp{Atentamente,}
\end{enumerate}
\end{corrige}

\begin{exercice}[Entraînement 2 : Transformation de registre]
Réécrivez les phrases suivantes du formel (\esp{usted}) vers l'informel (\esp{tú}) ou inversement.
\begin{enumerate}
    \item (Formel) \esp{Le ruego que me envíe el documento.} $\rightarrow$ (Informel) \esp{\_\_\_\_\_\_\_\_\_\_\_\_}
    \item (Informel) \esp{¿Puedes ayudarme con los deberes?} $\rightarrow$ (Formel) \esp{\_\_\_\_\_\_\_\_\_\_\_\_}
    \item (Formel) \esp{Quedo a la espera de su respuesta.} $\rightarrow$ (Informel) \esp{\_\_\_\_\_\_\_\_\_\_\_\_}
    \item (Informel) \esp{Avísame cuando llegues.} $\rightarrow$ (Formel) \esp{\_\_\_\_\_\_\_\_\_\_\_\_}
\end{enumerate}
\end{exercice}

\begin{corrige}[Exercice 2]
\begin{enumerate}
    \item \esp{Te ruego que me envíes el documento.} (ou \esp{Por favor, envíame el documento.})
    \item \esp{¿Puede ayudarme con los deberes?} (ou \esp{Le ruego que me ayude.})
    \item \esp{Avísame cuando contestes.} / \esp{Escríbeme pronto.}
    \item \esp{Le ruego que me avise cuando llegue.} (Subjonctif après \esp{cuando} futur).
\end{enumerate}
\end{corrige}

\begin{exercice}[Entraînement 3 : Rédaction libre (Type Bac)]
\textbf{Sujet :} Tu t'appelles Samuel. Tu habites à Yaoundé. Écris un courriel formel au Directeur de l'Alliance Française de Yaoundé (M. \esp{Dupont}) pour demander des informations sur les cours d'espagnol du soir (horaires, prix, inscription). Tu dois utiliser le registre formel (\esp{usted}). (90--110 mots).
\textit{Indications :} Objet clair. Formules d'ouverture/clôture. 3 paragraphes. Pièce jointe (CV) mentionnée.
\end{exercice}

\begin{corrige}[Exercice 3 -- Proposition de corrigé]
\begin{quote}
\textbf{Para:} direccion@afyaounde.org \\
\textbf{Asunto:} Solicitud de información -- Cursos de español vespertinos

Estimado Sr. Dupont:

Le escribo para solicitar información sobre los cursos de español para adultos que ofrece la Alianza Francesa para el próximo trimestre. Me interesa especialmente la modalidad vespertina (después de las 18h) debido a mi horario laboral.

¿Podría indicarme los horarios disponibles, el precio de la matrícula y el procedimiento de inscripción? Adjunto mi currículum vitae para que conozcan mi nivel actual (A2/B1).

Quedo a la espera de su respuesta.
Le saluda atentamente,
\textbf{Samuel Nguema} \\
Tel: +237 6XX XX XX XX
\end{quote}
\end{corrige}

\begin{savaistu}[L'étiquette épistolaire dans le monde hispanophone]
En Espagne, l'en-tête d'une lettre formelle place l'expéditeur en haut à gauche, le destinataire en haut à droite, la date à droite (ex. \esp{Madrid, 12 de octubre de 2025}). L'objet (\esp{Asunto:}) est en gras, aligné à gauche, avant la salutation. En Amérique latine, on omet souvent l'adresse de l'expéditeur sur la lettre (présente sur l'enveloppe). Au Cameroun, pour un courriel administratif, on suit la norme internationale : Objet clair, \esp{Estimado...}, corps structuré, \esp{Atentamente}. L'usage de \esp{Estimado} + Prénom (\esp{Estimado Juan:}) se répand dans les emails professionnels moins formels, mais \esp{Estimado Sr. Apellido} reste la norme sûre pour le Bac.
\end{savaistu}

\newpage

\coursec{Exercices}

\courssub{Je vérifie mes connaissances}

\begin{exercice}[QCM : registres et formules]
Pour chaque situation, choisissez la formule d'ouverture ou de clôture la plus appropriée.
\begin{enumerate}
\item Vous écrivez un courriel à votre ami Carlos pour lui donner des nouvelles.
      \begin{itemize}
      \item[a)] Estimado señor Carlos:
      \item[b)] Querido Carlos:
      \item[c)] Muy señor mío:
      \end{itemize}
\item Vous terminez une lettre formelle adressée au directeur d'un lycée à Madrid.
      \begin{itemize}
      \item[a)] Un abrazo fuerte,
      \item[b)] Besos,
      \item[c)] Le saluda atentamente,
      \end{itemize}
\item Vous commencez un courriel formel pour demander une bourse d'études à une fondation.
      \begin{itemize}
      \item[a)] Hola:
      \item[b)] Estimados señores:
      \item[c)] Queridos amigos:
      \end{itemize}
\item Vous finissez un message WhatsApp à votre correspondante mexicaine Ana.
      \begin{itemize}
      \item[a)] Atentamente,
      \item[b)] Un abrazo,
      \item[c)] Cordialmente,
      \end{itemize}
\end{enumerate}
\end{exercice}

\begin{exercice}[Vrai ou faux justifié]
Déterminez si les affirmations suivantes sont vraies (V) ou fausses (F). Justifiez votre réponse en français en une phrase.
\begin{enumerate}
\item Dans une lettre amicale, on utilise toujours le pronom \emph{usted} par respect.
\item La formule \esp{« Estimado señor: »} s'emploie pour écrire à un ami proche.
\item Le courriel formel conserve la structure classique : date, salutation, corps, formule de politesse, signature.
\item \esp{« Un abrazo »} est une formule de clôture adaptée à une lettre administrative.
\item En espagnol, la date s'écrit généralement : \esp{Yaoundé, 15 de mayo de 2025}.
\end{enumerate}
\end{exercice}

\begin{exercice}[Vocabulaire : la correspondance]
Complétez le tableau suivant en associant chaque terme espagnol à sa définition ou à son équivalent français.
\begin{center}
\begin{tabularx}{\linewidth}{|X|X|X|}
\hline
\rowcolor{popblueL} \textbf{Español} & \textbf{Français / Définition} & \textbf{Type de correspondance} \\ \hline
\esp{Remitente} & & \\ \hline
\esp{Destinatario} & & \\ \hline
\esp{Asunto} & & \\ \hline
\esp{Cuerpo de la carta} & & \\ \hline
\esp{Despedida} & & \\ \hline
\esp{Firma} & & \\ \hline
\esp{Adjunto} & & \\ \hline
\esp{Acuse de recibo} & & \\ \hline
\end{tabularx}
\end{center}
\end{exercice}

\begin{exercice}[Conjugaison : impératif de politesse (\emph{usted} / \emph{ustedes})]
Conjuguez les verbes entre parenthèses à l'impératif affirmatif de politesse (\emph{usted} ou \emph{ustedes} selon le contexte).
\begin{enumerate}
\item \esp{(Escribir) \_\_\_\_\_\_\_\_\_\_\_\_ a mí pronto, por favor.} (à une personne)
\item \esp{(Dime / Dígame) \_\_\_\_\_\_\_\_\_\_\_\_ la verdad, señor Director.}
\item \esp{(Enviar) \_\_\_\_\_\_\_\_\_\_\_\_ los documentos mañana.} (à plusieurs personnes)
\item \esp{(Esperar) \_\_\_\_\_\_\_\_\_\_\_\_ un momento, por favor.}
\item \esp{(Llamar) \_\_\_\_\_\_\_\_\_\_\_\_ al teléfono 222 33 44.} (à une personne)
\item \esp{(Traer) \_\_\_\_\_\_\_\_\_\_\_\_ el pasaporte el lunes.} (à plusieurs personnes)
\end{enumerate}
\end{exercice}

\courssub{Je m'entraîne}

\begin{exercice}[Traduction : du français vers l'espagnol (thème)]
Traduisez les phrases suivantes en espagnol en adaptant le registre (tú / usted).
\begin{enumerate}
\item Écris-moi une lettre quand tu arrives à Yaoundé. (amical)
\item Monsieur, dites-moi la date de l'examen, s'il vous plaît. (formel)
\item Envoyez-moi le formulaire par courriel, s'il vous plaît. (formel, à une personne)
\item Appelle-moi ce soir, on parlera du voyage. (amical)
\item Veuillez attendre ici un instant, messieurs. (formel, à plusieurs personnes)
\item Donnez-moi votre adresse e-mail pour rester en contact. (amical)
\end{enumerate}
\end{exercice}

\begin{exercice}[Transformation de registre]
Transformez la lettre amicale suivante en courriel formel adressé au directeur d'une école. Changez les pronoms, les verbes, les formules d'ouverture et de clôture, et le vocabulaire si nécessaire.
\begin{textebox}[Lettre amicale originale]
Querido Pablo:

Te escribo para contarte que llegaré a Madrid el 10 de julio. Quiero verte y pasar una semana contigo. Por favor, dime si tienes tiempo. Escríbeme pronto.

Un abrazo fuerte,

Marie
\end{textebox}
\end{exercice}

\begin{exercice}[Texte à trous : lettre formelle]
Complétez la lettre formelle ci-dessous avec les éléments manquants choisis parmi la liste : \esp{Estimado señor Director: / Le saluda atentamente / Yaoundé, 20 de octubre de 2025 / Adjunto / Asunto: Solicitud de información / Me dirijo a usted para / Quedo a la espera de su respuesta / Atentamente / Firma: Jean Mbarga}.
\begin{textebox}[Lettre formelle à compléter]
\_\_\_\_\_\_\_\_\_\_\_\_

\_\_\_\_\_\_\_\_\_\_\_\_

\_\_\_\_\_\_\_\_\_\_\_\_

\_\_\_\_\_\_\_\_\_\_\_\_, por la presente le solicito información sobre el programa de intercambio escolar entre Camerún y España. \_\_\_\_\_\_\_\_\_\_\_\_ mi currículum vitae y una carta de motivación.

\_\_\_\_\_\_\_\_\_\_\_\_ de sus noticias.

\_\_\_\_\_\_\_\_\_\_\_\_

\_\_\_\_\_\_\_\_\_\_\_\_
\end{textebox}
\end{exercice}

\begin{exercice}[Appariement : parties d'une lettre / courriel]
Remettez dans l'ordre logique les parties suivantes d'une lettre formelle en numérotant de 1 à 7.
\begin{enumerate}
\item \esp{Firma}
\item \esp{Fecha y lugar}
\item \esp{Cuerpo de la carta}
\item \esp{Saludo inicial (ej. Estimado señor:)}
\item \esp{Despedida (ej. Le saluda atentamente,)}
\item \esp{Asunto}
\item \esp{Datos del destinatario (nombre, cargo, dirección)}
\end{enumerate}
\end{exercice}

\courssub{Je me prépare au Bac}

\begin{exercice}[Compréhension et langue : courriel formel]
Lisez le courriel suivant et répondez aux questions en français.
\begin{textebox}[Courriel formel]
Asunto: Solicitud de beca de estudios « Excelencia 2025 »

Estimados señores de la Fundación « Puente Hispano »:

Por la presente, me dirijo a ustedes para solicitar información sobre la convocatoria de becas « Excelencia 2025 » destinadas a estudiantes africanos de lengua francesa que deseen cursar un máster en España.

Soy estudiante de Primera A en el Liceo Bilingüe de Yaoundé, tengo 17 años y mi expediente académico es excelente. Adjunto mi currículum vitae, mis notas certificadas y una carta de motivación en español.

Les agradecería que me indicaran los plazos de inscripción, los requisitos específicos y el procedimiento de selección.

Quedo a la espera de su amable respuesta.

Les saluda atentamente,

\emph{Pauline Ngo}
Tel.: +237 6XX XX XX XX
Correo: pauline.ngo@email.cm
\end{textebox}
\begin{enumerate}
\item Quel est l'objet de ce courriel ?
\item Identifiez l'expéditeur (nom, niveau, établissement, ville).
\item Quels documents sont joints au courriel ? (au moins trois)
\item Quelles informations précises demande Pauline ?
\item Relevez dans le texte deux formules de politesse caractéristiques du registre formel.
\item Trouvez dans le texte l'équivalent de : « je vous serais reconnaissante de bien vouloir m'indiquer ».
\item Conjuguez à l'indicatif présent : \esp{yo \_\_\_\_\_\_\_\_\_\_\_\_ (dirigir) ; yo \_\_\_\_\_\_\_\_\_\_\_\_ (adjuntar) ; yo \_\_\_\_\_\_\_\_\_\_\_\_ (agradecer) ; yo \_\_\_\_\_\_\_\_\_\_\_\_ (quedar)}.
\item Transformez la phrase \esp{« Les agradecería que me indicaran… »} en remplaçant \esp{indicaran} par l'infinitif (structure \esp{agradecer + infinitif}).
\end{enumerate}
\end{exercice}

\begin{exercice}[Thème et version]
\begin{enumerate}
\item \textbf{Thème (français $\rightarrow$ espagnol) :} Traduisez en espagnol (registre formel, \emph{usted}).
\begin{quote}
Monsieur le Directeur,
Je vous écris pour demander des informations sur l'échange scolaire entre le Cameroun et l'Espagne. Je suis élève de Première A à Yaoundé. Dans l'attente de votre réponse, je vous prie d'agréer, Monsieur, mes salutations distinguées.
\end{quote}
\item \textbf{Version (espagnol $\rightarrow$ français) :} Traduisez en français le courriel suivant.
\begin{textebox}[Courriel à traduire]
Estimada señora Directora:

Le escribo para confirmar mi participación en el programa de voluntariado en Guatemala el próximo mes de julio. Adjunto encontrará el formulario de inscripción debidamente cumplimentado y una copia de mi pasaporte.

Agradezco de antemano su confirmación y quedo a su disposición para cualquier información adicional.

Atentamente,

Carlos Méndez
\end{textebox}
\end{enumerate}
\end{exercice}

\begin{exercice}[Expression écrite : lettre amicale (sujet type Bac)]
\end{exercice}

\begin{methode}[Consignes de rédaction]
\begin{enumerate}
\item Respectez la structure de la lettre amicale : lieu + date, salutation, corps, formule de clôture, signature (prénom seulement).
\item Utilisez le \emph{tú} et le registre familier / amical.
\item Incluez les informations demandées dans le sujet.
\item Nombre de mots visé : \textbf{80 mots environ} (tolérance $\pm$ 10\%).
\item Relisez : accords, accents, ponctuation espagnole (¡ ¿), vocabulaire varié.
\end{enumerate}
\end{methode}
\textbf{Sujet :} Tu as un ami espagnol, Pablo, qui veut te rendre visite au Cameroun cet été. Écris-lui une lettre (environ 80 mots) pour lui donner des informations pratiques sur le voyage (vol, visa, vaccins, climat, accueil, activités prévues) et lui exprimer ta joie de le recevoir.


\newpage

\coursec{Corrigés des exercices}

\begin{corrige}[Exercice 1 — QCM : registres et formules]
\begin{enumerate}
\item \textbf{b) Querido Carlos:} — c'est un ami, donc registre amical. \esp{Estimado señor} et \esp{Muy señor mío} sont réservés au registre formel.
\item \textbf{c) Le saluda atentamente,} — clôture formelle obligatoire avec un directeur. \esp{Un abrazo} et \esp{Besos} sont familiers.
\item \textbf{b) Estimados señores:} — on ne connaît pas les destinataires de la fondation : salutation formelle au pluriel. \esp{Hola} et \esp{Queridos amigos} sont familiers.
\item \textbf{b) Un abrazo,} — message WhatsApp à une amie : registre amical. \esp{Atentamente} et \esp{Cordialmente} sont formels.
\end{enumerate}
\textbf{Point de vigilance :} le choix de la formule dépend du \cle{destinataire} et du \cle{registre}, jamais du contenu de la lettre.
\end{corrige}

\begin{corrige}[Exercice 2 — Vrai ou faux justifié]
\begin{enumerate}
\item \textbf{F} — Dans une lettre amicale, on utilise \esp{tú} (ou \esp{vosotros}) ; \esp{usted} est réservé au registre formel.
\item \textbf{F} — \esp{Estimado señor:} est une salutation formelle ; pour un ami proche on écrit \esp{Querido + prénom:}.
\item \textbf{V} — Le courriel formel reprend la structure de la lettre classique : date, salutation, corps, formule de politesse, signature (avec en plus l'\esp{asunto}).
\item \textbf{F} — \esp{Un abrazo} est une clôture amicale ; une lettre administrative exige \esp{Atentamente} ou \esp{Le saluda atentamente}.
\item \textbf{V} — En espagnol, la date suit le schéma : lieu, jour + \esp{de} + mois (en minuscules) + \esp{de} + année : \esp{Yaoundé, 15 de mayo de 2025}.
\end{enumerate}
\end{corrige}

\begin{corrige}[Exercice 3 — Vocabulaire : la correspondance]
\begin{center}
\begin{tabularx}{\linewidth}{|X|X|X|}
\hline
\rowcolor{popblueL} \textbf{Español} & \textbf{Français / Définition} & \textbf{Type de correspondance} \\ \hline
\esp{Remitente} & Expéditeur (celui qui envoie) & Les deux \\ \hline
\esp{Destinatario} & Destinataire (celui qui reçoit) & Les deux \\ \hline
\esp{Asunto} & Objet du courriel & Courriel \\ \hline
\esp{Cuerpo de la carta} & Corps de la lettre (le texte principal) & Les deux \\ \hline
\esp{Despedida} & Formule de clôture / au revoir & Les deux \\ \hline
\esp{Firma} & Signature & Les deux \\ \hline
\esp{Adjunto} & Pièce jointe & Courriel (et lettre formelle) \\ \hline
\esp{Acuse de recibo} & Accusé de réception & Formelle / administrative \\ \hline
\end{tabularx}
\end{center}
\textbf{Point de vigilance :} ne pas confondre \esp{remitente} (expéditeur) et \esp{destinatario} (destinataire) ; \esp{asunto} signifie « objet », pas « affaire » au sens commercial ici.
\end{corrige}

\begin{corrige}[Exercice 4 — Conjugaison : impératif de politesse]
\begin{enumerate}
\item \esp{Escríbame a mí pronto, por favor.} — \emph{usted}, verbe en \emph{-ir} : radical \esp{escrib-} + \esp{a} + accent sur l'antépénultième syllabe à cause du pronom \esp{me}.
\item \esp{Dígame la verdad, señor Director.} — forme irrégulière de \esp{decir} (subjonctif \esp{diga} + \esp{me}).
\item \esp{Envíen los documentos mañana.} — \emph{ustedes} : \esp{envíen} (accent sur le \emph{í}).
\item \esp{Espere un momento, por favor.} — \emph{usted}, verbe en \emph{-ar} : \esp{espere}.
\item \esp{Llame al teléfono 222 33 44.} — \emph{usted} : \esp{llame}.
\item \esp{Traigan el pasaporte el lunes.} — \emph{ustedes}, verbe irrégulier : subjonctif \esp{traiga} $\rightarrow$ \esp{traigan}.
\end{enumerate}
\textbf{Rappel de la règle :} l'impératif de politesse = subjonctif présent de la 3\up{e} personne : \esp{-ar} $\rightarrow$ \esp{-e/-en} ; \esp{-er/-ir} $\rightarrow$ \esp{-a/-an}. Avec un pronom joint, on ajoute un accent écrit : \esp{escríbame}, \esp{dígame}.
\end{corrige}

\begin{corrige}[Exercice 5 — Traduction : thème]
\begin{enumerate}
\item \esp{Escríbeme una carta cuando llegues a Yaoundé.} — impératif \emph{tú} de \esp{escribir} + \esp{me} ; \esp{cuando} + subjonctif (\esp{llegues}) pour un futur.
\item \esp{Señor, dígame la fecha del examen, por favor.} — impératif \emph{usted} irrégulier : \esp{dígame}.
\item \esp{Envíeme el formulario por correo electrónico, por favor.} — \emph{usted} : \esp{envíe} + \esp{me} $\rightarrow$ \esp{envíeme}.
\item \esp{Llámame esta noche, hablaremos del viaje.} — \emph{tú} : \esp{llama} + \esp{me} $\rightarrow$ \esp{llámame} ; futur simple \esp{hablaremos}.
\item \esp{Esperen aquí un momento, señores.} — \emph{ustedes} : \esp{esperen} ; « veuillez » se rend simplement par l'impératif de politesse.
\item \esp{Dame tu dirección de correo electrónico para seguir en contacto.} — \emph{tú} : \esp{da} + \esp{me} $\rightarrow$ \esp{dame} ; possessif \esp{tu} sans accent (à ne pas confondre avec \esp{tú}).
\end{enumerate}
\textbf{Points de vigilance :} l'accent écrit quand le pronom est joint à l'impératif ; \esp{por favor} en fin de phrase ; ne pas calquer « s'il vous plaît » par \esp{si usted quiere}.
\end{corrige}

\begin{corrige}[Exercice 6 — Transformation de registre]
\textbf{Proposition de courriel formel :}
\begin{textebox}[Courriel formel transformé]
Asunto: Solicitud de información sobre una visita

Estimado señor Director:

Le escribo para informarle de que llegaré a Madrid el 10 de julio. Me gustaría visitar su escuela y pasar una semana en la ciudad. Le agradecería que me indicara si usted está disponible durante ese período.

Quedo a la espera de su respuesta.

Le saluda atentamente,

Marie
\end{textebox}
\textbf{Transformations attendues :}
\begin{itemize}
\item \esp{Querido Pablo} $\rightarrow$ \esp{Estimado señor Director:} ; ajout de l'\esp{Asunto}.
\item \esp{Te escribo} $\rightarrow$ \esp{Le escribo} ; \esp{dime} $\rightarrow$ \esp{indíqueme} ou \esp{me indicara} (subjonctif après \esp{agradecería que}) ; \esp{tienes} $\rightarrow$ \esp{tiene / está disponible}.
\item \esp{Un abrazo fuerte} $\rightarrow$ \esp{Le saluda atentamente}.
\item Ton plus soutenu : \esp{Quiero verte} $\rightarrow$ \esp{Me gustaría visitar su escuela} (conditionnel de politesse).
\end{itemize}
\textbf{Barème indicatif (10 pts) :} salutation et clôture formelles (2) ; passage au \esp{usted} partout (3) ; vocabulaire adapté (2) ; correction grammaticale (2) ; structure du courriel avec asunto (1).
\end{corrige}

\begin{corrige}[Exercice 7 — Texte à trous : lettre formelle]
Ordre des éléments :
\begin{enumerate}
\item \esp{Yaoundé, 20 de octubre de 2025} (lieu et date, en haut).
\item \esp{Asunto: Solicitud de información} (objet).
\item \esp{Estimado señor Director:} (salutation).
\item \esp{Me dirijo a usted para}, por la presente le solicito información… (formule d'introduction).
\item \esp{Adjunto} mi currículum vitae y una carta de motivación. (pièces jointes).
\item \esp{Quedo a la espera} de sus noticias. (formule d'attente).
\item \esp{Le saluda atentamente,} (clôture) puis \esp{Firma: Jean Mbarga}.
\end{enumerate}
\textbf{Éléments non utilisés :} \esp{Atentamente} (doublon possible mais \esp{Le saluda atentamente} est attendu ici) — vérifier la cohérence : une seule formule de clôture.
\end{corrige}

\begin{corrige}[Exercice 8 — Appariement : parties d'une lettre]
Ordre logique d'une lettre formelle :
\begin{enumerate}
\item \esp{Fecha y lugar} (1)
\item \esp{Datos del destinatario (nombre, cargo, dirección)} (2)
\item \esp{Asunto} (3)
\item \esp{Saludo inicial (ej. Estimado señor:)} (4)
\item \esp{Cuerpo de la carta} (5)
\item \esp{Despedida (ej. Le saluda atentamente,)} (6)
\item \esp{Firma} (7)
\end{enumerate}
\textbf{Astuce :} retenir le schéma « date $\rightarrow$ destinataire $\rightarrow$ objet $\rightarrow$ salutation $\rightarrow$ corps $\rightarrow$ clôture $\rightarrow$ signature ».
\end{corrige}

\begin{corrige}[Exercice 9 — Compréhension et langue : courriel formel]
\begin{enumerate}
\item L'objet du courriel est une demande d'information sur la bourse d'études \esp{« Excelencia 2025 »} pour faire un master en Espagne.
\item Expéditrice : Pauline Ngo, élève de Première A (17 ans) au Lycée Bilingue de Yaoundé (Cameroun).
\item Documents joints : son curriculum vitae, ses notes certifiées (relevés de notes) et une lettre de motivation en espagnol.
\item Elle demande : les délais d'inscription (\esp{plazos de inscripción}), les conditions requises (\esp{requisitos específicos}) et la procédure de sélection.
\item Deux formules de politesse : \esp{« Les agradecería que me indicaran… »} et \esp{« Quedo a la espera de su amable respuesta »} (ou \esp{« Les saluda atentamente »}).
\item « Je vous serais reconnaissante de bien vouloir m'indiquer » = \esp{« Les agradecería que me indicaran »}.
\item \esp{yo dirijo (dirigir) ; yo adjunto (adjuntar) ; yo agradezco (agradecer) ; yo quedo (quedar)}. Attention : \esp{dirigir} $\rightarrow$ \esp{dirijo} et \esp{agradecer} $\rightarrow$ \esp{agradezco} (irréguliers à la 1\up{re} personne).
\item \esp{« Les agradecería indicarme los plazos de inscripción… »} — structure \esp{agradecer + infinitif} avec pronom joint (\esp{indicar} + \esp{me} $\rightarrow$ \esp{indicarme}).
\end{enumerate}
\textbf{Barème indicatif (16 pts) :} 2 pts par question ; exiger la citation exacte en espagnol pour les questions 5 et 6.
\end{corrige}

\begin{corrige}[Exercice 10 — Thème et version]
\textbf{1. Thème — proposition de traduction :}
\begin{textebox}[Traduction du thème]
Muy señor mío:

Le escribo para pedir información sobre el intercambio escolar entre Camerún y España. Soy alumno de Primera A en Yaoundé. En espera de su respuesta, le ruego acepte, señor, mis saludos más atentos.
\end{textebox}
Variante de clôture acceptée : \esp{« Le saluda atentamente »} ou \esp{« Reciba, señor Director, mis más cordiales saludos »}.
\textbf{Points de vigilance :} \esp{pedir información} (pas \esp{demandar}) ; \esp{Camerún} avec accent ; \esp{alumno de Primera A} ; le subjonctif \esp{acepte} après \esp{le ruego}.

\textbf{2. Version — proposition de traduction :}
\begin{quote}
Madame la Directrice,

Je vous écris pour confirmer ma participation au programme de volontariat au Guatemala au mois de juillet prochain. Vous trouverez ci-joint le formulaire d'inscription dûment rempli ainsi qu'une copie de mon passeport.

Je vous remercie d'avance de votre confirmation et reste à votre disposition pour tout renseignement complémentaire.

Cordialement,

Carlos Méndez
\end{quote}
\textbf{Points de vigilance :} \esp{Adjunto encontrará} = « vous trouverez ci-joint » (adverbe en tête) ; \esp{debidamente cumplimentado} = « dûment rempli » ; \esp{de antemano} = « d'avance » ; \esp{quedo a su disposición} = « je reste à votre disposition ».
\textbf{Barème indicatif (20 pts) :} thème 10 pts (exactitude grammaticale 5, registre formel 3, vocabulaire 2) ; version 10 pts (fidélité 5, qualité du français 3, formules 2).
\end{corrige}

\begin{corrige}[Exercice 11 — Expression écrite : lettre amicale]
\textbf{Proposition de rédaction modèle (environ 85 mots) :}
\begin{textebox}[Lettre amicale modèle]
Yaoundé, 5 de junio de 2025

Querido Pablo:

¡Qué alegría saber que vienes a Camerún este verano! Para el viaje, hay vuelos directos desde Madrid hasta Yaoundé. Necesitas un visado y la vacuna contra la fiebre amarilla. Aquí hace calor, pero en julio llueve mucho: trae ropa ligera y un impermeable.

Te esperaré en el aeropuerto. Visitaremos el mercado central, probaremos la comida camerunesa e iremos un fin de semana a Kribi, a la playa.

¡Tengo muchas ganas de verte!

Un abrazo fuerte,

Marie
\end{textebox}
\textbf{Barème indicatif (20 pts) :}
\begin{itemize}
\item Respect de la structure de la lettre amicale (date, salutation, clôture, prénom) : 4 pts.
\item Utilisation correcte du \esp{tú} et du registre amical : 4 pts.
\item Informations pratiques demandées (vol, visa, vaccins, climat, accueil, activités) : 6 pts.
\item Correction grammaticale (accords, conjugaisons, accents) : 4 pts.
\item Respect du nombre de mots (80 $\pm$ 10\%) et présentation : 2 pts.
\end{itemize}
\textbf{Points de vigilance :}
\begin{itemize}
\item \esp{¡Qué alegría!} avec point d'exclamation ouvrant ; \esp{Tengo muchas ganas de verte} pour « j'ai hâte de te voir » (ne pas calquer le français).
\item \esp{visado} (Espagne) ou \esp{visa} (Amérique) ; \esp{la vacuna contra la fiebre amarilla}.
\item Mois sans majuscule : \esp{julio} ; date : \esp{Yaoundé, 5 de junio de 2025}.
\item Éviter les calques : « bienvenue » = \esp{bienvenido} (pas \esp{bienvenida} pour Pablo) ; « je t'attendrai » = \esp{te esperaré}.
\end{itemize}
\end{corrige}

\newpage

\coursec{Fiche bilan}

\begin{popfigure}
\begin{tikzpicture}[
    mindmap,
    concept color=popblue,
    level 1 concept/.append style={level distance=4.5cm, sibling angle=360/7},
    level 2 concept/.append style={level distance=3cm},
    every node/.style={concept, execute at begin node=\small\bfseries, text width=2.8cm, align=center},
    grow cyclic,
    text=white
]
\node[concept, fill=popblue, align=center]{Correspondencia\\ e intercambios}
    child[concept color=poporange] { node[align=center]{Carta\\ amistosa} 
        child[concept color=poporange!70!white] { node {Querido/a...} }
        child[concept color=poporange!70!white] { node {Un abrazo} }
        child[concept color=poporange!70!white] { node {Tú} }
    }
    child[concept color=popgreen] { node[align=center]{Carta\\ formal} 
        child[concept color=popgreen!70!white] { node {Estimado Sr.} }
        child[concept color=popgreen!70!white] { node {Atentamente} }
        child[concept color=popgreen!70!white] { node {Usted} }
    }
    child[concept color=poppurple] { node[align=center]{Correo\\ electrónico} 
        child[concept color=poppurple!70!white] { node {Asunto} }
        child[concept color=poppurple!70!white] { node {Cuerpo} }
        child[concept color=poppurple!70!white] { node {Firma} }
    }
    child[concept color=poppink] { node[align=center]{Gramática:\\ Tú / Usted} 
        child[concept color=poppink!70!white] { node[align=center]{Imperativo\\ cortesía} }
        child[concept color=poppink!70!white] { node {Tratamiento} }
    }
    child[concept color=popgold] { node[align=center]{Estructura\\ general} 
        child[concept color=popgold!70!white] { node {Encabezado} }
        child[concept color=popgold!70!white] { node {Cuerpo} }
        child[concept color=popgold!70!white] { node {Despedida} }
    }
    child[concept color=popdark] { node[align=center]{Redacción\\ (80 palabras)} 
        child[concept color=popdark!70!white] { node {Registro} }
        child[concept color=popdark!70!white] { node {Coherencia} }
        child[concept color=popdark!70!white] { node {Ortografía} }
    }
    child[concept color=popblue!70!black] { node[align=center]{Contexto\\ camerounais} 
        child[concept color=popblue!50!white] { node {Corresponsal} }
        child[concept color=popblue!50!white] { node {Intercambios} }
    };
\end{tikzpicture}
\legende{Carte mentale : La correspondencia y los intercambios escritos (Leçon EA19)}
\end{popfigure}

\begin{bilan}

\textbf{Lexique clé : el mundo de la correspondencia}
\begin{center}
\begin{tabularx}{\linewidth}{|X|X|X|}
\hline
\rowcolor{popblueL}
\textbf{Español} & \textbf{Français} & \textbf{Contexto / Registro} \\
\hline
\esp{Querido / Querida} & Cher / Chère & Amical (tú) \\
\esp{Estimado / Estimada} & Estimé / Estimée & Formel (usted) \\
\esp{Distinguido / Distinguida} & Distingué / Distinguée & Très formel (administration) \\
\esp{Un abrazo / Un fuerte abrazo} & Une accolade / Une forte accolade & Clôture amicale \\
\esp{Saludos cordiales / Atentamente} & Salutations cordiales / Cordialement & Clôture formelle \\
\esp{Quedo a la espera de su respuesta} & J'attends votre réponse & Formel (fin de corps) \\
\esp{El asunto / El tema} & L'objet / Le sujet & Email / Lettre formelle \\
\esp{El remitente / El destinatario} & L'expéditeur / Le destinataire & En-tête \\
\esp{La firma} & La signature & Bas de page \\
\esp{Adjunto / Adjunto encontrará...} & Ci-joint / Vous trouverez ci-joint & Pièce jointe \\
\hline
\end{tabularx}
\end{center}

\textbf{Règles de grammaire à connaître par cœur}
\begin{center}
\begin{tabularx}{\linewidth}{|X|X|X|}
\hline
\rowcolor{popgreenL}
\textbf{Notion} & \textbf{Règle / Forme} & \textbf{Exemple traduit} \\
\hline
\textbf{Tú vs Usted} & \esp{Tú} = familiar (amis, famille, pairs). \esp{Usted} (Vd.) = respect, distance, hiérarchie, inconnu. Au Cameroun : \esp{usted} avec professeurs, autorités. & \esp{¿Cómo estás?} (ami) \\ \esp{¿Cómo está usted?} (professeur) \\
\hline
\textbf{Imperatif de politesse (Usted / Ustedes)} & Affirmatif = Subjonctif présent 3e pers. sing./plur. Négatif = \esp{No} + Subjonctif. Verbes irréguliers : \esp{sea, esté, vaya, tenga, sepa, dé}. & \esp{Sea amable de firmar.} (Soyez aimable de signer.) \\ \esp{No se preocupe.} (Ne vous inquiétez pas.) \\
\hline
\textbf{Imperatif familiar (Tú)} & Affirmatif = 3e pers. sing. indicatif présent (sauf 8 irréguliers). Négatif = \esp{No} + Subjonctif (2e pers. sing.). & \esp{Escribe pronto.} (Écris vite.) \\ \esp{No escribas mal.} (N'écris pas mal.) \\
\hline
\textbf{8 irréguliers (Tú affirmatif)} & \esp{Di, Haz, Ve, Pon, Sal, Sé, Ten, Ven}. Mnémotechnique : \textbf{« Di Haz Ve Pon Sal Sé Ten Ven »}. & \esp{Ven mañana.} (Viens demain.) \\
\hline
\textbf{Placement pronoms (Impératif)} & Affirmatif : pronoms \textbf{après} le verbe (liés, accent écrit si 3+ syllabes). Négatif : pronoms \textbf{avant} le verbe. & \esp{Escríbeme.} / \esp{No me escribas.} \\
\hline
\end{tabularx}
\end{center}

\textbf{Méthodes express}
\begin{enumerate}
\item \textbf{Rédiger une lettre/email amical (80 mots) :} 1. Lieu + date (Yaoundé, 12 de octubre de 2025). 2. Salut : \esp{Qerido/a [Prénom],}. 3. Corps : nouvelles perso, questions, lien avec le destinataire (vocabulaire courant, \esp{tú}). 4. Pré-clos : \esp{Espero tu carta / Escríbeme pronto.}. 5. Formule de fin : \esp{Un abrazo / Besos, [Ton prénom]}.
\item \textbf{Rédiger une lettre/email formel :} 1. En-tête expéditeur (gauche) + destinataire (droite) + lieu/date. 2. Objet : \esp{Asunto: Solicitud de información / Reclamación...}. 3. Salut : \esp{Estimado Sr. / Sra. [Apellido]:} ou \esp{Distinguido Sr. Director:}. 4. Corps : formules figées (\esp{Le escribo para...}, \esp{Quisiera solicitar...}, \esp{Adjunto encontrará...}), \esp{usted}. 5. Pré-clos : \esp{Quedo a la espera de su respuesta.}. 6. Formule de fin : \esp{Atentamente / Saludos cordiales, [Prénom Nom]}.
\item \textbf{Adapter le registre :} Identifier le destinataire $\rightarrow$ Choisir \esp{tú} ou \esp{usted} $\rightarrow$ Sélectionner formules d'ouverture/clôture correspondantes $\rightarrow$ Vérifier cohérence (pas de \esp{tú} avec \esp{Atentamente}).
\end{enumerate}

\end{bilan}

\begin{aretenir}[Checklist avant le Bac]
\begin{itemize}
\item $\square$ Je sais dater et localiser une lettre : \esp{Yaoundé, 15 de mayo de 2025}.
\item $\square$ Je distingue les formules d'ouverture : \esp{Querido/a} (ami) vs \esp{Estimado/a} / \esp{Distinguido/a} (formel).
\item $\square$ Je maîtrise les formules de clôture : \esp{Un abrazo} / \esp{Besos} (ami) vs \esp{Atentamente} / \esp{Saludos cordiales} (formel).
\item $\square$ J'utilise correctement \esp{tú} (verbes 2e pers. sg) et \esp{usted} (verbes 3e pers. sg).
\item $\square$ Je forme l'impératif de politesse (\esp{Usted}) : Subjonctif présent (\esp{Sea, Esté, Vaya, Tenga, Sepa, Dé}).
\item $\square$ Je place les pronoms après l'impératif affirmatif (\esp{Dígamelo}) et avant le négatif (\esp{No me lo diga}).
\item $\square$ Je connais les 8 irréguliers de l'impératif \esp{tú} affirmatif : \esp{Di, Haz, Ve, Pon, Sal, Sé, Ten, Ven}.
\item $\square$ Je structure
\end{itemize}
\end{aretenir}


\findecours

\end{document}
